% Les droits humains — ECM 1re Tech — Cours signé OnBuch+
% Programme officiel MINESEC. Compiler avec : tectonic N01-droits-humains.tex
\newcommand{\DOCMATIERE}{ECM}
\newcommand{\DOCNIVEAU}{1re Tech}
\newcommand{\DOCMODULE}{ECM – classe de Première technique}
\newcommand{\DOCLECON}{N01}
\newcommand{\DOCTITRE}{Les droits humains}
% OnBuch+ — préambule « cours signés » ENRICHI (compilé avec tectonic / XeLaTeX)
% Système visuel « Pop » d'OnBuch+ : fond crème (jamais blanc), bordures encre
% épaisses, ombres « dures » décalées, titres Archivo Black en capitales.
% Version MATHÉMATIQUES : graphes (pgfplots/tikz), boîtes pédagogiques
% (définition, propriété, méthode, attention, le savais-tu, exercice résolu,
% exercice, TP, bilan), page de garde et signature OnBuch+ ancrée en bas.
%
% Macros attendues dans main.tex AVANT \input{preamble} :
%   \DOCMATIERE  \DOCNIVEAU  \DOCMODULE  \DOCLECON (numéro)  \DOCTITRE
\documentclass[11pt]{article}

\usepackage{fontspec}
\usepackage[french]{babel}
\usepackage[a4paper,margin=2cm,top=3.4cm,bottom=2.8cm,headheight=1.4cm,footskip=1.3cm]{geometry}
\usepackage{amsmath,amssymb,amsfonts}
\usepackage{mathtools} % \xmapsto, \coloneqq... souvent utilisés par les modèles
\usepackage{cancel} % simplifications barrées (bilans de forces)
% \abs{z} (module, valeur absolue) et \norm{u} (norme), utilisés par les modèles
\providecommand{\abs}{}\let\abs\relax\DeclarePairedDelimiter{\abs}{\lvert}{\rvert}
\providecommand{\norm}{}\let\norm\relax\DeclarePairedDelimiter{\norm}{\lVert}{\rVert}
\usepackage[table]{xcolor} % [table] : fournit aussi \rowcolors (alternance de lignes)
\usepackage{pagecolor}
\usepackage{tikz}
\usetikzlibrary{arrows.meta,positioning,calc,shapes.geometric,shapes.misc,shapes.arrows,decorations.pathmorphing,decorations.pathreplacing,decorations.markings,patterns,shadows,mindmap,trees,fit,backgrounds,matrix,intersections,angles,quotes}
\usepackage{pgfplots}
\usepackage[version=4]{mhchem} % équations nucléaires \ce{^{235}_{92}U}
\usepackage[european,siunitx]{circuitikz} % schémas électriques
\pgfplotsset{compat=1.18}
\usepgfplotslibrary{fillbetween,groupplots} % aires entre courbes (calcul intégral)
\usepackage{enumitem}
\usepackage{fancyhdr}
% [most] charge toutes les bibliothèques tcolorbox (listings, minted, raster,
% documentation...) et gonfle inutilement la mémoire TeX sur un document long
% (~300 boîtes) : on ne charge que ce qui est réellement utilisé.
\usepackage[skins,breakable]{tcolorbox}
\usepackage{graphicx}
\usepackage{colortbl}
\usepackage{booktabs}
\usepackage{tabularx}
\usepackage{multirow}
\usepackage{diagbox} % case d'en-tête diagonale des tableaux à double entrée
% Colonnes extensibles centrée / gauche / droite (C, L, R), souvent utilisées par les modèles
\newcolumntype{C}{>{\centering\arraybackslash}X}
\newcolumntype{L}{>{\raggedright\arraybackslash}X}
\newcolumntype{R}{>{\raggedleft\arraybackslash}X}
\usepackage{array}
\usepackage{multicol}
\usepackage{siunitx}
% parse-numbers=false : accepte aussi \SI{2,5 \times 10^{-3}}{...}
\sisetup{locale=FR,output-decimal-marker={,},group-minimum-digits=4,parse-numbers=false}
\DeclareSIUnit\ohm{\ensuremath{\symup{\Omega}}} % U+2126 absent de la police
\DeclareSIUnit\division{div} % réglages d'oscilloscope (ms/div, V/div)
\DeclareSIUnit\tr{tr} % tours (vitesse de rotation, ex. \SI{1500}{\tr\per\minute})
\DeclareSIUnit\molar{M} % concentration molaire (ex. \SI{3}{\molar}, \SI{1}{\micro\molar})
\DeclareSIUnit\an{an} % année (le modèle confond parfois avec \year, un compteur TeX) — ex. \SI{5}{\kilo\meter\per\an}
\DeclareSIUnit\FCFA{FCFA} % franc CFA (ex. \SI{500}{\FCFA\per\kilo\gram})
\DeclareSIUnit\degreeBrix{\textdegree Brix} % teneur en sucre d'un sirop/jus (réfractométrie)
\DeclareSIUnit\month{mois} % le modèle utilise parfois \month, un compteur TeX, comme unité de durée
\DeclareSIUnit\microliter{\micro\liter} % le modèle écrit parfois \microliter en un seul token
\DeclareSIUnit\million{millions} % dénombrement (ex. \SI{15}{\million\per\mL} spermatozoïdes)
\usepackage{qrcode}
% \degree : commande fréquemment utilisée par les modèles pour un angle ou une
% température en dehors de \SI{}{} (non fournie par les paquets chargés).
\providecommand{\degree}{^{\circ}}
% \qed (fin de démonstration), utilisé par les modèles sans amsthm
\providecommand{\qed}{\hfill$\square$}
% \barre : double barre des valeurs interdites dans un tableau de variations (tkz-tab)
\providecommand{\barre}{\ensuremath{\|}}
% environnement proof (amsthm non chargé) utilisé par les modèles
\ifdefined\proof\else\newenvironment{proof}[1][Démonstration]{\par\noindent\textit{#1.}\ }{\qed\par}\fi
\usepackage{lastpage}
\usepackage{hyperref}
\hypersetup{hidelinks}

% ============================================================
% Palette « Pop » OnBuch+ — mêmes hex que la classe P (plus_ui.dart)
% ============================================================
\definecolor{poporange}{HTML}{F59321}
\definecolor{popdark}{HTML}{E07A0C}
\definecolor{popcream}{HTML}{FFF6E8}
\definecolor{popink}{HTML}{1C1714}
\definecolor{poppurple}{HTML}{7A5AE0}
\definecolor{popgreen}{HTML}{2E9E6A}
\definecolor{poppink}{HTML}{E0457B}
\definecolor{popblue}{HTML}{2F7DE1}
\definecolor{popgold}{HTML}{C98A12}
\definecolor{popmuted}{HTML}{8A7F76}
\definecolor{popline}{HTML}{EADFD2}
\definecolor{popgray}{HTML}{8A7F76}
\colorlet{popgrey}{popgray}
% Alias tolérants (le modèle nomme parfois les couleurs différemment)
\colorlet{popwhite}{white}
\colorlet{popblack}{popink}
\colorlet{popred}{poppink}
\colorlet{popyellow}{popgold}
% Teintes claires (fonds de tableaux/encadrés) — le modèle nomme parfois
% les couleurs avec un suffixe L (ex. popblueL) sans les avoir définies.
\colorlet{poporangeL}{poporange!20}
\colorlet{popdarkL}{popdark!20}
\colorlet{poppurpleL}{poppurple!20}
\colorlet{popgreenL}{popgreen!20}
\colorlet{poppinkL}{poppink!20}
\colorlet{popblueL}{popblue!20}
\colorlet{popgoldL}{popgold!20}
\colorlet{popredL}{popred!20}
\colorlet{popgrayL}{popgray!20}
% Teintes claires pour fonds de tableaux / zones
\colorlet{popblueL}{popblue!10!white}
\colorlet{poporangeL}{poporange!14!white}
\colorlet{popgreenL}{popgreen!12!white}
\colorlet{poppurpleL}{poppurple!12!white}
\colorlet{poppinkL}{poppink!12!white}
\colorlet{popgoldL}{popgold!14!white}

\pagecolor{popcream}

% ============================================================
% Typographie
% ============================================================
\defaultfontfeatures{Ligatures=TeX}
\setmainfont{PlusJakartaSans-Medium.ttf}[Path=../../../fonts/,BoldFont=PlusJakartaSans-Bold.ttf]
% Maths en sans-serif (Fira Math) avec chiffres et lettres droites de Plus
% Jakarta, pour que les formules se fondent dans le texte.
\usepackage[math-style=ISO,bold-style=ISO]{unicode-math}
\setmathfont{FiraMath-Regular.otf}[Path=../../../fonts/]
\setmathfont{PlusJakartaSans-Medium.ttf}[Path=../../../fonts/,range={up/num,up/{Latin,latin},\mathup}]
\usepackage{icomma} % virgule décimale sans espace en mode math
% Caractères mathématiques Unicode absents des polices de texte (Plus Jakarta,
% Archivo Black) : rendus par leur équivalent mathématique, en texte comme en
% formule — sinon ils s'affichent en carré vide (ex. « Arithmétique dans ℤ »).
\usepackage{newunicodechar}
\newunicodechar{ℝ}{\ensuremath{\mathbb{R}}}
\newunicodechar{ℤ}{\ensuremath{\mathbb{Z}}}
\newunicodechar{ℂ}{\ensuremath{\mathbb{C}}}
\newunicodechar{ℕ}{\ensuremath{\mathbb{N}}}
\newunicodechar{ℚ}{\ensuremath{\mathbb{Q}}}
\newunicodechar{𝔻}{\ensuremath{\mathbb{D}}}
\newunicodechar{α}{\ensuremath{\alpha}}
\newunicodechar{β}{\ensuremath{\beta}}
\newunicodechar{γ}{\ensuremath{\gamma}}
\newunicodechar{δ}{\ensuremath{\delta}}
\newunicodechar{ε}{\ensuremath{\varepsilon}}
\newunicodechar{θ}{\ensuremath{\theta}}
\newunicodechar{λ}{\ensuremath{\lambda}}
\newunicodechar{μ}{\ensuremath{\mu}}
\newunicodechar{σ}{\ensuremath{\sigma}}
\newunicodechar{τ}{\ensuremath{\tau}}
\newunicodechar{φ}{\ensuremath{\varphi}}
\newunicodechar{ψ}{\ensuremath{\psi}}
\newunicodechar{ω}{\ensuremath{\omega}}
\newunicodechar{Σ}{\ensuremath{\Sigma}}
% majuscules produites par \MakeUppercase dans les titres (α-aminés -> Α)
\newunicodechar{Α}{\ensuremath{\alpha}}
\newunicodechar{Β}{\ensuremath{\beta}}
\newunicodechar{Γ}{\ensuremath{\Gamma}}
\newunicodechar{Δ}{\ensuremath{\Delta}}
\newunicodechar{Ε}{\ensuremath{\varepsilon}}
\newunicodechar{Θ}{\ensuremath{\Theta}}
\newunicodechar{Λ}{\ensuremath{\Lambda}}
\newunicodechar{Μ}{\ensuremath{\mu}}
\newunicodechar{Τ}{\ensuremath{\tau}}
\newunicodechar{Φ}{\ensuremath{\Phi}}
\newunicodechar{Ψ}{\ensuremath{\Psi}}
\newunicodechar{Ω}{\ensuremath{\Omega}}
\newunicodechar{↦}{\ensuremath{\mapsto}}
\newunicodechar{∈}{\ensuremath{\in}}
\newunicodechar{∉}{\ensuremath{\notin}}
\newunicodechar{∧}{\ensuremath{\wedge}}
\newunicodechar{⇔}{\ensuremath{\Leftrightarrow}}
\newunicodechar{⟺}{\ensuremath{\Longleftrightarrow}}
\newunicodechar{⇒}{\ensuremath{\Rightarrow}}
\newunicodechar{⟹}{\ensuremath{\Longrightarrow}}
\newunicodechar{∘}{\ensuremath{\circ}}
\newunicodechar{∪}{\ensuremath{\cup}}
\newunicodechar{∩}{\ensuremath{\cap}}
\newunicodechar{≡}{\ensuremath{\equiv}}
\newunicodechar{⊂}{\ensuremath{\subset}}
\newunicodechar{⊥}{\ensuremath{\perp}}
\newunicodechar{∃}{\ensuremath{\exists}}
\newunicodechar{∀}{\ensuremath{\forall}}
\newunicodechar{∛}{\ensuremath{\sqrt[3]{\,}}}
\newunicodechar{∜}{\ensuremath{\sqrt[4]{\,}}}
\newunicodechar{⃗}{\ensuremath{{}^{\to}}}
\newunicodechar{ⁿ}{\ifmmode^{n}\else\textsuperscript{\ensuremath{n}}\fi}
\newunicodechar{ˣ}{\ifmmode^{x}\else\textsuperscript{\ensuremath{x}}\fi}
\newunicodechar{ᵇ}{\ifmmode^{b}\else\textsuperscript{\ensuremath{b}}\fi}
\newunicodechar{ʳ}{\ifmmode^{r}\else\textsuperscript{\ensuremath{r}}\fi}
\newunicodechar{ᵘ}{\ifmmode^{u}\else\textsuperscript{\ensuremath{u}}\fi}
\newunicodechar{ʸ}{\ifmmode^{y}\else\textsuperscript{\ensuremath{y}}\fi}
\newunicodechar{ᵃ}{\ifmmode^{a}\else\textsuperscript{\ensuremath{a}}\fi}
\newunicodechar{ᵉ}{\ifmmode^{e}\else\textsuperscript{\ensuremath{e}}\fi}
\newunicodechar{ᵏ}{\ifmmode^{k}\else\textsuperscript{\ensuremath{k}}\fi}
\newunicodechar{ᵖ}{\ifmmode^{p}\else\textsuperscript{\ensuremath{p}}\fi}
\newunicodechar{ᴺ}{\ifmmode^{N}\else\textsuperscript{\ensuremath{N}}\fi}
\newunicodechar{ᶜ}{\ifmmode^{c}\else\textsuperscript{\ensuremath{c}}\fi}
\newunicodechar{ʰ}{\ifmmode^{h}\else\textsuperscript{\ensuremath{h}}\fi}
\newunicodechar{ᵠ}{\ifmmode^{q}\else\textsuperscript{\ensuremath{q}}\fi}
\newunicodechar{ₙ}{\ifmmode_{n}\else\textsubscript{\ensuremath{n}}\fi}
\newunicodechar{ₐ}{\ifmmode_{a}\else\textsubscript{\ensuremath{a}}\fi}
\newunicodechar{ᵢ}{\ifmmode_{i}\else\textsubscript{\ensuremath{i}}\fi}
\newunicodechar{ᵣ}{\ifmmode_{r}\else\textsubscript{\ensuremath{r}}\fi}
\newunicodechar{ₖ}{\ifmmode_{k}\else\textsubscript{\ensuremath{k}}\fi}
\newunicodechar{ᵦ}{\ifmmode_{\beta}\else\textsubscript{\ensuremath{\beta}}\fi}
\newunicodechar{ₓ}{\ifmmode_{x}\else\textsubscript{\ensuremath{x}}\fi}
\newfontfamily\popdisplay{ArchivoBlack-Regular.ttf}[Path=../../../fonts/]
\newfontfamily\poplogo{SpaceGrotesk-Bold.ttf}[Path=../../../fonts/]
\newfontfamily\popsemi{PlusJakartaSans-SemiBold.ttf}[Path=../../../fonts/]
\newfontfamily\popxbold{PlusJakartaSans-ExtraBold.ttf}[Path=../../../fonts/]
\newfontfamily\popxboldls{PlusJakartaSans-ExtraBold.ttf}[Path=../../../fonts/,LetterSpace=3.0]

\setlist[enumerate]{leftmargin=1.7em,itemsep=5pt,topsep=4pt}
\setlist[itemize]{leftmargin=1.7em,itemsep=4pt,topsep=4pt,label={\color{poporange}\textbullet}}
\setlength{\parindent}{0pt}
\setlength{\parskip}{5pt}
\renewcommand{\arraystretch}{1.25}

% pgfplots : style Pop par défaut
\pgfplotsset{
  every axis/.append style={axis line style={popink,line width=1pt},tick style={popink},
    grid style={popline,line width=0.5pt},label style={font=\small},tick label style={font=\footnotesize}},
  popcurve/.style={poporange,line width=1.8pt,no markers,smooth},
}
% Même style disponible dans un \draw TikZ simple (hors pgfplots)
\tikzset{popcurve/.style={poporange,line width=1.8pt}}
\tikzset{popgrid/.style={popline,very thin}}
\colorlet{popcurve}{poporange} % le nom du style est parfois utilisé comme couleur

% ============================================================
% Pill / squiggle
% ============================================================
\newcommand{\poppill}[3]{%
  \tikz[baseline=-0.55ex]\node[draw=popink,line width=1.2pt,rounded corners=2.6pt,%
    fill=#2,inner xsep=6.5pt,inner ysep=3pt]{%
    \popxboldls\fontsize{7}{7}\selectfont\color{#3}\MakeUppercase{#1}};%
}
\newcommand{\popsquiggle}[1]{%
  \tikz[baseline=-0.4ex]{
    \draw[#1,line width=2.6pt,line cap=round]
      plot [smooth] coordinates {(0,0.09) (0.34,0.01) (0.68,0.21) (1.02,0.01) (1.36,0.10)};
  }%
}
% Mot-clé surligné (marqueur orange sous le texte)
\newcommand{\cle}[1]{{\popxbold\color{popdark}#1}}

% ============================================================
% En-tête / pied
% ============================================================
\pagestyle{fancy}
\fancyhf{}
\renewcommand{\headrulewidth}{0pt}
\renewcommand{\footrulewidth}{0pt}
\lhead{\raisebox{-0.15ex}{\poplogo\fontsize{13}{13}\selectfont\color{popink}OnBuch\color{poporange}+}}
\rhead{\poppill{\DOCMATIERE\ \textbullet\ \DOCNIVEAU\ \textbullet\ Leçon \DOCLECON}{white}{popink}}
\lfoot{\popxbold\fontsize{7.6}{7.6}\selectfont\color{popmuted}\MakeUppercase{Cours signé OnBuch+}\ \textbullet\ {\popsemi\DOCTITRE}}
\rfoot{\tikz[baseline=-0.6ex]\node[rounded corners=8.5pt,fill=popink,minimum height=17pt,minimum width=17pt,inner xsep=5pt,inner sep=0]{\popxbold\fontsize{8}{8}\selectfont\color{popcream}\thepage\,/\,\pageref*{LastPage}};}

% ============================================================
% Titres
% ============================================================
\newcommand{\chaptitre}[2]{%
  \begin{tcolorbox}[enhanced,colback=poporange,colframe=popink,boxrule=2.6pt,arc=11pt,
      drop shadow=popink,left=15pt,right=15pt,top=12pt,bottom=11pt,before skip=3pt,after skip=18pt]
    \poppill{#2}{popink}{popcream}\par\vspace{9pt}
    {\raggedright\hyphenpenalty=10000\exhyphenpenalty=10000\popdisplay\fontsize{22}{24}\selectfont\color{popcream}\MakeUppercase{#1}\par}
  \end{tcolorbox}
}
% Section numérotée : pastille orange avec le numéro + titre Archivo
\newcounter{popsec}
\newcommand{\coursec}[1]{%
  \stepcounter{popsec}\setcounter{popsub}{0}%
  \par\vspace{16pt}\noindent
  \begin{minipage}{\linewidth}
  \tikz[baseline=-0.7ex]\node[draw=popink,line width=1.6pt,fill=poporange,rounded corners=4pt,minimum size=22pt,inner sep=0]{\popdisplay\fontsize{12}{12}\selectfont\color{popcream}\thepopsec};%
  \hspace{8pt}{\popdisplay\fontsize{15}{17}\selectfont\color{popink}\MakeUppercase{#1}}\par
  \vspace{4pt}\noindent{\color{poporange}\rule{36pt}{3.6pt}}
  \end{minipage}\par\nopagebreak\vspace{8pt}
}
\newcounter{popsub}
\newcommand{\courssub}[1]{%
  \stepcounter{popsub}%
  \par\vspace{9pt}\noindent{\popxbold\fontsize{11.5}{13}\selectfont\color{popink}{\color{poporange}\thepopsec.\thepopsub}\hspace{6pt}#1}\par\nopagebreak\vspace{4pt}
}

% ============================================================
% Boîtes pédagogiques — même recette : fond blanc, bordure épaisse colorée,
% ombre dure encre, badge plein. Titre optionnel [#1].
% ============================================================
\tcbset{popbase/.style={breakable,enhanced,colback=white,boxrule=2.3pt,arc=9pt,
  shadow={3pt}{-3pt}{0pt}{popink},left=14pt,right=14pt,top=11pt,bottom=11pt,before skip=11pt,after skip=15pt,
  fontupper=\popsemi\fontsize{10.6}{14.5}\selectfont\color{popink},
  fontlower=\popsemi\fontsize{10.6}{14.5}\selectfont\color{popink}}}
% Coche dessinée (le glyphe ✓ n'existe pas dans les polices du document)
\newcommand{\popcheck}[1]{\tikz[baseline=-0.5ex]\draw[#1,line width=1.6pt,line cap=round,line join=round](-0.13,0.0)--(-0.04,-0.09)--(0.13,0.1);}
\providecommand{\coche}{\popcheck{popgreen}}
\providecommand{\resultat}[1]{\par\smallskip\noindent\fcolorbox{popgreen}{popgreenL}{\parbox{\dimexpr\linewidth-2\fboxsep-2\fboxrule\relax}{\textbf{Résultat.} #1}}\par\smallskip}
\newcommand{\popboxhead}[3]{\poppill{#1}{#2}{white}\ifx\relax#3\relax\else\hspace{6pt}{\popxbold\fontsize{10.6}{12}\selectfont\color{popink}#3}\fi\par\vspace{7pt}}

\newtcolorbox{aretenir}[1][]{popbase,colframe=popgreen,before upper={\popboxhead{À retenir}{popgreen}{#1}}}
\newtcolorbox{exemplebox}[1][]{popbase,colframe=poporange,before upper={\popboxhead{Exemple}{poporange}{#1}}}
\newtcolorbox{definition}[1][]{popbase,colframe=popblue,before upper={\popboxhead{Définition}{popblue}{#1}}}
\newtcolorbox{propriete}[1][]{popbase,colframe=poppurple,before upper={\popboxhead{Propriété}{poppurple}{#1}}}
\newtcolorbox{methode}[1][]{popbase,colframe=popgold,before upper={\popboxhead{Méthode}{popgold}{#1}}}
\newtcolorbox{attention}[1][]{popbase,colframe=poppink,before upper={\popboxhead{Attention}{poppink}{#1}}}
\newtcolorbox{experience}[1][]{popbase,colframe=popblue,colback=popblueL,before upper={\popboxhead{Expérience}{popblue}{#1}}}
% « Le savais-tu ? » : carte violette pleine (contexte, culture, Cameroun)
\newtcolorbox{savaistu}[1][]{popbase,colback=poppurple,colframe=popink,
  fontupper=\popsemi\fontsize{10.4}{14.2}\selectfont\color{white},
  before upper={\poppill{Le savais-tu\,?}{white}{poppurple}\ifx\relax#1\relax\else\hspace{6pt}{\popxbold\color{white}#1}\fi\par\vspace{7pt}}}
% Exercice résolu : énoncé en haut, \tcblower puis la solution sur fond crème
\newcounter{popexo}\newcounter{popexr}
\newtcolorbox{exoresolu}[1][]{popbase,colframe=poporange,colbacklower=popcream,
  segmentation style={popink,line width=1.2pt,dashed},
  before upper={\stepcounter{popexr}\popboxhead{Exercice résolu \thepopexr}{poporange}{#1}},
  before lower={\poppill{Solution}{popink}{popcream}\par\vspace{6pt}}}
% Exercice (non résolu en place) — la correction est dans la partie Corrigés
\newtcolorbox{exercice}[1][]{popbase,colframe=popink,
  before upper={\stepcounter{popexo}\popboxhead{Exercice \thepopexo}{popink}{#1}}}
% Corrigé
\newtcolorbox{corrige}[1][]{popbase,colframe=popgreen,colback=popgreenL,
  before upper={\popboxhead{Corrigé}{popgreen}{#1}}}
% Objectifs / prérequis / bilan
\newtcolorbox{objectifs}{popbase,colframe=popink,colback=poporangeL,
  before upper={\popboxhead{Objectifs de la leçon}{popink}{}}}
\newtcolorbox{prerequis}{popbase,colframe=popink,colback=popblueL,
  before upper={\popboxhead{Prérequis}{popblue}{}}}
\newtcolorbox{bilan}{popbase,colframe=popink,colback=white,boxrule=2.8pt,
  before upper={\popboxhead{Fiche bilan}{poporange}{L'essentiel à connaître pour l'examen}}}
% Figure encadrée avec légende
\newtcolorbox{popfigure}[1][]{enhanced,breakable=false,colback=white,colframe=popline,boxrule=1.4pt,arc=8pt,
  left=10pt,right=10pt,top=10pt,bottom=8pt,before skip=10pt,after skip=12pt,halign=center,
  lower separated=false,halign lower=center,
  fontlower=\popsemi\fontsize{9}{11.5}\selectfont\color{popmuted},
  #1}
\newcommand{\legende}[1]{\par\vspace{6pt}{\popsemi\fontsize{9}{11.5}\selectfont\color{popmuted}#1}}

% ============================================================
% Signature OnBuch+
% ============================================================
\newcommand{\poplogoinline}[1]{{\poplogo\fontsize{#1}{#1}\selectfont\color{popink}OnBuch\color{poporange}+}}
\newcommand{\signatureonbuch}{%
  \begin{tcolorbox}[enhanced,colback=popink,colframe=popink,boxrule=2.2pt,arc=10pt,
      drop shadow=poporange,left=14pt,right=14pt,top=10pt,bottom=10pt,before skip=0pt,after skip=0pt]
    \tikz[baseline=-0.7ex]\node[circle,fill=popgreen,draw=popcream,line width=1.3pt,minimum size=22pt,inner sep=0]{\popcheck{white}};%
    \hspace{9pt}\begin{minipage}[c]{0.62\linewidth}
      {\popxbold\fontsize{12}{13}\selectfont\color{popcream}Signé\ \textbullet\ {\poplogo OnBuch\color{poporange}+}}\par\vspace{2pt}
      {\raggedright\popsemi\fontsize{8}{10.5}\selectfont\color{popline}\MakeUppercase{Équipe pédagogique OnBuch+}\\\MakeUppercase{Conforme au programme officiel MINESEC}\par}
    \end{minipage}\hfill
    {\poplogo\fontsize{22}{22}\selectfont\color{popcream}OnBuch\color{poporange}+}
  \end{tcolorbox}%
}

% ============================================================
% Page de garde
% #1 titre · #2 matière · #3 niveau · #4 module · #5 n° leçon · #6 durée ·
% #7 description / objectifs (texte ou itemize)
% ============================================================
\newcommand{\pagegarde}[7]{%
  \thispagestyle{empty}
  \newgeometry{a4paper,margin=1.7cm,top=1.6cm,bottom=1.4cm}%
  \begin{tikzpicture}[remember picture,overlay]
    \fill[poporange!18] ([xshift=-3.2cm,yshift=-3.6cm]current page.north east) circle (4.6cm);
    \fill[poppurple!14] ([xshift=2.4cm,yshift=7.6cm]current page.south west) circle (3.8cm);
  \end{tikzpicture}%
  {\poplogo\fontsize{30}{30}\selectfont\color{popink}OnBuch\color{poporange}+}\hfill
  \raisebox{0.6ex}{\poppill{Cours signé \textbullet\ Premium}{popink}{popcream}}\par
  \vspace{1pt}\popsquiggle{poppurple}\par
  \vspace{8pt}
  \begin{tcolorbox}[enhanced,colback=poporange,colframe=popink,boxrule=2.8pt,arc=13pt,
      drop shadow=popink,left=18pt,right=18pt,top=12pt,bottom=13pt,after skip=10pt]
    \poppill{#2 \textbullet\ #3}{popink}{popcream}\hspace{5pt}\poppill{#4}{white}{popink}\par\vspace{8pt}
    {\popdisplay\fontsize{14}{14}\selectfont\color{popink}LEÇON #5}\par\vspace{4pt}
    {\raggedright\hyphenpenalty=10000\exhyphenpenalty=10000\popdisplay\fontsize{25}{27}\selectfont\color{popcream}\MakeUppercase{#1}\par}
    \vspace{8pt}
    {\popxbold\fontsize{9.5}{9.5}\selectfont\color{popink}\MakeUppercase{Durée conseillée : #6}}
  \end{tcolorbox}
  \begin{tcolorbox}[enhanced,colback=white,colframe=popink,boxrule=2.2pt,arc=10pt,
      drop shadow=popink,left=16pt,right=16pt,top=10pt,bottom=10pt,after skip=8pt]
    \poppill{Dans cette leçon}{poppurple}{white}\par\vspace{5pt}
    \popsemi\fontsize{9.3}{12.6}\selectfont\color{popink}#7
  \end{tcolorbox}
  \noindent\begin{minipage}[t]{0.487\linewidth}
    \begin{tcolorbox}[enhanced,colback=popgreen,colframe=popink,boxrule=2.2pt,arc=10pt,
        drop shadow=popink,left=11pt,right=11pt,top=10pt,bottom=10pt,height=3.1cm,valign=center]
      \tcbox[colback=white,colframe=popink,boxrule=1.3pt,arc=5pt,boxsep=0pt,nobeforeafter,
        left=3pt,right=3pt,top=3pt,bottom=3pt,valign=center]{\qrcode[height=1.9cm]{https://play.google.com/store/apps/details?id=com.luvvix.onbuch&hl=fr}}%
      \hspace{8pt}\begin{minipage}[c]{0.5\linewidth}
        {\popdisplay\fontsize{11}{12.5}\selectfont\color{white}\MakeUppercase{Télécharge OnBuch}}\par\vspace{4pt}
        {\raggedright\popsemi\fontsize{8.4}{11}\selectfont\color{white}Gratuit \textbullet\ Google Play\\Tuteur IA, annales, concours, résultats\par}
      \end{minipage}
    \end{tcolorbox}
  \end{minipage}\hfill
  \begin{minipage}[t]{0.487\linewidth}
    \begin{tcolorbox}[enhanced,colback=poppurple,colframe=popink,boxrule=2.2pt,arc=10pt,
        drop shadow=popink,left=11pt,right=11pt,top=10pt,bottom=10pt,height=3.1cm,valign=center]
      \tikz[baseline=-0.7ex]\node[circle,fill=white,draw=popink,line width=1.3pt,minimum size=1.6cm,inner sep=0]{\popdisplay\fontsize{24}{24}\selectfont\color{poppurple}+};%
      \hspace{8pt}\begin{minipage}[c]{0.6\linewidth}
        {\popdisplay\fontsize{11}{12.5}\selectfont\color{white}\MakeUppercase{Passe à OnBuch+}}\par\vspace{4pt}
        {\raggedright\popsemi\fontsize{8.4}{11}\selectfont\color{white}Tous les cours signés, Tuteur IA illimité, fiches PDF — onglet OnBuch+ de l'app\par}
      \end{minipage}
    \end{tcolorbox}
  \end{minipage}
  \vspace{8pt}
  \popcontenu
  \vfill
  \signatureonbuch
  \restoregeometry
  \newpage
}

% Rangée « Ce cours contient » (page de garde)
\newcommand{\poptile}[3]{% #1 couleur #2 gros texte #3 légende
  \begin{tcolorbox}[enhanced,colback=#1,colframe=popink,boxrule=2pt,arc=9pt,shadow={2.5pt}{-2.5pt}{0pt}{popink},
      left=6pt,right=6pt,top=9pt,bottom=9pt,halign=center,valign=center,height=1.75cm,width=\linewidth,nobeforeafter]
    {\popdisplay\fontsize{12.5}{13}\selectfont\color{white}\MakeUppercase{#2}}\par\vspace{3pt}
    {\centering\popsemi\fontsize{7.8}{9.5}\selectfont\color{white}#3\par}
  \end{tcolorbox}}
\newcommand{\popcontenu}{%
  \par\noindent{\popxboldls\fontsize{7.5}{7.5}\selectfont\color{popmuted}CE COURS SIGNÉ CONTIENT}\par\vspace{7pt}
  \noindent\begin{minipage}[t]{0.235\linewidth}\poptile{popblue}{Cours}{complet, illustré, pas à pas}\end{minipage}\hfill
  \begin{minipage}[t]{0.235\linewidth}\poptile{poporange}{Méthodes}{et exercices résolus}\end{minipage}\hfill
  \begin{minipage}[t]{0.235\linewidth}\poptile{poppink}{Exercices}{type examen + corrigés}\end{minipage}\hfill
  \begin{minipage}[t]{0.235\linewidth}\poptile{popgreen}{Fiche bilan}{l'essentiel à retenir}\end{minipage}\par}

% Sommaire
\newtcolorbox{sommaire}{enhanced,colback=white,colframe=popink,boxrule=2.2pt,arc=10pt,
  drop shadow=popink,left=18pt,right=18pt,top=13pt,bottom=13pt,before skip=6pt,after skip=20pt,
  before upper={\poppill{Sommaire}{poppurple}{white}\par\vspace{9pt}\popsemi\fontsize{10.6}{14.5}\selectfont\color{popink}}}

% Signature de fin de cours — poussée tout en bas de la dernière page
\newcommand{\findecours}{%
  \par\vfill\nopagebreak
  \begin{center}\popsquiggle{poporange}\end{center}\vspace{4pt}
  \signatureonbuch
}
\AtBeginDocument{\def\llbracket{\mathopen{[\mkern-3.5mu[}}\def\rrbracket{\mathclose{]\mkern-3.5mu]}}} % intervalles d'entiers (glyphe ⟦ absent de la police)
% Glyphes absents de Fira Math : redéfinis à partir de glyphes disponibles
\AtBeginDocument{%
  \def\setminus{\mathbin{\backslash}}\let\smallsetminus\setminus
  \def\vdots{\mathord{\vbox{\baselineskip3.2pt\lineskiplimit0pt\kern4pt\hbox{.}\hbox{.}\hbox{.}}}}%
  \def\ddots{\mathinner{\mkern1mu\raise6.5pt\hbox{.}\mkern2mu\raise3.6pt\hbox{.}\mkern2mu\raise0.7pt\hbox{.}\mkern1mu}}%
  \def\triangle{\mathord{\scalebox{0.9}{$\symup{\Delta}$}}}%
  \def\nabla{\mathord{\raisebox{\depth}{\scalebox{1}[-1]{$\symup{\Delta}$}}}}%
  \def\checkmark{\popcheck{popdark}}%
  \def\star{\mathbin{\ast}}\def\bigstar{\mathord{\ast}}%
  \def\diamond{\mathbin{\scalebox{0.6}{$\Diamond$}}}%
  \def\bigcup{\mathop{\mathchoice{\raisebox{-0.3em}{\scalebox{1.7}{$\cup$}}}{\scalebox{1.25}{$\cup$}}{\cup}{\cup}}\displaylimits}%
  \def\bigcap{\mathop{\mathchoice{\raisebox{-0.3em}{\scalebox{1.7}{$\cap$}}}{\scalebox{1.25}{$\cap$}}{\cap}{\cap}}\displaylimits}%
  \def\lesseqqgtr{\mathrel{\lessgtr}}\def\bigoplus{\mathop{\oplus}\displaylimits}%
}
\providecommand{\ssi}{si et seulement si} % abréviation fréquente
\AtBeginDocument{\providecommand{\wideparen}{\overparen}} % arc de cercle

%% FIGFIX-BEGIN v4 — correctifs globaux des figures (docs/onbuch-generation/tools/figfix)
\usepackage{adjustbox}\usepackage{etoolbox}
% toute figure TikZ/pgfplots plus large que la ligne est réduite à la largeur disponible
\BeforeBeginEnvironment{tikzpicture}{\begin{adjustbox}{max width=\linewidth}}
\AfterEndEnvironment{tikzpicture}{\end{adjustbox}}
% légendes pgfplots lisibles par-dessus les courbes
\pgfplotsset{every axis/.append style={legend style={fill=white,fill opacity=0.92,text opacity=1,draw=black!15,inner sep=3pt}}}
% étiquettes d'arêtes (syntaxe edge["..."]) avec fond blanc
\tikzset{every edge quotes/.append style={fill=white,inner sep=1.2pt}}
%% FIGFIX-END
\begin{document}

\pagegarde{Les droits humains}{ECM}{1re Tech}{ECM – classe de Première technique}{N01}{3 séances (fiche de progression harmonisée)}{Cette leçon explore les fondements, la classification et les mécanismes de protection des droits de l'homme, avec un focus sur le cadre juridique camerounais et l'étude de terrain d'une structure locale de promotion (TD1). Elle développe la capacité à analyser une situation concrète de violation et à rédiger une démarche de saisine argumentée.
\vspace{4pt}
\begin{multicols}{2}\small
\begin{itemize}
\item À la fin de cette leçon, je sais définir les droits de l'homme, en citer les caractéristiques (universalité, indivisibilité, interdépendance) et les classer en générations.
\item À la fin de cette leçon, je sais identifier les principaux instruments juridiques de protection (DUDH, PIDCP, PIDESC, Charte africaine, Constitution du Cameroun) et situer leur hiérarchie.
\item À la fin de cette leçon, je sais décrire l'architecture institutionnelle de protection (ONU : CDH, Comités ; UA : CADHP ; Cameroun : CNDHL, Tribunal administratif, Cours d'appel).
\item À la fin de cette leçon, je sais analyser le fonctionnement d'une structure locale de promotion (CNDHL, ONG, Association) à partir d'une grille d'observation (mission, moyens, actions, limites).
\item À la fin de cette leçon, je sais qualifier juridiquement une situation concrète de violation (ex: travail forcé, discrimination, arrestation arbitraire) en identifiant le droit violé et l'instance compétente.
\item À la fin de cette leçon, je sais rédiger une plainte ou une lettre de saisine structurée (faits, fondement juridique, demandes, pièces jointes) adressée à l'institution appropriée.
\end{itemize}\end{multicols}}

\begin{sommaire}
\begin{multicols}{2}
\begin{enumerate}
\item 1. Fondements, caractéristiques et classification des droits de l'homme
\item 2. Le cadre juridique : Hiérarchie des normes et justiciabilité au Cameroun
\item 3. Les institutions de protection : Architecture multiniveau (ONU, UA, National)
\item 4. TD1 : Étude de terrain d'une structure locale de promotion des droits de l'homme (CNDHL ou ONG)
\item 5. Mise en situation : Qualifier, Saisiner, Argumenter (Intégration des savoir-faire)
\item Activité d'intégration
\item Méthodes et exercices résolus
\item Exercices
\item Corrigés des exercices
\item Fiche bilan
\end{enumerate}
\end{multicols}
\end{sommaire}

\begin{objectifs}
\begin{itemize}
\item À la fin de cette leçon, je sais définir les droits de l'homme, en citer les caractéristiques (universalité, indivisibilité, interdépendance) et les classer en générations.
\item À la fin de cette leçon, je sais identifier les principaux instruments juridiques de protection (DUDH, PIDCP, PIDESC, Charte africaine, Constitution du Cameroun) et situer leur hiérarchie.
\item À la fin de cette leçon, je sais décrire l'architecture institutionnelle de protection (ONU : CDH, Comités ; UA : CADHP ; Cameroun : CNDHL, Tribunal administratif, Cours d'appel).
\item À la fin de cette leçon, je sais analyser le fonctionnement d'une structure locale de promotion (CNDHL, ONG, Association) à partir d'une grille d'observation (mission, moyens, actions, limites).
\item À la fin de cette leçon, je sais qualifier juridiquement une situation concrète de violation (ex: travail forcé, discrimination, arrestation arbitraire) en identifiant le droit violé et l'instance compétente.
\item À la fin de cette leçon, je sais rédiger une plainte ou une lettre de saisine structurée (faits, fondement juridique, demandes, pièces jointes) adressée à l'institution appropriée.
\item À la fin de cette leçon, je sais construire un organigramme fonctionnel des institutions de protection et un diagramme en barres représentant les types de violations les plus fréquents au Cameroun d'après les rapports annuels de la CNDHL.
\item À la fin de cette leçon, je sais argumenter à l'oral comme à l'écrit sur l'effectivité des droits au Cameroun en confrontant les textes aux réalités de terrain observées lors du TD1.
\end{itemize}
\end{objectifs}

\begin{prerequis}
\begin{itemize}
\item Notion de citoyenneté et de civisme (programme de 4ème / 3ème) : droits et devoirs du citoyen.
\item Connaissance de l'organisation politique du Cameroun (institutions de la République, séparation des pouvoirs).
\item Notions de base en droit : distinction droit national / droit international, notion de traité.
\item Maîtrise de la lecture de documents statistiques (tableaux, graphiques) et de textes juridiques simplifiés.
\item Compétences en expression écrite : structurer un paragraphe argumentatif (thèse, arguments, exemples).
\end{itemize}
\end{prerequis}

\newpage

\coursec{Fondements, caractéristiques et classification des droits de l'homme}

\begin{exemplebox}[Situation d'accroche : Deux visages de l'injustice au Cameroun]
\textbf{À Douala, quartier Bonabéri.} Jean, 17 ans, apprenti menuisier dans un atelier informel, réclame depuis des semaines le port de lunettes de protection et de gants pour manipuler la scie circulaire. Son patron, excédé par ces « exigences », le licencie sur-le-champ, sans préavis, sans certificat de travail, sans indemnité. Jean se retrouve sans revenu, sans papier, sans voix.

\textbf{À Yaoundé.} L'association « Femmes Debout » dépose une demande officielle pour louer la salle des fêtes de la mairie afin d'organiser une causerie éducative sur les mariages précoces et les grossesses en milieu scolaire. Le maire adjoint refuse net, arguant d'un « trouble à l'ordre public » non précisé, alors que la même salle a été accordée la semaine précédente à un meeting politique.

Ces deux situations, banales en apparence, révèlent une même racine : \textbf{l'arbitraire face à la dignité humaine}. Jean voit bafoués son droit au travail dans des conditions justes et sa sécurité au travail (droits économiques et sociaux). Les femmes de l'association voient entraver leur liberté de réunion et d'expression (droits civils et politiques).

\medskip
\textbf{Problématique centrale :} \emph{Comment, concrètement, un citoyen camerounais peut-il identifier ses droits, en comprendre la portée juridique et saisir les mécanismes de protection pour transformer l'indignation en action citoyenne efficace ?}
\end{exemplebox}

\courssub{Définition : des prérogatives inhérentes à la personne humaine}

\begin{definition}[Droits de l'homme]
Les \cle{droits de l'homme} (ou droits humains) sont l'ensemble des \textbf{prérogatives attachées à la personne humaine du seul fait de sa nature}, reconnues par la conscience universelle et \textbf{garanties par le droit} (national et international). Ils ne sont pas des faveurs accordées par l'État, mais des attributs inhérents à la dignité humaine.
\end{definition}

\textbf{Trois idées structurent cette définition :}
\begin{enumerate}
    \item \textbf{L'universalité du titulaire :} tout être humain, sans distinction de race, de sexe, de langue, de religion, d'opinion politique, d'origine nationale ou sociale, de fortune ou de naissance (art. 2 de la DUDH).
    \item \textbf{La source :} la dignité inhérente à la personne humaine (Préambule de la DUDH ; Préambule de la Constitution camerounaise de 1996). L'État ne \emph{crée} pas le droit, il le \emph{reconnaît} et s'oblige à le \emph{garantir}.
    \item \textbf{La garantie juridique :} un droit sans possibilité de recours (accès à un juge, réparation) n'est qu'une déclaration morale. La justiciabilité est la condition de l'effectivité.
\end{enumerate}

\begin{attention}[Piège fréquent : droits de l'homme et droits du citoyen]
Ne confondez pas :
\begin{itemize}
    \item \textbf{Droits de l'homme :} universels, attachés à la \cle{personne humaine} (ex. : droit à la vie, interdiction de la torture, liberté d'opinion). Un étranger au Cameroun en bénéficie aussi.
    \item \textbf{Droits du citoyen :} liés à la \cle{nationalité} et au statut politique (ex. : droit de vote, éligibilité, accès à la fonction publique). Un étranger ne vote pas aux élections camerounaises.
\end{itemize}
\textbf{Astuce Probatoire :} dans une dissertation, précisez toujours : « Les droits de l'homme sont le socle ; les droits du citoyen en sont une déclinaison politique réservée aux nationaux. »
\end{attention}

\courssub{Les quatre caractéristiques fondamentales}

Ces caractéristiques, solennellement réaffirmées par la \textbf{Conférence mondiale sur les droits de l'homme de Vienne (1993)}, forment le « code génétique » des droits humains. Elles s'appliquent \textbf{indissociablement}.

\begin{popfigure}
\begin{center}
\begin{tikzpicture}[
    box/.style={draw=popblue, thick, rounded corners, fill=popblueL, minimum width=4.2cm, minimum height=1.6cm, align=center, font=\small},
    arrow/.style={-Stealth, thick, popdark}
]
\node[box, align=center] (uni) at (0,0){\textbf{UNIVERSALITÉ}\\\emph{Valables pour tous, partout}};
\node[box, fill=popgreenL, draw=popgreen, align=center] (indiv) at (6,0){\textbf{INDIVISIBILITÉ}\\\emph{Un bloc indissociable}};
\node[box, fill=poporangeL, draw=poporange, align=center] (inter) at (0,-3){\textbf{INTERDÉPENDANCE}\\\emph{Chaque droit soutient les autres}};
\node[box, fill=poppurpleL, draw=poppurple, align=center] (inal) at (6,-3){\textbf{INALIÉNABILITÉ}\\\emph{On ne peut y renoncer}};

\draw[arrow] (uni) -- (indiv);
\draw[arrow] (indiv) -- (inal);
\draw[arrow] (inal) -- (inter);
\draw[arrow] (inter) -- (uni);
\draw[arrow, dashed] (uni) -- (inal);
\draw[arrow, dashed] (indiv) -- (inter);
\end{tikzpicture}
\end{center}
\legende{Les quatre piliers indissociables des droits humains (Conférence de Vienne, 1993).}
\end{popfigure}

\textbf{1. Universalité.}\par
Les droits s'appliquent à \emph{tous les êtres humains}, en tout lieu, sans discrimination. C'est le sens de l'article 1\ier{} de la DUDH : « Tous les êtres humains naissent libres et égaux en dignité et en droits. » Au Cameroun, le Préambule de la Constitution proclame « l'attachement du peuple camerounais aux libertés fondamentales inscrites dans la Déclaration universelle des droits de l'homme ».

\textbf{2. Indivisibilité.}\par
On ne peut pas « faire le marché » : accepter les droits civils (vote, presse) et refuser les droits économiques (santé, éducation), ou l'inverse. L'ensemble forme un \textbf{tout cohérent}.

\begin{propriete}[Indivisibilité : le bloc de la dignité]
La violation d'une catégorie de droits entraîne l'érosion des autres.
\textbf{Exemple concret :} un enfant non scolarisé (violation du droit à l'éducation) devient un adulte analphabète. Il ne peut pas lire la Constitution, comprendre un bulletin de vote ni saisir un tribunal (atteinte aux droits civils et politiques). La pauvreté extrême \emph{désarme} le citoyen de ses droits civils.
\end{propriete}

\textbf{3. Interdépendance.}\par
Au-delà du lien logique (indivisibilité), il existe un \textbf{lien fonctionnel} : la réalisation d'un droit \emph{conditionne} l'exercice d'un autre.
\begin{itemize}
    \item \textbf{Droit à la santé $\rightarrow$ droit à la vie :} sans accès aux soins, le droit à la vie est illusoire.
    \item \textbf{Liberté d'expression $\rightarrow$ droit au développement :} la transparence et le contrôle citoyen de l'action publique sont indispensables pour orienter les ressources vers le développement.
\end{itemize}

\textbf{4. Inaliénabilité.}\par
Nul ne peut \emph{renoncer} à ses droits, ni les \emph{aliéner} (les vendre, les abandonner). Un employeur ne peut pas faire signer une clause de « renonciation au droit de grève » dans un contrat de travail : une telle clause est \textbf{nulle}.

\begin{savaistu}[Le « relativisme culturel » : un faux débat]
Certains soutiennent que les droits humains seraient une « importation occidentale » incompatible avec les « valeurs africaines » (solidarité communautaire, respect des aînés). La \textbf{Charte africaine des droits de l'homme et des peuples (1981)} a tranché : elle \emph{intègre} les devoirs de l'individu envers la communauté (art. 27-29) et les droits des peuples (art. 19-24) \emph{à l'intérieur} du cadre universel. Le Cameroun, en ratifiant la Charte et en inscrivant les droits humains dans le Préambule de sa Constitution de 1996, a opéré cette synthèse : universalité des principes, contextualisation des mises en œuvre.
\end{savaistu}

\courssub{Classification : les trois générations de droits (Karel Vasak, 1979)}

La doctrine (Karel Vasak, juriste franco-tchèque, ancien directeur à l'UNESCO) propose une lecture \textbf{historique et logique} en trois « générations », inspirée de la devise \emph{Liberté, Égalité, Fraternité}. Attention : il ne s'agit pas de générations qui s'effacent l'une l'autre, mais de \textbf{strates cumulatives}.

\begin{popfigure}
\begin{center}
\begin{tikzpicture}[
    gen/.style={draw=popblue, very thick, fill=popblueL, rounded corners, minimum width=5.2cm, minimum height=1.3cm, align=center, font=\small\bfseries},
    ex/.style={draw=popmuted, fill=popcream, rounded corners, align=left, font=\footnotesize, text width=6.2cm},
    arrow/.style={-Stealth, thick, popdark}
]
\node[gen, align=center] (g1) at (0,0){1\textsuperscript{re} GÉNÉRATION\\ Droits civils et politiques};
\node[ex, right=0.6cm of g1, anchor=west] (e1) {Liberté d'aller et venir ; droit à la vie et à l'intégrité physique ; liberté d'opinion et de presse ; procès équitable ; droit de vote.};
\draw[arrow] (g1) -- (e1);

\node[gen, draw=popgreen, fill=popgreenL, align=center] (g2) at (0,-2.6){2\textsuperscript{e} GÉNÉRATION\\ Droits économiques, sociaux et culturels};
\node[ex, right=0.6cm of g2, anchor=west] (e2) {Droit au travail et à des conditions justes ; droit à la santé ; droit à l'éducation ; droit au logement ; sécurité sociale.};
\draw[arrow] (g2) -- (e2);

\node[gen, draw=poporange, fill=poporangeL, align=center] (g3) at (0,-5.2){3\textsuperscript{e} GÉNÉRATION\\ Droits de solidarité (droits des peuples)};
\node[ex, right=0.6cm of g3, anchor=west] (e3) {Droit au développement ; droit à un environnement sain ; droit à la paix ; patrimoine commun de l'humanité.};
\draw[arrow] (g3) -- (e3);
\end{tikzpicture}
\end{center}
\legende{Les trois générations de droits selon Karel Vasak (1979), avec exemples.}
\end{popfigure}

\textbf{Analyse comparative des trois générations}\par

\begin{center}
\begin{tabularx}{\linewidth}{|l|X|X|X|}
\hline
\rowcolor{popblueL}
\textbf{Critère} & \textbf{1\textsuperscript{re} génération} & \textbf{2\textsuperscript{e} génération} & \textbf{3\textsuperscript{e} génération} \\
\hline
\textbf{Devise} & Liberté & Égalité & Fraternité / Solidarité \\
\hline
\textbf{Objet} & Protéger la personne contre l'arbitraire de l'État & Assurer à chacun une vie digne (biens et services) & Gérer les biens communs de l'humanité \\
\hline
\textbf{Titulaire} & Individu & Individu et groupes (travailleurs, familles) & Peuples, humanité, générations futures \\
\hline
\textbf{Obligation de l'État} & \textbf{Respecter} (s'abstenir de violer) & \textbf{Réaliser} (agir : lois, budget, infrastructures) & \textbf{Coopérer} (solidarité nationale et internationale) \\
\hline
\textbf{Textes clés} & DUDH (art. 3-21), PIDCP & DUDH (art. 22-27), PIDESC & Charte africaine (art. 19-24) \\
\hline
\textbf{Exemples au Cameroun} & Liberté d'association et de réunion (loi n\textsuperscript{o} 90/053 du 19/12/1990) & Gratuité de l'école primaire publique ; couverture sanitaire & Loi-cadre n\textsuperscript{o} 96/12 du 05/08/1996 relative à la gestion de l'environnement \\
\hline
\end{tabularx}
\end{center}

\begin{attention}[Erreur fréquente au Probatoire]
Ne dites pas que les générations « se succèdent et se remplacent ». Les trois générations \textbf{coexistent} et se renforcent mutuellement : c'est une conséquence directe de l'indivisibilité et de l'interdépendance.
\end{attention}

\courssub{Le cadre juridique au Cameroun : la place des traités dans la hiérarchie des normes}

Le Cameroun a construit une architecture juridique dans laquelle le droit international des droits de l'homme occupe une place élevée.

\begin{propriete}[Article 45 de la Constitution : la supériorité des traités sur les lois]
\begin{quote}
\emph{« Les traités et accords internationaux régulièrement ratifiés ou approuvés ont, dès leur publication, une autorité supérieure à celle des lois, sous réserve, pour chaque accord ou traité, de son application par l'autre partie. »}
\end{quote}
\textbf{Conséquence directe :} le PIDCP, le PIDESC, la Charte africaine, ratifiés par le Cameroun, sont \textbf{supérieurs aux lois nationales}. Un juge camerounais peut écarter une loi interne contraire à un traité ratifié.
\end{propriete}

\begin{exemplebox}[Le « bloc de constitutionnalité » camerounais]
Le Préambule de la Constitution du 18 janvier 1996 dispose :
\begin{quote}
\emph{« Le peuple camerounais [...] affirme son attachement aux libertés fondamentales inscrites dans la Déclaration universelle des droits de l'homme, la Charte des Nations Unies, la Charte africaine des droits de l'homme et des peuples et tous les instruments juridiques internationaux dûment ratifiés y relatifs. »}
\end{quote}
Ce Préambule a \textbf{valeur constitutionnelle} : il fait partie intégrante de la Constitution. Il forme, avec l'article 45, le \cle{bloc de constitutionnalité} : tout texte de loi doit lui être conforme.
\end{exemplebox}

\begin{savaistu}[La hiérarchie des normes au Cameroun]
De haut en bas : (1) la \textbf{Constitution} (avec son Préambule) ; (2) les \textbf{traités internationaux} ratifiés et publiés (art. 45) ; (3) les \textbf{lois} votées par le Parlement ; (4) les \textbf{règlements} (décrets, arrêtés). Chaque norme doit respecter celles qui lui sont supérieures : c'est le principe de \cle{hiérarchie des normes}.
\end{savaistu}

\courssub{Droits absolus et droits relatifs : le régime des restrictions}

Tous les droits ne supportent pas le même régime juridique. La distinction est essentielle pour analyser une situation concrète.

\begin{center}
\begin{tabularx}{\linewidth}{|l|X|X|}
\hline
\rowcolor{popblueL}
\textbf{Critère} & \textbf{Droits ABSOLUS (non dérogables)} & \textbf{Droits RELATIFS (susceptibles de restrictions)} \\
\hline
\textbf{Nature} & Noyau dur de la dignité humaine & Libertés d'exercice, qui peuvent être encadrées \\
\hline
\textbf{Restriction possible ?} & \textbf{JAMAIS}, même en état d'urgence ou de guerre (art. 4 PIDCP) & \textbf{OUI}, mais strictement encadrée (test tripartite ci-dessous) \\
\hline
\textbf{Exemples} & Droit à la vie ; interdiction de la torture et des traitements inhumains ; interdiction de l'esclavage ; non-rétroactivité de la loi pénale ; liberté de pensée et de conscience & Liberté de mouvement ; liberté d'expression ; liberté de réunion et d'association ; droit de vote ; droit de propriété \\
\hline
\end{tabularx}
\end{center}

\begin{methode}[Le test tripartite de légitimité d'une restriction]
Pour qu'une restriction à un droit relatif soit \textbf{légitime}, elle doit remplir \textbf{trois conditions cumulatives} :
\begin{enumerate}
    \item \textbf{Légalité :} la restriction est prévue par une \textbf{loi} accessible, précise et non rétroactive. Une simple circulaire ou un arrêté municipal ne suffit pas.
    \item \textbf{Finalité légitime :} la loi poursuit un but reconnu (sécurité nationale, ordre public, santé publique, protection des droits d'autrui). Interdire une marche pour un réel risque de troubles est légitime ; l'interdire parce qu'elle critique le gouvernement ne l'est pas.
    \item \textbf{Nécessité et proportionnalité :} la mesure est \textbf{indispensable} (il n'existe pas de moyen moins restrictif) et \textbf{proportionnée} au but visé.
\end{enumerate}
\textbf{Charge de la preuve :} c'est à l'autorité qui restreint de démontrer que les trois conditions sont remplies.
\end{methode}

\begin{exoresolu}[Qualification juridique : le cas de l'apprenti menuisier de Douala]
\textbf{Faits :} Jean, apprenti, est licencié verbalement, sans préavis ni indemnité, pour avoir réclamé des équipements de protection individuelle (EPI).

\textbf{Travail à faire :} identifier les droits en cause, qualifier la violation et proposer une voie de recours.
\tcblower
\textbf{Solution détaillée :}
\begin{enumerate}
    \item \textbf{Droits en cause :}
        \begin{itemize}
            \item droit au travail dans des conditions justes et favorables (2\textsuperscript{e} génération ; art. 6-7 PIDESC ; Code du travail, loi n\textsuperscript{o} 92/007 du 14/08/1992) ;
            \item droit à la santé et à la sécurité au travail (2\textsuperscript{e} génération ; art. 12 PIDESC) ;
            \item droit à un recours effectif (1\textsuperscript{re} génération ; art. 2 PIDCP).
        \end{itemize}
    \item \textbf{Application du test tripartite au licenciement :}
        \begin{itemize}
            \item \emph{Légalité ?} Le Code du travail impose une procédure (motif réel et sérieux, préavis, indemnité). Le licenciement verbal est \textbf{illégal}.
            \item \emph{Finalité ?} Réclamer des EPI n'est pas une faute : le motif est \textbf{illégitime} (punition de l'exercice d'un droit).
            \item \emph{Proportionnalité ?} Le licenciement, sanction maximale, est \textbf{disproportionné} face à une simple demande de sécurité.
        \end{itemize}
    \item \textbf{Conclusion :} la violation est caractérisée. Jean peut saisir l'inspection du travail puis le tribunal compétent pour obtenir le paiement des salaires dus, une indemnité de licenciement abusif et des dommages-intérêts.
\end{enumerate}
\end{exoresolu}

\courssub{Obligations immédiates et réalisation progressive}

C'est une distinction technique majeure, source de nombreux malentendus.

\begin{propriete}[Deux types d'obligations de l'État]
\begin{itemize}
    \item \textbf{Droits civils et politiques (art. 2 PIDCP) :} l'État s'engage à \emph{respecter} et à \emph{assurer} ces droits \textbf{immédiatement}. Si une loi autorisait la torture, l'État violerait le Pacte \textbf{dès maintenant}, quel que soit son niveau de développement.
    \item \textbf{Droits économiques, sociaux et culturels (art. 2 PIDESC) :} l'État s'engage en vue d'une \textbf{réalisation progressive}, « au maximum de ses ressources disponibles ». Mais cela implique des \textbf{obligations immédiates} : adopter des lois et des politiques dès la ratification, ne pas discriminer, ne pas prendre de mesures délibérément régressives (ex. : supprimer la gratuité de l'école primaire sans alternative).
\end{itemize}
\end{propriete}

\begin{attention}[Piège classique : « les droits sociaux ne sont pas justiciables » — FAUX]
\begin{itemize}
    \item \textbf{Erreur :} « Le droit à la santé est progressif, donc on ne peut rien exiger du juge. »
    \item \textbf{Vérité :} le juge peut sanctionner l'absence de politique ou de plan (obligation de moyens), la \textbf{discrimination} dans l'accès aux services (interdite immédiatement), les \textbf{mesures régressives} non justifiées, et l'atteinte au \textbf{minimum vital} (nourriture, eau, soins essentiels, éducation primaire), exigible immédiatement.
\end{itemize}
\textbf{Exemple :} refuser l'inscription d'un enfant à l'école primaire publique parce que son acte de naissance est en cours d'établissement viole le droit à l'éducation et le principe de non-discrimination : l'accueil de l'enfant est immédiat, la régularisation administrative suit en parallèle.
\end{attention}

\courssub{Repères historiques : de 1789 à la Constitution camerounaise de 1996}

\begin{popfigure}
\begin{center}
\begin{tikzpicture}[
    event/.style={draw=popblue, thick, fill=popblueL, rounded corners, align=left, font=\footnotesize, text width=7.2cm},
    date/.style={font=\bfseries\small, anchor=east},
    dot/.style={circle, fill=popblue, minimum size=5pt, inner sep=0}
]
\draw[popdark, thick] (0,0.3) -- (0,-11.6);

\node[dot] at (0,0) {}; \node[date] at (-0.3,0) {1789};
\node[event, anchor=west] at (0.3,0) {Déclaration des droits de l'homme et du citoyen (France) : libertés individuelles, résistance à l'oppression.};

\node[dot] at (0,-1.6) {}; \node[date] at (-0.3,-1.6) {1945};
\node[event, anchor=west] at (0.3,-1.6) {Charte des Nations Unies (San Francisco) : promotion du respect des droits de l'homme.};

\node[dot] at (0,-3.2) {}; \node[date] at (-0.3,-3.2) {1948};
\node[event, anchor=west] at (0.3,-3.2) {\textbf{Déclaration universelle des droits de l'homme (DUDH)}, adoptée à Paris le 10 décembre 1948 : socle commun à tous les peuples.};

\node[dot] at (0,-4.8) {}; \node[date] at (-0.3,-4.8) {1966};
\node[event, anchor=west] at (0.3,-4.8) {Pactes de l'ONU : PIDCP et PIDESC (en vigueur en 1976), juridiquement contraignants pour les États parties.};

\node[dot] at (0,-6.4) {}; \node[date] at (-0.3,-6.4) {1981};
\node[event, anchor=west] at (0.3,-6.4) {\textbf{Charte africaine des droits de l'homme et des peuples} (Banjul) : originalité africaine — devoirs de l'individu et droits des peuples.};

\node[dot] at (0,-8) {}; \node[date] at (-0.3,-8) {1990};
\node[event, anchor=west] at (0.3,-8) {Cameroun : lois sur les libertés d'association, de réunion et de la communication sociale.};

\node[dot] at (0,-9.6) {}; \node[date] at (-0.3,-9.6) {1996};
\node[event, anchor=west] at (0.3,-9.6) {\textbf{Constitution du 18 janvier 1996} : Préambule (droits fondamentaux) et art. 45 (supériorité des traités).};

\node[dot] at (0,-11.2) {}; \node[date] at (-0.3,-11.2) {2004};
\node[event, anchor=west] at (0.3,-11.2) {Cameroun : loi n\textsuperscript{o} 2004/016 du 22/07/2004 portant création de la \textbf{Commission Nationale des Droits de l'Homme et des Libertés (CNDHL)}.};
\end{tikzpicture}
\end{center}
\legende{Chronologie des grands jalons : de la philosophie des Lumières à l'ancrage constitutionnel camerounais.}
\end{popfigure}

\courssub{Entraînement : classer, qualifier, argumenter}

\begin{exercice}[Cas pratiques de classification et de qualification]
Pour chacun des cas ci-dessous, en vous appuyant sur la méthode ci-après :
\begin{enumerate}
    \item Identifiez le ou les droits en cause (libellé précis).
    \item Classez-les par génération (1\textsuperscript{re}, 2\textsuperscript{e}, 3\textsuperscript{e}).
    \item Précisez s'il s'agit d'un droit absolu ou relatif.
    \item Indiquez l'obligation principale de l'État (respecter / protéger / réaliser).
    \item Citez une base textuelle applicable au Cameroun.
\end{enumerate}

\textbf{Cas A :} un journaliste est arrêté et détenu 72 heures sans être présenté à un juge, pour avoir publié une enquête sur la gestion de fonds publics.

\textbf{Cas B :} un quartier périphérique de Bamenda n'a plus accès à l'eau potable depuis six mois ; les forages sont en panne et aucune réparation n'est entreprise.

\textbf{Cas C :} une entreprise minière déverse des résidus toxiques dans une rivière, détruisant la pêche et les cultures des villages riverains.

\textbf{Cas D :} un projet de loi envisage d'interdire totalement le droit de grève à tous les fonctionnaires.
\end{exercice}

\begin{methode}[Classifier un droit en 5 étapes (fiche outil)]
\begin{enumerate}
    \item \textbf{Objet du droit :} de quoi s'agit-il ? (circuler librement ? se soigner ? vivre dans un environnement sain ?)
    \item \textbf{Titulaire :} qui en bénéficie ? (tout individu ? les travailleurs ? un peuple ? les générations futures ?)
    \item \textbf{Obligation de l'État :} s'abstenir (respecter) ? empêcher les tiers de nuire (protéger) ? mettre en œuvre des moyens (réaliser) ?
    \item \textbf{Référentiel textuel :} quel article de la Constitution, de la DUDH, du PIDCP, du PIDESC ou de la Charte africaine ?
    \item \textbf{Régime :} droit absolu (aucune restriction) ou relatif (test tripartite) ?
\end{enumerate}
\end{methode}

\begin{corrige}[Exercice — Cas pratiques]
\textbf{Cas A (journaliste détenu) :}
\begin{itemize}
    \item Droits : liberté d'expression et de la presse ; liberté et sécurité de la personne ; garanties d'une procédure équitable.
    \item Génération : \textbf{1\textsuperscript{re}} (droits civils et politiques).
    \item Nature : la liberté d'expression est \textbf{relative} (restrictions possibles mais encadrées) ; la détention sans présentation à un juge viole une garantie qui ne souffre aucune restriction arbitraire.
    \item Obligation : \textbf{respecter} (ne pas arrêter arbitrairement) et \textbf{protéger} (le juge contrôle l'action de la police).
    \item Base : Préambule de la Constitution ; art. 9 et 19 PIDCP ; loi n\textsuperscript{o} 90/052 du 19/12/1990 sur la communication sociale.
\end{itemize}

\textbf{Cas B (eau à Bamenda) :}
\begin{itemize}
    \item Droit : droit à un niveau de vie suffisant, incluant l'accès à l'eau potable ; droit à la santé.
    \item Génération : \textbf{2\textsuperscript{e}} (DESC).
    \item Nature : droit \textbf{relatif} à réalisation progressive, mais le \textbf{minimum vital} (eau potable) est exigible immédiatement, sans discrimination entre quartiers.
    \item Obligation : \textbf{réaliser} (investir, réparer) et \textbf{protéger} (contrôler le service de l'eau).
    \item Base : art. 11 et 12 PIDESC ; loi n\textsuperscript{o} 98/005 du 14/04/1998 portant régime de l'eau.
\end{itemize}

\textbf{Cas C (pollution minière) :}
\begin{itemize}
    \item Droits : droit à un environnement sain ; droit au développement ; droits à l'alimentation, à l'eau et à la vie des riverains.
    \item Génération : \textbf{3\textsuperscript{e}} (environnement, développement), avec atteinte aux 1\textsuperscript{re} (vie) et 2\textsuperscript{e} (alimentation, santé) générations — illustration de l'indivisibilité.
    \item Nature : droit à l'environnement \textbf{relatif}, mais le devoir de prévention des dommages graves est immédiat ; le droit à la vie est \textbf{absolu}.
    \item Obligation : \textbf{protéger} (contrôler l'entreprise, sanctionner, faire réparer) et \textbf{réaliser} (études d'impact, législation environnementale).
    \item Base : art. 24 Charte africaine ; loi-cadre n\textsuperscript{o} 96/12 du 05/08/1996 sur la gestion de l'environnement.
\end{itemize}

\textbf{Cas D (grève des fonctionnaires) :}
\begin{itemize}
    \item Droit : liberté syndicale et droit de grève.
    \item Génération : \textbf{1\textsuperscript{re}} (liberté d'association) et \textbf{2\textsuperscript{e}} (droits des travailleurs).
    \item Nature : droit \textbf{relatif} : des restrictions sont admises pour les agents exerçant des fonctions essentielles de sécurité, mais une \textbf{interdiction générale et absolue} serait disproportionnée et donc illégitime (échec au test tripartite).
    \item Obligation : \textbf{respecter} (ne pas interdire purement et simplement) et \textbf{protéger} (garantir l'exercice effectif).
    \item Base : art. 22 PIDCP ; art. 8 PIDESC ; Code du travail (loi n\textsuperscript{o} 92/007 du 14/08/1992).
\end{itemize}
\end{corrige}

\begin{aretenir}[Synthèse de la section : ce qu'il faut maîtriser pour le Probatoire]
\begin{itemize}
    \item \textbf{Définition :} prérogatives inhérentes à la personne humaine, garanties par le droit national et international.
    \item \textbf{4 caractéristiques (Vienne, 1993) :} universalité, indivisibilité, interdépendance, inaliénabilité.
    \item \textbf{3 générations (Vasak) :} 1\textsuperscript{re} — libertés civiles et politiques (obligation de \emph{respecter}) ; 2\textsuperscript{e} — droits économiques, sociaux et culturels (obligation de \emph{réaliser} progressivement, avec minimum vital immédiat) ; 3\textsuperscript{e} — droits de solidarité (obligation de \emph{coopérer}).
    \item \textbf{Cameroun :} Préambule + art. 45 de la Constitution de 1996 = supériorité des traités ratifiés sur les lois.
    \item \textbf{Absolu / relatif :} droits absolus (vie, interdiction de la torture et de l'esclavage, non-rétroactivité pénale, liberté de conscience) : jamais de restriction. Droits relatifs : test tripartite (loi + finalité légitime + nécessité/proportionnalité).
    \item \textbf{Justiciabilité des DESC :} réelle — non-discrimination immédiate, interdiction des régressions, minimum vital exigible.
\end{itemize}
\end{aretenir}

\coursec{Le cadre juridique : Hiérarchie des normes et justiciabilité au Cameroun}

Après avoir identifié les fondements et la classification des droits de l'homme (générations de droits, universalité, indivisibilité), il est indispensable de comprendre comment ces droits s'insèrent dans l'ordre juridique camerounais et deviennent \cle{opposables}, c'est-à-dire invoquables devant un juge. Un droit proclamé sans garantie juridictionnelle effective reste une simple déclaration d'intention. Cette section analyse l'architecture normative qui donne force de loi aux droits humains au Cameroun, les mécanismes de contrôle de leur respect par le juge national, et les voies de recours ouvertes aux justiciables, tant sur le plan interne qu'international.

\courssub{La pyramide des normes camerounaise : de la Constitution aux actes administratifs}

Toute norme inférieure tire sa validité de la norme supérieure qui l'autorise. En cas de conflit, la norme supérieure l'emporte. Cette \cle{hiérarchie des normes} détermine la valeur juridique des instruments protégeant les droits de l'homme au Cameroun.

\begin{popfigure}
\begin{tikzpicture}[
    norm/.style={draw=popdark, thick, rounded corners=4pt, fill=popblueL, text width=6.5cm, align=center, font=\small\bfseries, minimum height=1.1cm},
    arrow/.style={-Stealth, thick, popdark},
    ctrl/.style={draw=poporange, thick, dashed, rounded corners=3pt, fill=poporangeL, text width=4.2cm, align=center, font=\scriptsize, minimum height=0.9cm}
]
\node[norm, align=center] (const){CONSTITUTION (18 janvier 1996, révisée en 2008)\\ Sommet de la pyramide};
\node[norm, below=0.5cm of const, align=center] (trait){TRAITÉS INTERNATIONAUX RATIFIÉS (art. 45 Const.)\\ DUDH, PIDCP, PIDESC, CEDAW, Charte africaine...};
\node[norm, below=0.5cm of trait, align=center] (loiorg){LOIS ORGANIQUES\\ (majorité qualifiée, contrôle obligatoire)};
\node[norm, below=0.5cm of loiorg, align=center] (loiord){LOIS ORDINAIRES\\ (votées par le Parlement)};
\node[norm, below=0.5cm of loiord, align=center] (decret){DÉCRETS PRÉSIDENTIELS\\ (domaine réglementaire)};
\node[norm, below=0.5cm of decret, align=center] (arrete){ARRÊTÉS ET ACTES ADMINISTRATIFS\\ (mesures d'application)};

\draw[arrow] (const) -- (trait) node[midway, right, xshift=3mm, font=\scriptsize, popgreen] {supériorité (art. 45)};
\draw[arrow] (trait) -- (loiorg);
\draw[arrow] (loiorg) -- (loiord);
\draw[arrow] (loiord) -- (decret);
\draw[arrow] (decret) -- (arrete);

\node[ctrl, right=1.2cm of const, align=center] (cc){CONSEIL CONSTITUTIONNEL\\ contrôle de constitutionnalité des lois};
\node[ctrl, right=1.2cm of loiord, align=center] (ja){JUGE ADMINISTRATIF\\ contrôle de légalité des actes administratifs};

\draw[arrow, poporange, dashed] (cc.west) -- (const.east);
\draw[arrow, poporange, dashed] (ja.west) -- (loiord.east);
\end{tikzpicture}
\legende{Pyramide des normes au Cameroun. La Constitution domine ; les traités régulièrement ratifiés (art. 45) ont une autorité supérieure à celle des lois. Le Conseil Constitutionnel contrôle la conformité des lois à la Constitution ; le juge administratif contrôle la légalité des actes administratifs.}
\end{popfigure}

\begin{definition}[Hiérarchie des normes dans l'ordre juridique camerounais]
L'ordre juridique camerounais est structuré en niveaux hiérarchiques :
\begin{enumerate}
    \item \textbf{La Constitution} du 18 janvier 1996 (révisée en 2008) : norme suprême. Son Préambule affirme l'attachement du peuple camerounais aux libertés fondamentales inscrites dans la Déclaration universelle des droits de l'homme (DUDH), la Charte des Nations unies et la Charte africaine des droits de l'homme et des peuples (CADHP).
    \item \textbf{Les traités et accords internationaux régulièrement ratifiés} (art. 45 Const.) : ils ont une autorité supérieure à celle des lois, sous réserve de réciprocité.
    \item \textbf{Les lois organiques} : votées selon une procédure renforcée, soumises obligatoirement au contrôle du Conseil Constitutionnel avant promulgation.
    \item \textbf{Les lois ordinaires} : votées par le Parlement (Assemblée nationale et Sénat).
    \item \textbf{Les décrets} : actes du Président de la République (domaine réglementaire).
    \item \textbf{Les arrêtés et décisions} : actes du Premier Ministre, des ministres et des autres autorités administratives (mesures d'exécution).
\end{enumerate}
Toute norme inférieure contraire à une norme supérieure est inconstitutionnelle ou illégale et peut être écartée ou annulée par le juge compétent.
\end{definition}

\begin{propriete}[Supériorité des traités ratifiés — Article 45 de la Constitution]
\textbf{Énoncé :} « Les traités ou accords internationaux régulièrement approuvés ou ratifiés ont, dès leur publication, une autorité supérieure à celle des lois, sous réserve, pour chaque traité ou accord, de son application par l'autre partie. »
\textbf{Portée :} c'est le principe de \cle{supériorité des traités sur les lois nationales}, même postérieures. Le juge camerounais doit écarter la loi nationale incompatible avec un traité régulièrement ratifié et publié : c'est le \cle{contrôle de conventionnalité}.
\textbf{Conditions :} un traité est « régulièrement ratifié » lorsqu'il a été signé, autorisé par le Parlement (loi d'autorisation), ratifié par le Président de la République, puis publié au Journal Officiel.
\end{propriete}

\begin{attention}[Piège fréquent : signature $\neq$ ratification]
La simple \textbf{signature} d'un traité par un représentant de l'État n'engage pas définitivement l'État sur le plan juridique : elle manifeste seulement l'intention de se lier. Seule la \textbf{ratification}, suivie de la \textbf{publication} au Journal Officiel, rend le traité obligatoire et opposable dans l'ordre interne. Dans une copie, vérifiez toujours que le traité invoqué a bien été ratifié et publié par le Cameroun.
\end{attention}

\begin{savaistu}[Le Préambule : une norme à part entière]
Le Préambule de la Constitution camerounaise de 1996 ne se limite pas à une déclaration politique : il fait \textbf{partie intégrante de la Constitution} et possède la même valeur juridique que ses articles. Les droits et libertés qu'il proclame (droit à la vie, à la liberté, à l'éducation, à un environnement sain, liberté d'expression, de réunion...) sont donc des \textbf{normes constitutionnelles} que le juge et le citoyen peuvent invoquer. C'est l'un des fondements majeurs de la protection des droits humains au Cameroun.
\end{savaistu}

\courssub{Le contrôle de constitutionnalité : le Conseil Constitutionnel}

Au sommet de la pyramide, le \cle{Conseil Constitutionnel} est le gardien de la suprématie de la Constitution. Il veille à ce que les lois respectent les droits et libertés garantis par la Constitution et son Préambule.

\begin{methode}[Les voies de saisine du Conseil Constitutionnel]
\begin{enumerate}
    \item \textbf{Saisine obligatoire (avant promulgation) :} les lois organiques et les règlements intérieurs de l'Assemblée nationale et du Sénat lui sont obligatoirement soumis.
    \item \textbf{Saisine facultative (avant promulgation) :} pour les lois ordinaires, peuvent saisir le Conseil : le Président de la République, le Premier Ministre, le Président de l'Assemblée nationale, le Président du Sénat, ainsi qu'\textbf{un dixième (1/10) des députés ou des sénateurs}.
    \item \textbf{Exception d'inconstitutionnalité (en cours de procès) :} devant une juridiction, une partie peut soutenir qu'une loi porte atteinte aux droits et libertés garantis par la Constitution ; la juridiction sursoit à statuer et renvoie la question au Conseil Constitutionnel. Le justiciable ne saisit donc pas directement le Conseil : c'est le juge qui le fait.
    \item \textbf{Règlement des conflits :} le Conseil arbitre également les conflits de compétence entre les institutions de l'État.
\end{enumerate}
\end{methode}

\begin{exemplebox}[Application : le contrôle d'une loi électorale]
Avant la promulgation d'une nouvelle loi électorale, le Président de la République ou un dixième des parlementaires peuvent la déférer au Conseil Constitutionnel pour vérifier sa conformité aux principes constitutionnels (égalité du suffrage, liberté de vote). Si le Conseil déclare une disposition contraire à la Constitution, celle-ci ne peut entrer en vigueur : la loi doit être modifiée avant promulgation.
\end{exemplebox}

\courssub{La justiciabilité des droits : droits directement applicables et droits programmatiques}

Tous les droits proclamés par la Constitution ou les traités n'ont pas la même portée devant le juge. Cette distinction est essentielle pour comprendre pourquoi certains droits peuvent être immédiatement sanctionnés et d'autres non.

\begin{definition}[Justiciabilité]
La \cle{justiciabilité} est la possibilité pour un justiciable de saisir un juge indépendant et impartial pour faire constater et sanctionner la violation d'un droit et obtenir réparation. Elle suppose que la norme invoquée crée une obligation suffisamment précise et exigible pour les autorités publiques.
\end{definition}

\begin{propriete}[Classification des droits selon leur justiciabilité]
\begin{itemize}
    \item \textbf{Droits directement applicables (auto-exécutoires) :} leur contenu est précis et ne nécessite aucune mesure d'application pour produire ses effets. Le juge peut les appliquer \textbf{directement}. \emph{Exemples :} droit à la vie, interdiction de la torture et des traitements inhumains, liberté d'aller et venir, non-discrimination, droit à un procès équitable, droit d'être informé des motifs de son arrestation.
    \item \textbf{Droits programmatiques (de mise en œuvre progressive) :} ils énoncent un objectif général (éducation, santé, logement, travail, environnement sain) dont la réalisation dépend de lois, de budgets et de politiques publiques. Le juge ne peut pas se substituer au législateur pour définir ces politiques, mais il peut sanctionner :
    \begin{itemize}
        \item l'\textbf{absence totale de mesure} (carence fautive de l'État) ;
        \item le caractère \textbf{manifestement discriminatoire} des mesures prises ;
        \item les \textbf{mesures de régression} qui remettent en cause le niveau de protection déjà acquis.
    \end{itemize}
\end{itemize}
\end{propriete}

\begin{exoresolu}[Qualifier la justiciabilité de deux situations]
\textbf{Énoncé :} M. X est arrêté par la police et maintenu en garde à vue au-delà de la durée légale, sans être présenté à un magistrat ni informé de ses droits. Il saisit le juge judiciaire aux fins de libération. Parallèlement, une association saisit le juge administratif pour exiger la construction d'un hôpital dans son arrondissement, en invoquant le droit à la santé.
\tcblower
\textbf{Solution détaillée :}
\begin{enumerate}
    \item \textbf{Cas de M. X (détention irrégulière) :} le droit à la liberté individuelle et le droit d'être informé des motifs de son arrestation sont des droits \textbf{directement applicables} (Préambule de la Constitution, art. 9 du PIDCP ratifié par le Cameroun). Le juge judiciaire peut ordonner la mise en liberté sur le fondement direct de ces textes. \textbf{Recours recevable et fondé.}
    \item \textbf{Cas de l'association (hôpital) :} le droit à la santé est un droit \textbf{programmatique}. Le juge ne peut pas ordonner la construction d'un hôpital, car cela relève des choix budgétaires et de la politique publique de santé. Il pourrait seulement sanctionner une \textbf{carence totale} de l'État ou une discrimination manifeste entre les populations. \textbf{Recours en principe irrecevable sur la demande de construction.}
\end{enumerate}
\textbf{Conclusion :} la qualification du droit (directement applicable ou programmatique) détermine le succès du recours.
\end{exoresolu}

\courssub{Les recours internes : architecture contentieuse administrative et judiciaire}

Face à une violation de ses droits, le justiciable camerounais dispose d'une panoplie de recours. Le choix de la voie de droit dépend de l'\textbf{auteur de la violation} (administration ou particulier) et de la \textbf{nature du droit violé}.

\begin{center}
\begin{tabularx}{\linewidth}{|l|X|X|X|}
\hline
\rowcolor{popblueL}
\textbf{Voie de recours} & \textbf{Objet} & \textbf{Conditions essentielles} & \textbf{Issue possible} \\
\hline
\textbf{Recours gracieux} & Demande adressée à l'auteur de l'acte (ministre, maire, directeur) pour qu'il retire ou modifie sa décision. & Intérêt personnel et légitime à agir. Aucune forme particulière exigée. & Retrait ou maintien de l'acte. \\
\hline
\textbf{Recours hiérarchique} & Demande adressée au supérieur de l'auteur de l'acte (ex. : ministre $\rightarrow$ Premier Ministre). & Mêmes conditions que le recours gracieux. & Annulation, modification ou rejet. \\
\hline
\textbf{Recours devant le juge administratif} & \textbf{Excès de pouvoir} (annulation d'un acte administratif illégal) ; \textbf{responsabilité administrative} (réparation du dommage causé par l'administration). & Qualité et intérêt à agir ; recours formé dans le \textbf{délai légal de 2 mois} à compter de la notification ou de la publication de l'acte. & Annulation de l'acte ; dommages et intérêts. \\
\hline
\textbf{Référé devant le juge administratif} & Mesures d'urgence face à une atteinte grave et manifestement illégale à une liberté fondamentale (référé-liberté) ou à un risque de dommage irréversible (référé-suspension). & Urgence ; atteinte grave et manifestement illégale. & Ordonnance rapide : suspension de l'acte, mesures conservatoires. \\
\hline
\textbf{Appel et cassation} & La Cour d'appel réexamine les décisions des tribunaux ; la Cour Suprême statue en cassation (violation de la loi). & Respect des délais d'appel et de pourvoi. & Réformation ou annulation de la décision attaquée. \\
\hline
\textbf{Juge judiciaire} & Violations commises par des \textbf{particuliers} ou des agents agissant hors de leurs fonctions (violences, discriminations privées, atteintes à la liberté individuelle). & Saisine par requête ou citation ; en matière pénale, plainte ou citation directe. & Sanction pénale ; dommages et intérêts ; mesures de protection. \\
\hline
\end{tabularx}
\end{center}

\begin{methode}[Vérifier la recevabilité d'un recours contre un acte administratif]
Avant de saisir le juge administratif, suivez cette démarche en quatre étapes :
\begin{enumerate}
    \item \textbf{Qualité et intérêt à agir :} le requérant est-il directement concerné par l'acte ? Une association de défense des droits humains peut-elle défendre un intérêt collectif ?
    \item \textbf{Voies amiables :} un recours gracieux ou hiérarchique a-t-il été tenté ? (Recommandé, et parfois nécessaire, avant le contentieux.)
    \item \textbf{Délai :} le recours contentieux est-il introduit dans le \textbf{délai de 2 mois} ? Passé ce délai, le recours est forclos (irrecevable).
    \item \textbf{Compétence :} l'acte attaqué est-il bien un acte administratif relevant du juge administratif, et non un litige entre particuliers relevant du juge judiciaire ?
\end{enumerate}
\end{methode}

\begin{attention}[Les pièges de la procédure]
De nombreux recours échouent non pas parce que le droit invoqué n'existe pas, mais pour des raisons de \textbf{procédure} : délai de recours dépassé (forclusion), juridiction incompétente (juge judiciaire saisi à la place du juge administratif), absence d'intérêt à agir. Dans une copie d'examen, vérifiez toujours la recevabilité \textbf{avant} d'examiner le fond.
\end{attention}

\courssub{Les réserves et déclarations interprétatives du Cameroun}

La ratification d'un traité n'est pas toujours pure et simple. L'État peut formuler des \cle{réserves} (exclusion ou modification de l'effet juridique de certaines dispositions) ou des \cle{déclarations interprétatives} (précision de la portée qu'il entend donner à un article). Ces actes, faits lors de la ratification, modifient la portée du traité pour le Cameroun.

\begin{center}
\begin{tabularx}{\linewidth}{|l|X|X|}
\hline
\rowcolor{popblueL}
\textbf{Instrument} & \textbf{Réserves / déclarations du Cameroun} & \textbf{Impact sur la justiciabilité} \\
\hline
\textbf{CEDAW} (Convention sur l'élimination de toutes les formes de discrimination à l'égard des femmes, ratifiée en 1994) & Réserves portant notamment sur l'article 2 (politique d'élimination des discriminations) et l'article 16 (égalité dans le mariage et la famille), au motif de la conformité avec le droit interne et les coutumes. & Le juge ne peut pas fonder une décision uniquement sur ces articles réservés pour imposer une égalité totale en matière matrimoniale sans loi interne de mise en conformité. \\
\hline
\textbf{Protocole de Maputo} (droits des femmes en Afrique, ratifié par le Cameroun en 2009, entré en vigueur pour le Cameroun après dépôt) & Déclarations interprétatives, notamment sur l'âge du mariage et sur l'interruption médicale de grossesse, renvoyant à la législation nationale. & Les dispositions du droit interne (Code civil, Code pénal) continuent de s'appliquer sur ces points ; le juge ne peut pas invoquer le Protocole pour contraindre le législateur à modifier la loi. \\
\hline
\textbf{PIDCP} (Pacte international relatif aux droits civils et politiques, en vigueur pour le Cameroun depuis 1984) & Aucune réserve substantielle. & Pleine justiciabilité des droits civils et politiques (liberté, procès équitable, expression...). \\
\hline
\end{tabularx}
\end{center}

\begin{attention}[Effet des réserves sur le juge national]
Une réserve régulièrement formulée lie l'État sur le plan international : le traité s'applique \textbf{moins la disposition réservée}. Le juge camerounais ne peut pas appliquer la disposition réservée comme si la réserve n'existait pas ; il applique alors la loi interne. En cas de doute, il faut consulter l'instrument de ratification publié au Journal Officiel.
\end{attention}

\courssub{Ouverture vers l'international : les voies de recours au-delà du juge national}

Lorsque toutes les voies de recours internes ont été épuisées sans succès, le justiciable peut se tourner vers les mécanismes internationaux de protection. C'est la règle de l'\cle{épuisement préalable des voies de recours internes}.

\begin{propriete}[Les mécanismes internationaux accessibles aux Camerounais]
\begin{itemize}
    \item \textbf{Le Comité des droits de l'homme de l'ONU :} le Cameroun est partie au PIDCP et à son Premier Protocole facultatif, qui autorise les \textbf{communications individuelles}. Une victime peut donc saisir le Comité après épuisement des recours internes. Le Comité adopte des « constatations » (vues) dont l'État est invité à tirer les conséquences.
    \item \textbf{La Commission africaine des droits de l'homme et des peuples} (siège à Banjul, Gambie) : toute personne ou ONG peut lui adresser une communication. Elle adopte des recommandations, juridiquement non contraignantes mais dotées d'une forte autorité morale et politique.
    \item \textbf{La Cour africaine des droits de l'homme et des peuples} (siège à Arusha, Tanzanie) : elle rend des \textbf{arrêts contraignants}. Toutefois, la saisine directe par les individus et les ONG n'est possible que si l'État a fait la déclaration prévue à l'article 34(6) du Protocole. \textbf{Le Cameroun n'a pas déposé cette déclaration} : les individus camerounais ne peuvent donc pas saisir directement la Cour ; l'affaire doit passer par la Commission africaine, qui peut la renvoyer devant la Cour.
\end{itemize}
\end{propriete}

\begin{attention}[Erreur fréquente à éviter]
Ne confondez pas \textbf{Commission africaine} (recommandations, saisine ouverte à tous) et \textbf{Cour africaine} (arrêts contraignants, saisine directe des individus seulement contre les États ayant fait la déclaration de l'article 34(6) — ce qui n'est pas le cas du Cameroun). De même, seuls les États ayant ratifié le Protocole facultatif se rapportant au PIDCP — ce qui est le cas du Cameroun — peuvent voir leurs ressortissants saisir le Comité des droits de l'homme de l'ONU.
\end{attention}

\begin{exoresolu}[Stratégie de saisine : du local à l'international]
\textbf{Énoncé :} Une journaliste est condamnée à une peine de prison ferme pour « propagation de fausses nouvelles » (infraction prévue par le Code pénal camerounais) après la publication d'une enquête sur la gestion de fonds publics. Elle a exercé l'appel (confirmé) et le pourvoi en cassation (rejeté). Quelle stratégie contentieuse internationale recommandez-vous ?
\tcblower
\textbf{Solution détaillée :}
\begin{enumerate}
    \item \textbf{Vérifier l'épuisement des voies internes :} appel et cassation épuisés $\rightarrow$ la condition est remplie.
    \item \textbf{Identifier les organes accessibles :}
    \begin{itemize}
        \item \textbf{Comité des droits de l'homme (ONU) :} accessible, car le Cameroun est partie au Premier Protocole facultatif du PIDCP $\rightarrow$ communication individuelle possible.
        \item \textbf{Commission africaine des droits de l'homme et des peuples :} accessible (le Cameroun est partie à la Charte africaine) $\rightarrow$ communication possible, aboutissant à des recommandations.
        \item \textbf{Cour africaine :} saisine directe \textbf{impossible} (pas de déclaration de l'article 34(6) par le Cameroun) ; seule la Commission africaine pourrait renvoyer l'affaire devant la Cour.
    \end{itemize}
    \item \textbf{Construire l'argumentaire :} violation de la liberté d'expression (art. 19 PIDCP, art. 9 Charte africaine) ; peine disproportionnée ; restriction non nécessaire dans une société démocratique.
    \item \textbf{Respecter les délais :} saisine dans un délai raisonnable après la décision interne définitive.
\end{enumerate}
\textbf{Conclusion :} la voie la plus appropriée est la saisine du Comité des droits de l'homme de l'ONU et/ou de la Commission africaine, la saisine directe de la Cour africaine étant fermée aux individus camerounais.
\end{exoresolu}

\begin{aretenir}[Synthèse : cadre juridique et justiciabilité au Cameroun]
\begin{itemize}
    \item \textbf{Pyramide des normes :} Constitution $>$ traités régulièrement ratifiés (art. 45) $>$ lois organiques $>$ lois ordinaires $>$ décrets $>$ arrêtés. Le Préambule de la Constitution a valeur constitutionnelle.
    \item \textbf{Contrôle de constitutionnalité :} Conseil Constitutionnel, saisi avant promulgation (Président, Premier Ministre, présidents des chambres, 1/10 des parlementaires) ou par exception d'inconstitutionnalité renvoyée par le juge.
    \item \textbf{Justiciabilité :} distinction entre droits \textbf{directement applicables} (sanction immédiate par le juge) et droits \textbf{programmatiques} (contrôle limité : carence, discrimination, régression).
    \item \textbf{Recours internes :} recours gracieux et hiérarchique, puis juge administratif (excès de pouvoir, responsabilité, référés ; délai de 2 mois), appel, cassation ; juge judiciaire pour les violations commises par des particuliers.
    \item \textbf{Réserves :} les réserves du Cameroun (CEDAW, Protocole de Maputo) limitent l'application directe de certaines dispositions relatives aux droits des femmes.
    \item \textbf{Recours internationaux :} Comité des droits de l'homme de l'ONU (Protocole facultatif ratifié) et Commission africaine accessibles après épuisement des recours internes ; saisine directe de la Cour africaine non ouverte aux individus camerounais (absence de déclaration de l'article 34(6)).
\end{itemize}
\end{aretenir}

\coursec{Les institutions de protection : Architecture multiniveau (ONU, UA, National)}

\courssub{Transition : du texte à l'action}

Dans la section précédente, nous avons vu que le Cameroun, en ratifiant les grands traités internationaux, a fait entrer les droits de l'homme dans son ordre juridique interne (\emph{hiérarchie des normes}, \emph{justiciabilité}). Mais un droit sans juge ni gardien reste une « lettre morte ». Comment un élève de l'ENIET de Douala, un artisan du quartier Mvog-Ada à Yaoundé ou une association de défense des droits des femmes à Maroua peuvent-ils faire valoir leurs droits quand ils sont violés ? Par quelle porte frapper ? C'est tout l'objet de cette section : découvrir l'\textbf{architecture multiniveau de protection}, un système imbriqué où chaque niveau (universel, régional, sous-régional, national) a sa compétence, sa procédure et son rôle, dans une logique de \cle{subsidiarité} (le niveau supérieur n'intervient qu'après épuisement des recours internes).

\begin{experience}[Enquête de terrain : « À quelle porte frapper ? »]
\textbf{Contexte :} Un groupe d'élèves de Première Technique mène une mini-enquête dans leur quartier. Ils interrogent 20 personnes : « Si vos droits sont violés (licenciement abusif, refus d'inscription à l'école, violence policière, expulsion sans titre), vers qui vous tournez-vous en premier ? »\par
\textbf{Résultats observés :} 12 citent le chef de quartier ou le chef de village ; 5 citent la police/gendarmerie ; 2 citent un avocat ; 1 cite une ONG ; 0 ne cite la CNDHL, la CADHP ou l'ONU.\par
\textbf{Interprétation :} La méconnaissance des institutions formelles de protection pousse vers des recours informels (coutumiers, amicaux) qui, s'ils sont utiles, ne garantissent pas toujours le respect des standards juridiques (égalité, proportionnalité, droit à un recours effectif).\par
\textbf{Leçon :} Connaître l'architecture institutionnelle, c'est s'outiller pour choisir le bon niveau de saisine, au bon moment, avec la bonne procédure.
\end{experience}

\courssub{Niveau universel : Le système des Nations Unies (ONU)}

L'ONU constitue le « toit » commun. Deux piliers principaux assurent la protection :

\begin{itemize}
  \item \textbf{Le Conseil des Droits de l'Homme (CDH)} (siège à Genève, 47 États membres élus par l'AG pour 3 ans) :
    \begin{itemize}
      \item \textbf{L'Examen Périodique Universel (EPU)} : chaque État est passé au crible tous les 4,5 ans sur l'ensemble de ses obligations DH. Trois documents de base : rapport national, compilation ONU, rapport des parties prenantes (ONG, INDH). Aboutit à des \emph{recommandations} (acceptées ou notées).
      \item \textbf{Les Procédures Spéciales (Rapporteurs spéciaux, Experts indépendants, Groupes de travail)} : mandats thématiques (ex. torture, liberté d'expression, droits des migrants) ou pays. Ils mènent des enquêtes, effectuent des visites pays, envoient des \emph{communications} (lettres d'allégation, appels urgents) aux États.
    \end{itemize}
  \item \textbf{Les Comités des Traités (Treaty Bodies)} : 10 comités d'experts indépendants surveillent l'application de chaque traité principal. Pour le Cameroun, les plus pertinents sont :
    \begin{itemize}
      \item \textbf{Comité des Droits de l'Homme (CCPR)} : Pacte civil et politique (PIDCP) — \emph{communications individuelles} possibles (Protocole facultatif ratifié en 1984).
      \item \textbf{Comité des Droits Économiques, Sociaux et Culturels (CESCR)} : PIDESC — communications possibles (Protocole 2008 non ratifié par le Cameroun $\to$ pas de communications individuelles).
      \item \textbf{Comité CEDAW} : Convention élimination discrimination femmes — communications (Protocole facultatif ratifié 2004).
      \item \textbf{Comité CAT} : Convention contre la torture — communications (Art. 22 ratifié 1986).
    \end{itemize}
    \textbf{Procédure type :} État soumet rapport périodique $\to$ dialogue constructif $\to$ observations finales (recommandations). Si protocole facultatif ratifié : communication individuelle $\to$ recevabilité $\to$ fond $\to$ vues (non contraignantes mais autorité morale forte).
\end{itemize}

\begin{aretenir}[Repères pour l'examen]
\begin{itemize}
  \item EPU = examen \emph{global}, tous les 4,5 ans, \emph{tous} les États, \emph{politique} (recommandations).
  \item Comités des traités = examen \emph{spécialisé} par traité, périodique (4-5 ans), \emph{juridique} (observations finales).
  \item Communications individuelles = quasi-contentieux, seulement si État a ratifié le protocole facultatif correspondant.
  \item Cameroun : PIDCP (OUI), CEDAW (OUI), CAT (OUI) $\to$ communications possibles. PIDESC (NON pour Protocole) $\to$ pas de communications.
\end{itemize}
\end{aretenir}

\courssub{Niveau régional : Le système africain (Union Africaine)}

L'UA a construit un système original, à deux bras : une Commission (quasi-juridictionnelle, promotion/protection) et une Cour (juridictionnelle contraignante).

\begin{itemize}
  \item \textbf{Commission Africaine des Droits de l'Homme et des Peuples (CADHP)} (siège à Banjul, Gambie, 11 commissaires élus 6 ans) :
    \begin{itemize}
      \item \textbf{Communications} : toute personne, ONG, État peut saisir (Art. 55 Charte). Condition : \emph{épuisement des recours internes} (sauf si recours ineffectifs, indûment prolongés). Aboutit à des \emph{recommandations} (non contraignantes mais suivi par la Commission).
      \item \textbf{Rapports d'État} : Art. 62 Charte $\to$ rapport tous les 2 ans. Dialogue $\to$ observations finales.
      \item \textbf{Missions de promotion/protection} : visites pays, enquêtes sur situations graves, rapporteurs spéciaux (ex. droits des femmes, défenseurs DH, prisons).
    \end{itemize}
  \item \textbf{Cour Africaine des Droits de l'Homme et des Peuples (CADHP-Cour)} (siège à Arusha, Tanzanie, 11 juges) :
    \begin{itemize}
      \item \textbf{Contentieux} : saisine par États parties, Commission, \emph{ONG ayant statut d'observateur} et \emph{individus} \textbf{seulement si l'État a fait la déclaration Art. 34(6) Protocole}. \textbf{Le Cameroun n'a pas fait cette déclaration} $\to$ les individus/ONG ne peuvent pas saisir directement la Cour.
      \item \textbf{Avis consultatifs} : demandés par l'UA, ses organes, États, ONG reconnues.
      \item \textbf{Arrêts contraignants} : exécution surveillée par le Conseil Exécutif de l'UA.
    \end{itemize}
\end{itemize}

\begin{attention}[Piège fréquent : Confusion Commission / Cour]
\begin{itemize}
  \item \textbf{Commission (Banjul)} : accessible à \textbf{tous} (individus, ONG) pour communications, mais décisions = \emph{recommandations} non contraignantes.
  \item \textbf{Cour (Arusha)} : décisions = \emph{arrêts contraignants}, mais accès direct individus/ONG \textbf{bloqué} sans déclaration Art. 34(6) (Cameroun ne l'a pas faite).
  \item \textbf{Stratégie pratique :} Saisir d'abord la Commission (accessible). Si violation grave et État récalcitrant, la Commission \emph{peut} déférer l'affaire à la Cour (Art. 5(1) Protocole Cour).
\end{itemize}
\end{attention}

\courssub{Niveau sous-régional : CEMAC et CEEAC}

Deux cours suprêmes communautaires coexistent, avec des compétences partiellement chevauchantes mais des bases juridiques distinctes.

\begin{center}
\begin{tabularx}{\linewidth}{|l|X|X|}
\hline
\rowcolor{popblueL}
\textbf{Critère} & \textbf{Cour de Justice de la CEMAC} & \textbf{Cour de Justice de la CEEAC} \\
\hline
\textbf{Siège} & N'Djaména (Tchad) & N'Djaména (Tchad) \\
\hline
\textbf{Base juridique} & Traité CEMAC 1994, Protocole 2000 & Traité CEEAC 1983 (révisé), Protocole 2000/2005/2006 \\
\hline
\textbf{Compétence DH} & Droits \emph{communautaires} (liberté circulation, établissement, non-discrimination sur base nationalité) & Droits de l'homme \emph{généraux} (Protocole 2005/2006 étend compétence aux violations DH) \\
\hline
\textbf{Saisine individus/ONG} & Oui, pour violation droits communautaires & Oui, pour violations DH (après épuisement recours internes) \\
\hline
\textbf{Arrêts} & Contraignants pour États parties & Contraignants \\
\hline
\end{tabularx}
\end{center}

\begin{savaistu}[Deux Cours pour un même territoire ?]
Le Cameroun appartient aux deux organisations. La CEMAC (6 États : Cameroun, Tchad, RCA, Gabon, Guinée équatoriale, Congo) vise l'intégration monétaire/économique profonde. La CEEAC (11 États : CEMAC + Angola, Burundi, RDC, Rwanda, Sao Tomé-et-Principe) vise l'intégration politique/sécurité. Un justiciable camerounais victime de discrimination fondée sur sa nationalité (ex. refus de droit d'établissement au Gabon) saisira la Cour CEMAC. Une victime de torture ou d'arrestation arbitraire saisira la Cour CEEAC (si recours internes épuisés). La jurisprudence tend à coordonner : principe \emph{non-bis-in-idem} (pas deux cours pour même fait).
\end{savaistu}

\courssub{Niveau national : L'institution phare — La CNDHL}

C'est le maillon central, le plus accessible, le premier recours effectif pour tout résident du Cameroun.

\begin{definition}[Institution Nationale des Droits de l'Homme (INDH) — Principes de Paris 1993]
Organe indépendant créé par la loi (ou la Constitution) pour \textbf{promouvoir} et \textbf{protéger} les droits de l'homme au niveau national. Critères d'indépendance (statut « A ») : mandat large, pluralisme composition, indépendance fonctionnelle/financière, pouvoir d'auto-saisine, accès aux lieux de détention, rapport annuel public. \textbf{La CNDHL est l'INDH du Cameroun} (accréditée statut A depuis 2010, renouvelé 2015, 2020).
\end{definition}

\begin{propriete}[Indépendance fonctionnelle de la CNDHL (Loi 2004/016 modifiée par Loi 2019/014)]
\begin{itemize}
  \item \textbf{Nomination :} Président et 18 Commissaires nommés par décret du Président de la République pour \textbf{5 ans non renouvelables} (garantie contre pressions pour renouvellement).
  \item \textbf{Immunité d'opinion :} Aucune poursuite pour avis émis dans l'exercice des fonctions.
  \item \textbf{Budget autonome :} Ligne budgétaire distincte (théorique — en pratique, dépendance au Minfi reste un défi).
  \item \textbf{Pluralisme :} Représentation : magistrats, avocats, enseignants-chercheurs, ONG DH, syndicats, confessions religieuses, personnes handicapées, jeunesse, femmes.
  \item \textbf{Auto-saisine :} Peut agir sans plainte (ex. émeutes, conditions de détention).
\end{itemize}
\end{propriete}

\courssub{Missions de la CNDHL (Art. 2 Loi 2004/016 modifiée)}

\begin{itemize}
  \item \textbf{Promotion :} Éducation, formation, sensibilisation (écoles, prisons, médias, langues nationales), publication de guides, organisation journées thématiques.
  \item \textbf{Protection :} Réception plaintes, enquête, audition, visite lieux de détention (police, gendarmerie, prisons, centres rétention), médiation, \emph{recommandations} aux autorités, transmission au Procureur si infraction pénale.
  \item \textbf{Conseil :} Avis sur projets de lois/règlements (conformité DH), avis au Gouvernement/Parlement.
  \item \textbf{Rapport annuel :} Publié, remis au Président, au Parlement, au public. Baromètre de l'état des DH.
\end{itemize}

\begin{methode}[Comment saisir la CNDHL ? (Procédure gratuite, sans avocat obligatoire)]
\begin{enumerate}
  \item \textbf{Rédiger un courrier simple} (français, anglais ou langue locale) : identité, faits, date, lieu, auteurs présumés, preuves (photos, certificats médicaux, témoignages), droit violé invoqué.
  \item \textbf{Déposer} : Siège à Yaoundé (quartier Bastos) \textbf{OU} dans l'une des \textbf{10 antennes régionales} (voir carte) \textbf{OU} par courrier recommandé / email (cndhl@cndhl.cm).
  \item \textbf{Accusé de réception} : Numéro de dossier attribué, délai de traitement indicatif communiqué.
  \item \textbf{Instruction} : Enquêteur désigné $\to$ audition parties, visite lieux, demande documents aux administrations (obligation de collaboration).
  \item \textbf{Issue} : \textbf{Médiation} (accord écrit, suivi) \textbf{OU} \textbf{Recommandation} (adressée à l'autorité concernée, suivi à 6 mois) \textbf{OU} \textbf{Transmission au Procureur} (si faits constitutifs d'infraction pénale).
\end{enumerate}
\end{methode}

\begin{exemplebox}[Cas concret : Licenciement abusif d'une déléguée du personnel]
Madame M., déléguée du personnel dans une usine agroalimentaire à Bassa (Douala), est licenciée sans motivation écrite ni procédure disciplinaire, 3 jours après avoir déposé une revendication sur les EPI (équipements de protection individuelle). Elle saisit l'antenne CNDHL Littoral.
\begin{itemize}
  \item \textbf{Enquête :} Audition DRH, représentants syndicaux, consultation registre du personnel, convention collective.
  \item \textbf{Médiation :} 3 séances $\to$ accord : réintégration + paiement salaires échus + engagement écrit de respecter procédure disciplinaire.
  \item \textbf{Suivi :} CNDHL vérifie exécution à 3 et 6 mois. Cas clos sans contentieux judiciaire.
\end{itemize}
\end{exemplebox}

\begin{exoresolu}[Qualifier la compétence : CNDHL ou Tribunal ?]
\textbf{Énoncé :} M. K., agent contractuel d'une mairie, n'a pas perçu 8 mois de salaire. Il a adressé une mise en demeure au Maire (restée sans suite). Il hésite : saisir le Tribunal Administratif (contentieux) ou la CNDHL (médiation/recommandation) ?\par
\tcblower
\textbf{Solution détaillée :}
\begin{enumerate}
  \item \textbf{Analyse du fait :} Non-paiement salaire = violation droit au travail (Art. 6 PIDESC, Art. 15 Charte Africaine), droit à une rémunération équitable. Responsabilité : administration publique (mairie).
  \item \textbf{Voie CNDHL :} Gratuite, rapide (médiation 1-3 mois), recommandation non contraignante mais pression morale/politique forte (rapport annuel public). Adaptée si relation de travail pas encore rompue, volonté de préserver l'emploi.
  \item \textbf{Voie Tribunal Administratif :} Contentieux, contraignant (astreinte, condamnation pécuniaire), mais plus long (12-24 mois), frais d'huissier/avocat, risque de rupture définitive.
  \item \textbf{Stratégie optimale (complémentarité) :} Saisir \textbf{d'abord la CNDHL} (délai court, préservation relation). Si médiation échoue ou recommandation non suivie $\to$ saisir \textbf{Tribunal Administratif} (recours interne épuisé = condition pour saisine ultérieure CADHP/Cour CEEAC). Les deux ne sont pas exclusifs : la recommandation CNDHL peut servir de pièce au dossier judiciaire.
\end{enumerate}
\textbf{Conclusion :} La CNDHL est le \textbf{premier recours effectif, gratuit, non contentieux}. Le Tribunal est le \textbf{recours contentieux contraignant}. L'intelligence du justiciable est d'articuler les deux.
\end{exoresolu}

\begin{exemplebox}[Rapport d'activités CNDHL 2023 (données publiques synthétisées)]
\begin{center}
\begin{tabular}{|l|c|c|}
\hline
\rowcolor{popblueL}
\textbf{Indicateur} & \textbf{Chiffre} & \textbf{Commentaire} \\
\hline
Dossiers reçus & 1 245 & +12\% vs 2022 \\
\hline
Traités par médiation & 42\% (523) & Principal mode de règlement \\
\hline
Transmis au Parquet & 18\% (224) & Faits relevant infraction pénale (torture, viol, détention arbitraire) \\
\hline
Recommandations générales & 5 & Adressées au Gouvernement/Parlement (ex. surpeuplement prisons, droit nationalité, protection défenseurs DH) \\
\hline
Visites lieux de détention & 87 & Commissariats, brigades, prisons principales \\
\hline
Sensibilisation (personnes touchées) & $\approx$ 45 000 & Écoles, marchés, médias, langues nationales \\
\hline
\end{tabular}
\end{center}
\par\small\textit{Source : Rapport annuel CNDHL 2023, synthèse publique.}
\end{exemplebox}

\courssub{Autres acteurs nationaux du dispositif de protection}

La CNDHL n'est pas seule. L'architecture nationale repose sur une \textbf{complémentarité} :

\begin{itemize}
  \item \textbf{Le Médiateur de la République} (Loi 2019/010) : \emph{Non contentieux}. Règle les conflits \textbf{administration / administré} (dysfonctionnements, lenteurs, abus de pouvoir). Ne traite pas les violations DH \emph{per se} (torture, discrimination) mais les \emph{maladministrations}. Saisine gratuite, simple courrier. Recommandations, médiation. \textbf{Ne juge pas}.
  \item \textbf{Juridictions administratives} :
    \begin{itemize}
      \item \textbf{Tribunaux Administratifs} (10, chefs-lieux régions) : premier ressort contentieux administratif (excès de pouvoir, responsabilité administration, contrats publics).
      \item \textbf{Cours d'Appel — Chambres Administratives} (10, mêmes sièges) : appel des TA, certains contentieux directs (élections locales, certains fonctionnaires).
      \item \textbf{Cour Suprême — Section Administrative} : cassation, quelques contentieux directs (élections législatives/présidentielle, régularité actes Gouvernement).
    \end{itemize}
  \item \textbf{Police Judiciaire (PJ) / Gendarmerie} : Enquête infractions pénales (torture, violences, discriminations réprimées par Code Pénal). \textbf{Acteur clé} : la CNDHL transmet au Procureur, la PJ enquête sous autorité du Parquet.
  \item \textbf{ONG agréées (société civile organisée)} : REDHAC (Réseau Défenseurs DH Afrique Centrale), ACDH (Action Chrétienne DH), CONGEH (Coordination ONG Genre/DH), etc. Rôles : \emph{plaidoyer, assistance juridique, rapports alternatifs (shadow reports) aux Comités ONU/CADHP, accompagnement victimes, alertes précoces}.
\end{itemize}

\courssub{Articulation et complémentarité : La subsidiarité en action}

Le principe directeur est la \cle{subsidiarité} : \textbf{les recours internes (nationaux) doivent être épuisés avant de saisir les instances régionales/internationales}. Mais « épuisement » ne signifie pas « attendre indéfiniment » : recours \emph{effectifs}, \emph{accessibles}, \emph{dans un délai raisonnable}.

\begin{propriete}[Règles d'articulation (jurisprudence CADHP, Comités ONU, Cour CEEAC)]
\begin{itemize}
  \item \textbf{Épuisement des recours internes} : Condition de recevabilité (sauf si recours inexistants, ineffectifs, indûment prolongés, ou risque de dommages irreparables).
  \item \textbf{Non-bis-in-idem} : On ne peut pas saisir deux instances internationales pour le \emph{même fait} (ex. CADHP puis Comité CCPR). Mais on peut saisir séquentiellement : national $\to$ régional $\to$ universel.
  \item \textbf{Rôle des rapports alternatifs (Shadow Reports)} : Les ONG/INDH soumettent leurs propres rapports aux Comités ONU et à la CADHP \emph{en parallèle} du rapport d'État. Cela permet de \emph{contre-vérifier}, d'alerter sur angles morts, d'obtenir des recommandations plus ciblées. La CNDHL (statut A) a droit de parole propre devant la CADHP et les Comités ONU.
\end{itemize}
\end{propriete}

\begin{savaistu}[Le Cameroun à l'EPU 2023 : 3ème cycle, novembre 2023]
\begin{itemize}
  \item \textbf{210 recommandations reçues} (record).
  \item \textbf{185 acceptées} (88\%) : abolition peine de mort (moratoire maintenu), liberté de la presse (loi 2023 sur communication), droits des personnes handicapées, lutte contre traite personnes.
  \item \textbf{25 notées (refusées ou reportées)} : dépénalisation homosexualité (Art. 347-1 Code Pénal), légalisation avortement élargi, ratification Protocole Facultatif PIDESC, statut défenseurs DH.
  \item \textbf{Suivi :} Mise en place Comité national de suivi (Premier Ministre), rapport mi-parcours 2026. La CNDHL et les ONG surveillent l'exécution.
\end{itemize}
\end{savaistu}

\courssub{Illustration : Architecture multiniveau de protection — Organigramme de saisine et de rapportage}

\begin{popfigure}
\begin{tikzpicture}[
  node distance=1.2cm and 1.5cm,
  level/.style={sibling distance=3.5cm},
  inst/.style={rectangle, rounded corners, draw=popblue, thick, fill=popblueL, text width=3.2cm, align=center, minimum height=1.8cm, font=\small},
  instnat/.style={rectangle, rounded corners, draw=popgreen, thick, fill=popgreenL, text width=3.2cm, align=center, minimum height=1.8cm, font=\small},
  instsub/.style={rectangle, rounded corners, draw=poporange, thick, fill=poporangeL, text width=3.2cm, align=center, minimum height=1.8cm, font=\small},
  instreg/.style={rectangle, rounded corners, draw=poppurple, thick, fill=poppurpleL, text width=3.2cm, align=center, minimum height=1.8cm, font=\small},
  arrow/.style={->, >=Stealth, thick, draw=popdark},
  dasharrow/.style={->, >=Stealth, thick, draw=popmuted, dashed},
]
% Niveau Universel
\node[inst, align=center] (onu) at (0,0){\textbf{NIVEAU UNIVERSEL (ONU)}\\\vspace{0.2cm} Conseil DH (EPU, Proc. Spéciales)\\\vspace{0.1cm} Comités Traités (CCPR, CESCR,\\ CEDAW, CAT — Comm. indiv.)};
% Niveau Régional
\node[instreg, align=center] (cadhp) at (-5,-3.5){\textbf{NIVEAU RÉGIONAL (UA)}\\\vspace{0.2cm} Commission Africaine DH\\(Communications, Rapports,\\ Missions, Rapporteurs)};
\node[instreg, align=center] (couraf) at (5,-3.5){\textbf{Cour Africaine DH}\\\textbf{(Arusha)}\\\vspace{0.2cm} Contentieux (arrêts contraignants)\\ Avis consultatifs\\ \textit{Accès direct ind./ONG : non (Cmr)}};
% Niveau Sous-régional
\node[instsub, align=center] (cemac) at (-7,-6.5){\textbf{SOUS-RÉGIONAL}\\\vspace{0.2cm} Cour Justice CEMAC\\(N'Djaména)\\ Droits communautaires\\ Saisine ind./ONG : OUI};
\node[instsub, align=center] (ceeac) at (7,-6.5){\textbf{Cour Justice CEEAC}\\\textbf{(N'Djaména)}\\\vspace{0.2cm} Droits humains généraux\\ Saisine ind./ONG : OUI\\ (après épuisement intern.)};
% Niveau National
\node[instnat, text width=4cm] (cndhl) at (-5,-10) {\textbf{NATIONAL — INSTITUTION PHARE}\\\vspace{0.2cm} \textbf{CNDHL} (Statut A)\\\vspace{0.1cm} Promotion, Protection, Conseil,\\ Enquête, Médiation, Rapport annuel\\ 10 Antennes régionales};
\node[instnat, align=center] (mediateur) at (0,-10){\textbf{Médiateur de la République}\\\vspace{0.2cm} Maladministration\\ Administration/Administré\\ Médiation, Recommandations};
\node[instnat, align=center] (justice) at (5,-10){\textbf{Juridictions Admin.}\\\vspace{0.2cm} TA (1er ressort)\\ CA-Chambres Adm. (Appel)\\ CS-Sect. Adm. (Cassation)\\ \textit{Contentieux contraignant}};
\node[instnat, align=center] (pj) at (9,-10){\textbf{Police Judiciaire /}\\\textbf{ Gendarmerie}\\\vspace{0.2cm} Enquête infractions pénales\\ Sous autorité Parquet};
\node[instnat, align=center] (ong) at (-9,-10){\textbf{ONG Agréées}\\\vspace{0.2cm} (REDHAC, ACDH, CONGEH...)\\ Plaidoyer, Assistance jurid.,\\ Rapports alternatifs, Alerte};
% Flèches de saisine (montantes)
\draw[arrow] (cndhl) -- node[left, font=\tiny, text width=2.5cm, align=center] {Saisine gratuite, sans avocat} (cadhp);
\draw[arrow] (justice) -- node[right, font=\tiny, text width=2.5cm, align=center] {Recours contentieux exhaustion} (cadhp);
\draw[arrow] (cndhl) -- node[left, font=\tiny, text width=2.5cm, align=center] {Signalement / Plaidoyer pour déféré} (couraf);
\draw[arrow] (justice) -- node[right, font=\tiny, text width=2.5cm, align=center] {Épuisement\\ recours} (ceeac);
\draw[arrow] (cndhl) -- node[left, font=\tiny, text width=2.5cm, align=center] {Droits communautaires} (cemac);
\draw[arrow] (cadhp) -- node[left, font=\tiny, text width=2.5cm, align=center] {Défère affaire\\ Art.5 Protocole} (couraf);
\draw[arrow] (cadhp) -- node[above, font=\tiny, align=center]{Rapportage,\\Shadow reports} (onu);
\draw[arrow] (cndhl) -- node[above, font=\tiny, align=center]{Rapport annuel,\\Contribution EPU} (onu);
\draw[arrow] (ong) -- node[above, font=\tiny, align=center]{Shadow reports,\\Plaidoyer} (cadhp);
\draw[arrow] (ong) -- node[above, font=\tiny, align=center]{Shadow reports,\\Communications} (onu);
% Flèches de coordination nationales
\draw[dasharrow] (cndhl) -- node[above, font=\tiny] {Transmission Parquet} (pj);
\draw[dasharrow] (cndhl) -- node[above, font=\tiny, align=center]{Recommandations,\\Échange infos} (mediateur);
\draw[dasharrow] (cndhl) -- node[above, font=\tiny, align=center]{Recommandations,\\Pièces dossier} (justice);
\draw[dasharrow] (mediateur) -- node[above, font=\tiny] {Renvoi si DH} (cndhl);
\draw[dasharrow] (justice) -- node[above, font=\tiny, align=center]{Demande avis,\\Exécution} (cndhl);
\draw[dasharrow] (pj) -- node[above, font=\tiny, align=center]{Information,\\Collaboration} (cndhl);
% Légende
\node[draw=popdark, fill=popcream, rounded corners, font=\footnotesize, align=left, anchor=south east] at (11,-11) {
\textbf{Légende :}\\
\textcolor{popblue}{■} Universel (ONU)\\
\textcolor{poppurple}{■} Régional (UA)\\
\textcolor{poporange}{■} Sous-régional (CEMAC/CEEAC)\\
\textcolor{popgreen}{■} National\\
\textcolor{popdark}{$\rightarrow$} Saisine / Rapportage principal\\
\textcolor{popmuted}{$\dashrightarrow$} Coordination / Complémentarité nationale
};
\end{tikzpicture}
\legende{Architecture multiniveau de protection des droits de l'homme applicable au Cameroun : flux de saisine (flèches pleines) et coordination nationale (flèches pointillées).}
\end{popfigure}

\courssub{Illustration : Carte du Cameroun — Antennes CNDHL et Cours d'Appel (Chambres Administratives)}

\begin{popfigure}
\begin{tikzpicture}[
  scale=0.75,
  region/.style={draw=popblue, thick, fill=popblue!10},
  antenna/.style={circle, fill=popgreen, draw=popgreen!80!black, minimum size=0.35cm},
  cour/.style={rectangle, fill=poporange, draw=poporange!80!black, minimum size=0.35cm},
  label/.style={font=\tiny, text width=1.8cm, align=center},
  arrowN/.style={->, >=Stealth, thick, draw=popdark},
]
% Contour très simplifié du Cameroun (polygone approximatif)
\draw[region] (0,0) -- (1.5,0.5) -- (3,0) -- (4.5,0.5) -- (5.5,1.5) -- (5,3) -- (4,4) -- (3.5,5.5) -- (2.5,6.5) -- (1,6) -- (0,4.5) -- (-0.5,3) -- (0,1.5) -- cycle;
% Villes / Sièges (coordonnées approximatives relatives)
% Extrême-Nord : Maroua
\coordinate (MAR) at (0.8, 5.8);
% Nord : Garoua
\coordinate (GAR) at (1.8, 4.5);
% Adamaoua : Ngaoundéré
\coordinate (NGA) at (3.2, 3.8);
% Est : Bertoua
\coordinate (BER) at (4.8, 2.8);
% Centre : Yaoundé
\coordinate (YAO) at (2.8, 1.8);
% Sud : Ebolowa
\coordinate (EBO) at (2.0, 0.8);
% Littoral : Douala
\coordinate (DOU) at (1.2, 1.2);
% Ouest : Bafoussam
\coordinate (BAF) at (1.5, 2.8);
% Nord-Ouest : Bamenda
\coordinate (BAM) at (0.5, 3.5);
% Sud-Ouest : Buea
\coordinate (BUE) at (0.8, 2.0);
% Antennes CNDHL (cercles verts) + Cours d'Appel (carrés orange) — mêmes villes
\foreach \city/\name in {MAR/Maroua, GAR/Garoua, NGA/Ngaoundéré, BER/Bertoua, YAO/Yaoundé, EBO/Ebolowa, DOU/Douala, BAF/Bafoussam, BAM/Bamenda, BUE/Buea} {
   \node[antenna] at (\city) {};
   \node[cour] at ($(\city)+(0.35,0)$) {};
   \node[label, anchor=north] at ($(\city)+(0.175,-0.5)$) {\name};
}
% Légende
\node[draw=popdark, fill=popcream, rounded corners, font=\footnotesize, align=left, anchor=north west] at (-1.5, 6.5) {
\textbf{Cameroun — 10 Régions}\\
\textcolor{popgreen}{●} Antenne régionale CNDHL\\
\textcolor{poporange}{■} Cour d'Appel (Chambre Administrative)\\
\textit{Sièges identiques (chefs-lieux régionaux)}\\
\medskip
\textbf{Accès :} Tout justiciable peut se rendre\\
à l'antenne CNDHL ou au TA/CA de sa région.
};
% Flèche Nord
\draw[arrowN] (-1, 6.2) -- (-1, 5.7) node[midway, left, font=\tiny] {N};
\end{tikzpicture}
\legende{Localisation des 10 antennes régionales de la CNDHL (cercles verts) et des 10 Cours d'Appel siégeant en Chambres Administratives (carrés orange). Les deux réseaux maillent le territoire aux chefs-lieux régionaux, garantissant une accessibilité territoriale de la protection administrative et quasi-juridictionnelle.}
\end{popfigure}

\courssub{Synthèse opérationnelle : Choisir son niveau de saisine}

\begin{center}
\begin{tabularx}{\linewidth}{|l|X|X|X|}
\hline
\rowcolor{popblueL}
\textbf{Situation type} & \textbf{1er recours (National)} & \textbf{Recours régional (si échec/épuisement)} & \textbf{Recours universel (subsidiaire)} \\
\hline
Violence policière, torture & CNDHL (enquête, transmission Parquet) + PJ/Procureur & CADHP (Communication) + Cour CEEAC & CAT (Com. indiv.), CCPR (Com. indiv.) \\
\hline
Licenciement abusif / Droit du travail & CNDHL (médiation) $\to$ TA/CA & Cour CEEAC (si droit communautaire) & CESCR (pas com. indiv. Cmr), CCPR \\
\hline
Discrimination femme / Genre & CNDHL + TA & CADHP (Com.) + Cour CEEAC & CEDAW (Com. indiv. — Prot. ratifié) \\
\hline
Détention arbitraire / Justice & CNDHL (visite, réclam.) + TA (référé liberté) & CADHP (Com. urgente) + Cour CEEAC & CCPR (Com. indiv.), CAT \\
\hline
Droit nationalité / Apatridie & CNDHL (conseil, médiation) + TA & CADHP (Com.) + Cour CEMAC (droit communaut.) & CCPR, CERD \\
\hline
\end{tabularx}
\end{center}

\begin{attention}[Erreurs à ne pas commettre à l'examen]
\begin{itemize}
  \item \textbf{Confondre CNDHL et Tribunal :} La CNDHL \textbf{ne juge pas}, n'émet pas d'ordonnances, ne condamne pas. Elle recommande, médiatise, transmet au Parquet.
  \item \textbf{Oublier la condition d'épuisement des recours internes :} Saisir la CADHP ou la Cour CEEAC sans avoir saisi CNDHL/TA/CA $\to$ \textbf{irrecevabilité} (sauf exceptions : recours ineffectifs, délais déraisonnables, risque vie).
  \item \textbf{Croire que la Cour Africaine est accessible directement :} Le Cameroun n'a \textbf{pas} fait la déclaration Art. 34(6). Seule la Commission peut déférer l'affaire à la Cour.
  \item \textbf{Négliger les rapports alternatifs (Shadow Reports) :} C'est le levier principal de la société civile pour influencer les recommandations ONU/UA sans être partie à un contentieux.
  \item \textbf{Mélanger Médiateur et CNDHL :} Médiateur = maladministration (dysfonctionnement service public). CNDHL = violation droits humains (atteinte droits fondamentaux). Chevauchement possible mais mandats distincts.
\end{itemize}
\end{attention}

\begin{exercice}[Mise en situation : Qualifier, Saisiner, Argumenter]
\textbf{Cas :} Dans une entreprise de BTP à Kribi, 15 ouvriers (dont 3 femmes) n'ont pas reçu leur salaire depuis 4 mois. Le chef de chantier répond aux réclamations par des menaces de licenciement. Une ouvrière, Mme E., est renvoyée après avoir osé réclamer. Le délégué syndical contacte une ONG locale (membre de la REDHAC).
\begin{enumerate}
  \item Identifiez les droits violés (référez-vous aux textes : Constitution, Code Travail, PIDESC, Charte Africaine).
  \item Proposez une \textbf{stratégie de saisine multiniveau} (ordre chronologique, institutions concernées, arguments clés, pièces à joindre).
  \item Expliquez pourquoi la saisine de la Cour Africaine \emph{directement} par Mme E. est impossible aujourd'hui, et par quel biais elle pourrait dennoch y aboutir.
  \item Quel rôle l'ONG peut-elle jouer aux niveaux national, régional, universel (rapports alternatifs, plaidoyer, assistance) ?
\end{enumerate}
\end{exercice}

\begin{corrige}[Exercice — Mise en situation Kribi]
\begin{enumerate}
  \item \textbf{Droits violés :} Droit au travail (Art. 6 PIDESC, Art. 15 Charte), droit à une rémunération équitable (Art. 7 PIDESC), interdiction travail forcé (Art. 8 PIDCP, Art. 5 Charte), protection contre licenciement abusif/représailles syndicales (Code Travail Art. 35, 36, Convention OIT 98), non-discrimination femme (Art. 3 Charte, CEDAW), droit à un recours effectif (Art. 7 Charte, Art. 2 PIDCP).
  \item \textbf{Stratégie chronologique :}
    \begin{enumerate}
      \item \textbf{CNDHL Antenne Sud (Ebolowa) ou Littoral (Douala) :} Courrier collectif (15 ouvriers) + individuel (Mme E.). Pièces : contrats, fiches de paie antérieures, mises en demeure, témoignages, certificats médicaux si stress. Demande : médiation urgente, paiement arriérés, réintégration Mme E., sanction employeur.
      \item \textbf{Inspection du Travail (Ministère Travail) :} Saisine parallèle pour constat d'infraction (Code Travail), procès-verbal $\to$ transmission Procureur.
      \item \textbf{Tribunal Administratif (si employeur = personne publique) OU Tribunal de Grande Instance (privé) :} Recours contentieux paiement salaires + dommages-intérêts + réintégration (si privé : TGI ; si public : TA). \emph{Épuisement recours interne}.
      \item \textbf{CADHP (Communication) :} Si TA/TGI lent/ineffectif $\to$ Communication CADHP (Art. 55 Charte) au nom des 15 + Mme E. Arguments : violation Art. 7, 15, 26 Charte. Demande : mesures provisoires (Art. 98 Règlement) pour paiement urgent.
      \item \textbf{Cour CEEAC :} Après épuisement (jugement définitif non exécuté ou déni de justice) $\to$ Saisine Cour CEEAC (Protocole 2005). Arrêt contraignant.
      \item \textbf{Comité CCPR (Communication) :} Parallèle possible (PIDCP ratifié + Protocole). Violation Art. 8 (travail forcé), 26 (non-discrimination), 2 (recours effectif).
    \end{enumerate}
  \item \textbf{Cour Africaine inaccessible directement :} Cameroun n'a pas déposé déclaration Art. 34(6) Protocole Cour. \textbf{Biais possible :} La CADHP, saisie en communication, peut \emph{déférer l'affaire à la Cour} (Art. 5(1) Protocole Cour) si elle estime que l'affaire révèle des violations graves/massives ou que l'État ne se conforme pas à ses recommandations. L'ONG peut plaider ce déféré.
  \item \textbf{Rôle ONG (REDHAC locale) :}
    \begin{itemize}
      \item \textbf{National :} Assistance rédaction plaintes CNDHL/Inspection/TGI, accompagnement victimes, médiation paritaire, alerte médias, plaidoyer auprès MinTravail.
      \item \textbf{Régional :} Rédaction \emph{Shadow Report} pour CADHP (session ordinaire), participation Forum ONG pré-session, communication conjointe CADHP, demande mission pays Rapporteur Spécial Droits Économiques/Sociaux.
      \item \textbf{Universel :} Contribution rapport alternatif EPU 2027 (3ème cycle suivi), communication Comités CCPR/CEDAW/CAT, plaidoyer auprès Rapporteurs Spéciaux ONU (travail forcé, droits migrants si ouvriers migrants).
    \end{itemize}
\end{enumerate}
\end{corrige}

\coursec{TD1 : Étude de terrain d'une structure locale de promotion des droits de l'homme (CNDHL ou ONG)}

Après avoir présenté l'architecture multiniveau des institutions de protection (ONU, Union africaine, niveau national), nous passons à l'investigation de terrain. Ce travail dirigé (TD1) permet de confronter les normes juridiques à la réalité d'une structure locale — la Commission nationale des droits de l'homme et des libertés (CNDHL) ou une organisation non gouvernementale (ONG) — et de développer les savoir-faire « qualifier, saisir, argumenter » dans un cadre rigoureux.

\begin{aretenir}[Objectifs du TD1]
\begin{itemize}
    \item \textbf{Savoir} : vérifier la connaissance des institutions de protection des droits de l'homme par l'observation directe.
    \item \textbf{Savoir-faire} : appliquer les notions de l'unité à une situation concrète ; rédiger et justifier une démarche, une réponse ou une conclusion.
    \item \textbf{Savoir-être} : adopter une posture d'observateur neutre, respectueux des personnes et de l'éthique de la recherche.
\end{itemize}
\end{aretenir}

\courssub{Préparation collective de la grille d'enquête}

La première étape consiste à élaborer, en groupe classe, une grille d'enquête structurée. Chaque binôme propose des indicateurs ; le professeur synthétise et valide la version finale. La grille couvre neuf dimensions essentielles :

\begin{itemize}
    \item \textbf{Identité} : dénomination, statut juridique, date de création, rattachement administratif.
    \item \textbf{Cadre légal} : textes fondateurs (loi n°2004/016 du 22 juillet 2004 pour la CNDHL, statuts et récépissé pour une ONG), décrets d'application, agréments.
    \item \textbf{Missions} : promotion, protection, éducation aux droits humains, veille, plaidoyer, suivi des recommandations internationales.
    \item \textbf{Organisation} : organigramme, composition du bureau, nombre de commissaires ou de permanents, antennes régionales.
    \item \textbf{Moyens humains et financiers} : effectifs, qualifications, budget annuel (fonctionnement, investissement, activités), sources de financement.
    \item \textbf{Activités réelles} : nombre de plaintes reçues et traitées, enquêtes menées, visites des lieux de détention, formations, campagnes de sensibilisation.
    \item \textbf{Partenariats} : institutions étatiques (MINJUSTICE, MINAT, forces de sécurité, juridictions), organisations internationales (ONU, Union africaine, ONG internationales), société civile locale.
    \item \textbf{Difficultés} : sous-financement, manque de personnel, pressions diverses, méconnaissance du public, locaux inadaptés.
    \item \textbf{Perception du public} : accessibilité (horaires, localisation, langue), confiance, notoriété, retours des usagers.
\end{itemize}

\begin{popfigure}
\begin{tikzpicture}[scale=0.85, every node/.style={font=\small}]
% En-têtes
\draw[popblue, line width=1.2pt] (0,0) rectangle (16,-0.7);
\node[anchor=west] at (0.2,-0.35) {\textbf{Indicateur}};
\node[anchor=west] at (3.2,-0.35) {\textbf{Question guide}};
\node[anchor=west] at (7.2,-0.35) {\textbf{Réponse observée}};
\node[anchor=west] at (11.2,-0.35) {\textbf{Source}};
\node[anchor=west] at (14.2,-0.35) {\textbf{Appréciation}};
% Lignes d'exemple
\foreach \i/\ind/\q/\r/\s/\a in {
1/Identité/Quel est le statut juridique de la structure ?/CNDHL -- loi 2004-016/Statuts et décret/Fort,
2/Missions/Quelles missions sont effectivement menées ?/Promotion ; plaintes ; détention/Rapports annuels/Moyen,
3/Moyens/Quel est le budget annuel de fonctionnement ?/12 M FCFA (2023)/Comptabilité/Faible,
4/Activités/Combien de plaintes traitées en 2023 ?/85 plaintes/Registre plaintes/Moyen,
5/Partenariats/Quels partenaires techniques et financiers ?/PNUD ; MINJUSTICE ; ONG locales/Conventions/Fort}
{
\pgfmathsetmacro{\y}{-0.7 - \i*0.9};
\draw[popline] (0,\y) rectangle (16,\y-0.9);
\node[anchor=west] at (0.2,\y-0.45) {\ind};
\node[anchor=west, text width=3.8cm, align=left] at (3.2,\y-0.45) {\q};
\node[anchor=west, text width=3.8cm, align=left] at (7.2,\y-0.45) {\r};
\node[anchor=west, text width=2.8cm, align=left] at (11.2,\y-0.45) {\s};
\node[anchor=west] at (14.2,\y-0.45) {\a};
}
\end{tikzpicture}
\legende{Modèle de grille d'enquête (extrait) : cinq colonnes — Indicateur, Question guide, Réponse observée, Source, Appréciation (Fort / Moyen / Faible). Les données chiffrées sont fictives et servent d'exemple de remplissage.}
\end{popfigure}

\begin{methode}[Conduite d'un entretien semi-directif]
\textbf{1. Thèmes préparés} : lister à l'avance les neuf dimensions de la grille.\\
\textbf{2. Questions ouvertes} : « Pouvez-vous décrire comment s'organise le traitement d'une plainte ? » plutôt que « Traitez-vous les plaintes ? ».\\
\textbf{3. Relances} : approfondir par « Pourquoi ? », « Comment ? », « Quel exemple concret ? ».\\
\textbf{4. Prise de notes et enregistrement} : désigner un secrétaire de séance ; l'enregistrement audio n'est possible qu'avec l'accord écrit de la personne interrogée.\\
\textbf{5. Remerciements et compte rendu validé} : envoyer le résumé à l'interlocuteur pour validation avant intégration au rapport.
\end{methode}

\begin{attention}[Éthique de la recherche]
\begin{itemize}
    \item \textbf{Anonymat des victimes} : ne jamais citer de noms ni de détails permettant d'identifier une personne dans le rapport des élèves.
    \item \textbf{Neutralité de l'observateur} : ne pas prendre parti ; distinguer le fait (chiffre, document) de l'opinion (appréciation).
    \item \textbf{Restitution des résultats} : remettre une synthèse à la structure visitée (respect de la réciprocité).
\end{itemize}
\end{attention}

\courssub{Réalisation de la visite (ou visioconférence / exploitation d'un rapport d'activité)}

Deux modalités sont possibles selon le contexte :

\begin{itemize}
    \item \textbf{Visite physique} (à privilégier) : déplacement à l'antenne régionale de la CNDHL ou au siège de l'ONG. Observation des locaux (accessibilité, affichage des droits, salle d'accueil), entretien avec un commissaire ou un permanent, collecte de documents (rapports d'activité, affiches de sensibilisation, formulaires de plainte).
    \item \textbf{Visioconférence ou exploitation documentaire} : si le déplacement est impossible, organiser un échange en visioconférence avec le responsable et exploiter le dernier rapport d'activité publié (souvent disponible sur le site web de la CNDHL ou de l'ONG).
\end{itemize}

Durant la visite, les élèves remplissent la grille en temps réel, photographient (avec autorisation) les panneaux d'information et les registres anonymisés, et notent leurs impressions sur l'accueil et la réactivité du personnel.

\courssub{Analyse croisée : les textes face à la réalité de terrain}

L'analyse compare les dispositions de la loi n°2004/016 du 22 juillet 2004 (création, organisation et fonctionnement de la CNDHL) aux constats de terrain.

\begin{definition}[Effectivité des droits]
L'\cle{effectivité} des droits est la mesure dans laquelle les droits proclamés dans les textes sont réellement exercés par les bénéficiaires sur le terrain. Elle s'apprécie selon quatre critères : \textbf{disponibilité} (existence de structures et de services), \textbf{accessibilité} (géographique, économique, informationnelle), \textbf{acceptabilité} (respect des cultures et des langues) et \textbf{qualité} (compétence, diligence, suivi).
\end{definition}

\begin{exemplebox}[Antenne régionale de la CNDHL (données fictives d'entraînement, 2023)]
\begin{itemize}
    \item \textbf{Effectifs} : 3 enquêteurs pour 10 départements (soit environ 1 enquêteur pour 3,3 départements).
    \item \textbf{Budget de fonctionnement} : 12 millions de FCFA par an.
    \item \textbf{Plaintes traitées} : 85 dossiers (dont 60\,\% concernant des droits civils, 25\,\% des droits économiques et sociaux, 15\,\% des droits politiques).
    \item \textbf{Visites de lieux de détention} : 2 visites dans des maisons d'arrêt de la région.
    \item \textbf{Locaux} : bureaux au 2\textsuperscript{e} étage sans ascenseur, absence de signalétique en langues locales.
\end{itemize}
\end{exemplebox}

\begin{aretenir}[Grille d'appréciation croisée (exemple de remplissage)]
\begin{center}
\begin{tabularx}{\linewidth}{|l|X|X|X|}\hline
\textbf{Critère} & \textbf{Prescription légale} & \textbf{Constat de terrain (exemple)} & \textbf{Écart} \\ \hline
Composition & Commissaires et enquêteurs en nombre suffisant & 1 commissaire + 3 enquêteurs pour 10 départements & \textbf{Fort} (sous-effectif) \\ \hline
Budget & Moyens nécessaires inscrits au budget de l'État & 12 M FCFA de fonctionnement, 0 M d'investissement & \textbf{Fort} (sous-financement) \\ \hline
Accessibilité & Locaux accessibles à tous & Étage sans ascenseur, horaires 8 h -- 15 h & \textbf{Moyen} \\ \hline
Traitement des plaintes & Délai raisonnable fixé par le règlement intérieur (30 jours) & Délai moyen de 45 jours (manque de personnel) & \textbf{Moyen} \\ \hline
Suivi des recommandations & Rapports transmis aux pouvoirs publics & 2 recommandations suivies sur 5 & \textbf{Fort} \\ \hline
\end{tabularx}
\end{center}
\end{aretenir}

\courssub{Identification des forces et des faiblesses}

À partir de l'analyse croisée, les élèves dressent un diagnostic simplifié de type SWOT (Forces, Faiblesses, Opportunités, Menaces), centré ici sur les Forces et les Faiblesses.

\begin{exemplebox}[Synthèse Forces / Faiblesses (exemple d'antenne régionale)]
\begin{itemize}
    \item \textbf{Forces} : légitimité institutionnelle (loi n°2004/016), ancrage territorial (présence au chef-lieu de région), expertise juridique des commissaires, partenariats avec le MINJUSTICE et le PNUD.
    \item \textbf{Faiblesses} : sous-financement chronique (budget de fonctionnement seul), effectifs insuffisants (3 enquêteurs pour 10 départements), méconnaissance du public (peu de plaintes spontanées), absence de pouvoir contraignant (la CNDHL émet des avis et recommandations, sans pouvoir de sanction), locaux non accessibles aux personnes à mobilité réduite.
\end{itemize}
\end{exemplebox}

\begin{popfigure}
\begin{tikzpicture}
\begin{axis}[
    ybar,
    bar width=18pt,
    width=\linewidth,
    height=7cm,
    enlarge x limits=0.3,
    symbolic x coords={Fonctionnement,Investissement,Activités},
    xtick=data,
    xlabel={Poste budgétaire},
    ylabel={Montant (millions de FCFA)},
    ymin=0,
    ymax=14,
    nodes near coords,
    nodes near coords align={vertical},
    tick label style={font=\small},
    label style={font=\small}
]
\addplot[draw=popblue,fill=popblue!30] coordinates {(Fonctionnement,12) (Investissement,0) (Activités,2)};
\end{axis}
\end{tikzpicture}
\legende{Répartition budgétaire d'une antenne régionale de la CNDHL (données fictives d'entraînement, 2023) : fonctionnement 12 millions de FCFA, investissement 0 million, activités 2 millions.}
\end{popfigure}

\courssub{Rédaction du rapport de TD1}

Le rapport suit une structure imposée, garantissant la rigueur de la démarche et la lisibilité pour l'évaluation.

\begin{methode}[Structure type du rapport de TD1]
\textbf{1. Introduction} : objet du TD, structure visitée, dates, méthode (grille, entretien, documents).\\
\textbf{2. Présentation de la structure} : identité, cadre légal, missions, organigramme, moyens.\\
\textbf{3. Analyse du fonctionnement} : confrontation textes / terrain (tableau croisé), effectivité (disponibilité, accessibilité, acceptabilité, qualité).\\
\textbf{4. Évaluation de l'effectivité} : appréciation globale (Fort / Moyen / Faible) par critère, justification chiffrée.\\
\textbf{5. Propositions} : recommandations concrètes (ex. : recrutement de 2 enquêteurs, ligne d'investissement de 5 M FCFA, campagne radio en langues locales).\\
\textbf{6. Conclusion} : synthèse des forces et faiblesses, portée pour la protection des droits au niveau local.\\
\textbf{7. Annexes} : grille d'enquête complétée, photos autorisées (locaux, affichages), documents collectés (extrait du rapport d'activité, formulaire de plainte anonymisé), fiche de synthèse d'une plainte type.
\end{methode}

\begin{aretenir}[Boîte à outils -- Fiche de synthèse d'une plainte type (annexe, exemple fictif)]
\begin{center}
\begin{tabularx}{\linewidth}{|l|X|}\hline
\textbf{Champ} & \textbf{Contenu} \\ \hline
N° d'enregistrement & 2023-0042 \\ \hline
Date de réception & 12/03/2023 \\ \hline
Auteur (anonymisé) & M. K., 34 ans, commerçant \\ \hline
Droit violé & Droit au travail / non-paiement du salaire (art. 6 du Pacte international relatif aux droits économiques, sociaux et culturels ; art. 23 de la DUDH) \\ \hline
Acteur mis en cause & Entreprise X SARL (secteur du BTP) \\ \hline
Suite donnée & Médiation de la CNDHL $\rightarrow$ accord transactionnel (paiement de 3 mois d'arriérés) \\ \hline
Délai de traitement & 38 jours (objectif : 30 jours) \\ \hline
\end{tabularx}
\end{center}
\end{aretenir}

\courssub{Restitution orale et débat}

Chaque groupe présente une synthèse de 5 minutes (diaporama de 6 à 8 diapositives) devant la classe. La restitution met en avant :

\begin{itemize}
    \item les constats majeurs (effectivité, forces et faiblesses) ;
    \item deux propositions prioritaires argumentées ;
    \item une réflexion critique sur le rôle respectif de la société civile (ONG) et des institutions de l'État (CNDHL) dans la protection des droits humains.
\end{itemize}

Un débat, modéré par l'enseignant, suit les exposés. Points de discussion suggérés :

\begin{itemize}
    \item \textbf{Complémentarité ou substitution} : l'ONG peut-elle combler les lacunes de la CNDHL sans affaiblir sa légitimité ?
    \item \textbf{Indépendance} : comment garantir l'autonomie financière et politique des commissaires ?
    \item \textbf{Accessibilité numérique} : intérêt d'une plateforme de plainte en ligne (application mobile) pour réduire les délais de traitement.
\end{itemize}

\begin{savaistu}[La CNDHL et le système africain des droits de l'homme]
La CNDHL participe aux travaux de la Commission africaine des droits de l'homme et des peuples (CADHP), notamment lors de l'examen du rapport périodique du Cameroun. Comme d'autres institutions nationales des droits de l'homme, elle peut présenter des informations complémentaires (rapports dits « alternatifs »), ce qui renforce le dialogue entre l'État et la société civile au niveau continental.
\end{savaistu}

\begin{exoresolu}[Qualification d'une plainte reçue par l'antenne]
\textbf{Énoncé} : Un citoyen dépose une plainte pour « refus d'inscription sur les listes électorales » auprès de l'antenne régionale de la CNDHL. Le formulaire mentionne : nom, bureau de vote, date de dépôt, pièce d'identité. Procédez à la qualification de la plainte et décrivez la procédure de traitement.
\tcblower
\textbf{Solution détaillée} :
\begin{enumerate}
    \item \textbf{Qualification juridique} : violation du droit de participer à la direction des affaires publiques (art. 25 du Pacte international relatif aux droits civils et politiques ; art. 13 de la Charte africaine des droits de l'homme et des peuples ; préambule de la Constitution camerounaise du 18 janvier 1996).
    \item \textbf{Compétence} : la CNDHL peut être saisie de toute violation des droits de l'homme, y compris des droits civiques, conformément à ses missions de promotion et de protection (loi n°2004/016).
    \item \textbf{Procédure} : enregistrement de la plainte (N° 2023-0087), accusé de réception, transmission au commissaire régional pour enquête.
    \item \textbf{Délai} : le règlement intérieur fixe un objectif de 30 jours ; en pratique, le délai moyen constaté est de 38 jours (cf. fiche de synthèse).
    \item \textbf{Issue possible} : recommandation adressée à l'organe de gestion des élections (ELECAM) de régulariser l'inscription ; suivi par la CNDHL jusqu'à l'exécution.
\end{enumerate}
\end{exoresolu}

\begin{exercice}[Mise en situation -- Élaborer une recommandation]
À partir des faiblesses identifiées (sous-financement, effectifs insuffisants, accessibilité), rédigez une recommandation adressée au Président de la CNDHL, structurée en quatre points : (1) constat chiffré, (2) base légale (loi n°2004/016), (3) proposition concrète (budget, recrutement, locaux), (4) indicateurs de suivi (délai, taux de traitement). \emph{Travail en binôme, à rendre en annexe du rapport.}
\end{exercice}

\begin{corrige}[Exercice n°1 -- Mise en situation]
\textbf{1. Constat} : budget de fonctionnement de 12 M FCFA, 0 M d'investissement, 3 enquêteurs pour 10 départements, délai moyen de traitement de 45 jours.\\
\textbf{2. Base légale} : la loi n°2004/016 du 22 juillet 2004 confie à la CNDHL des missions de promotion et de protection des droits de l'homme et prévoit qu'elle dispose des moyens nécessaires à leur accomplissement, son budget étant inscrit au budget de l'État.\\
\textbf{3. Proposition} : (a) inscrire une ligne d'investissement de 50 M FCFA (locaux accessibles, véhicule) ; (b) recruter 7 enquêteurs supplémentaires (pour atteindre environ 1 enquêteur par département) ; (c) ouvrir des permanences décentralisées dans les chefs-lieux de département.\\
\textbf{4. Indicateurs de suivi} : taux d'exécution budgétaire supérieur à 90\,\% ; nombre d'enquêteurs par département au moins égal à 1 ; délai moyen de traitement des plaintes inférieur ou égal à 30 jours ; au moins 5 permanences opérationnelles d'ici la fin de l'année suivante.
\end{corrige}

\coursec{Mise en situation : Qualifier, Saisir, Argumenter (Intégration des savoir-faire)}

Après l'immersion de terrain réalisée lors du TD1 (section précédente), où vous avez observé le fonctionnement concret d'une structure comme la CNDHL ou une ONG partenaire, nous passons maintenant à l'appropriation pratique des outils juridiques. Savoir \emph{où} aller ne suffit plus : il faut maîtriser \emph{comment} qualifier une violation, \emph{comment} rédiger une saisine recevable et \emph{comment} construire un argumentaire solide devant un juge ou un médiateur. Cette section est le cœur de l'intégration des savoir-faire exigés au programme : \cle{qualifier}, \cle{saisir}, \cle{argumenter}.

\courssub{Exercice de qualification juridique : de la situation vécue à la norme violée}

La \cle{qualification juridique} est l'opération intellectuelle qui consiste à faire correspondre un fait brut (une expulsion, un licenciement, un refus d'inscription) à une catégorie juridique protégée (droit au logement, liberté syndicale, droit à l'éducation) et à identifier le texte précis qui le garantit. C'est la boussole qui oriente le choix de la voie de recours.

\begin{methode}[Démarche de qualification en 4 étapes]
\begin{enumerate}
    \item \textbf{Isoler les faits matériels} : Qui ? Quoi ? Où ? Quand ? Comment ? (Dates, lieux, auteurs présumés, témoins, preuves matérielles : photos, courriers, certificats médicaux).
    \item \textbf{Identifier le droit violé} : s'agit-il d'un droit civil et politique (liberté, vie privée, procès équitable) ou d'un droit économique, social et culturel (travail, éducation, logement, santé) ?
    \item \textbf{Rattacher au texte applicable} : Constitution du 18 janvier 1996 (Préambule), Pactes internationaux de 1966 (droits civils et politiques ; droits économiques, sociaux et culturels), Charte africaine des droits de l'homme et des peuples (1981), lois nationales (Code du travail, Code pénal, lois sur la liberté d'association).
    \item \textbf{Déterminer l'urgence et l'instance compétente} : danger imminent $\rightarrow$ juge des référés ; violation par une administration $\rightarrow$ CNDHL ou juridiction administrative ; infraction pénale $\rightarrow$ Procureur de la République ; voies internes épuisées $\rightarrow$ Commission africaine ou organes de traités de l'ONU.
\end{enumerate}
\end{methode}

\begin{exoresolu}[Cas pratique n°1 : expulsion forcée dans un quartier de Douala]
\textbf{Énoncé.} Le 15 mars 2024, à 5 heures du matin, une compagnie de sécurité privée, escortée par des éléments de police, détruit les habitations de 50 familles installées depuis 12 ans sur un terrain revendiqué par un promoteur immobilier. Aucune décision de justice, aucun préavis, aucune proposition de relogement. Trois personnes sont légèrement blessées.

\textbf{Travail à faire :} qualifiez juridiquement cette situation (droits violés, textes applicables, urgence, instance compétente).
\tcblower
\textbf{Solution commentée.}
\begin{itemize}
    \item \textbf{Droits violés :} droit au logement (art. 11 du Pacte international relatif aux droits économiques, sociaux et culturels), droit à la propriété (art. 17 de la Déclaration universelle des droits de l'homme), interdiction des traitements inhumains ou dégradants (art. 5 de la Charte africaine), droit à un procès équitable (art. 7 de la Charte africaine).
    \item \textbf{Textes applicables :} Préambule de la Constitution de 1996, Pactes de 1966 et Charte africaine (ratifiés par le Cameroun), Code pénal (destruction de biens, violences).
    \item \textbf{Urgence :} extrême (familles à la rue, risque de récidive, preuves qui disparaissent).
    \item \textbf{Instances compétentes :} le juge des référés pour faire cesser l'opération illégale ; le Procureur de la République pour les infractions pénales (destructions, violences) ; la CNDHL en parallèle pour enquête, médiation et accompagnement des victimes.
\end{itemize}
La qualification « expulsion forcée » est plus forte que « démolition » : l'absence de décision judiciaire rend l'opération illégale et engage la responsabilité de l'État (police présente). La saisine de la CNDHL complète l'action judiciaire par la médiation en vue du relogement.
\end{exoresolu}

\begin{exercice}[Cas pratiques à qualifier (travail en binôme)]
Analysez les deux cas ci-dessous en appliquant la méthode en 4 étapes.
\begin{enumerate}
    \item \textbf{Cas 2 (licenciement d'un syndicaliste) :} M. Kamga, délégué du personnel élu dans une entreprise agro-industrielle de Njombé, est licencié pour « faute grave » 48 heures après avoir déposé un préavis de grève pour non-paiement de primes. L'employeur n'a pas demandé l'autorisation préalable de l'Inspection du travail, pourtant obligatoire pour le licenciement d'un délégué du personnel.
    \item \textbf{Cas 3 (refus d'inscription d'un enfant handicapé) :} Mme Abena tente d'inscrire son fils de 10 ans, malvoyant, à l'école publique de son quartier. Le directeur refuse : « Nous n'avons pas de matériel adapté, allez à l'école spécialisée ». Aucune notification écrite n'est remise.
\end{enumerate}
\end{exercice}

\begin{corrige}[Exercice n°1]
\begin{itemize}
    \item \textbf{Cas 2 :} droit violé = liberté syndicale et protection du délégué du personnel (Convention n°87 de l'OIT, art. 10 de la Charte africaine, Code du travail camerounais qui subordonne le licenciement d'un délégué à l'autorisation de l'inspecteur du travail). Urgence = forte (précarité immédiate, effet dissuasif sur l'action syndicale). Instances = juridiction du travail, Inspection du travail (constat de l'infraction), CNDHL.
    \item \textbf{Cas 3 :} droit violé = droit à l'éducation sans discrimination (art. 13 du Pacte international relatif aux droits économiques, sociaux et culturels, art. 17 de la Charte africaine, loi n°2010/002 du 13 avril 2010 relative à la protection et au bien-être des personnes handicapées, qui impose l'accès à l'éducation inclusive). Urgence = moyenne (rentrée scolaire proche). Instances = CNDHL (médiation rapide), recours hiérarchique auprès du ministère de l'Éducation de base, juge administratif si le refus persiste.
\end{itemize}
\end{corrige}

\begin{popfigure}
\begin{tikzpicture}[
    node distance=1.1cm and 1.4cm,
    decision/.style={diamond, draw=popdark, thick, fill=poporangeL, text width=3.2cm, align=center, inner sep=3pt},
    process/.style={rectangle, draw=popblue, thick, fill=popblueL, text width=3.4cm, align=center, rounded corners, inner sep=4pt},
    startstop/.style={ellipse, draw=popgreen, thick, fill=popgreenL, text width=2.8cm, align=center, inner sep=4pt},
    arrow/.style={-Stealth, thick, draw=popdark},
    yesno/.style={font=\small\bfseries, pos=0.5, above}
]
\node (start) [startstop] {Violation d'un droit constatée};
\node (q1) [decision, below of=start] {Atteinte grave à une liberté fondamentale ?};
\node (q2) [decision, below of=q1, yshift=-0.8cm] {Urgence absolue ? (danger imminent)};
\node (q3) [decision, right of=q1, xshift=4.5cm] {Auteur : administration ou service public ?};
\node (q4) [decision, below of=q3, yshift=-0.8cm] {Voies internes épuisées ?};
\node (p1) [process, left of=q2, xshift=-4.2cm] {Référé devant le juge (suspension immédiate)};
\node (p2) [process, below of=q2, yshift=-0.8cm] {Saisine CNDHL (enquête, médiation)};
\node (p3) [process, right of=q4, xshift=4.2cm] {Plainte au Procureur (voie pénale)};
\node (p4) [process, below of=q4, yshift=-0.8cm] {Commission africaine / organes ONU (après épuisement)};
\draw [arrow] (start) -- (q1);
\draw [arrow] (q1) -- node[yesno] {OUI} (q2);
\draw [arrow] (q1) -- node[yesno] {NON} (q3);
\draw [arrow] (q2) -- node[yesno] {OUI} (p1);
\draw [arrow] (q2) -- node[yesno] {NON} (p2);
\draw [arrow] (q3) -- node[yesno] {NON} (p3);
\draw [arrow] (q3) -- node[yesno] {OUI} (q4);
\draw [arrow] (q4) -- node[yesno] {OUI} (p4);
\draw [arrow] (q4) -- node[yesno] {NON} (p2);
\end{tikzpicture}
\legende{Algorithme de décision : « Quelle voie de recours choisir ? » Ce schéma guide le justiciable vers la procédure la plus adaptée selon la nature du droit, l'urgence et l'auteur de la violation.}
\end{popfigure}

\begin{attention}[Piège fréquent : confondre qualification et narration]
Au Probatoire technique, l'élève qui raconte l'histoire (« M. Kamga a été renvoyé parce qu'il a fait la grève ») sans qualifier (« licenciement abusif d'un délégué du personnel en violation du Code du travail et de la liberté syndicale ») perd l'essentiel des points. La qualification est un \textbf{acte juridique}, pas un récit. Utilisez le vocabulaire exact : « discrimination », « détention arbitraire », « violation du droit à un procès équitable ».
\end{attention}

\courssub{Atelier rédaction : la saisine efficace (CNDHL ou juge)}

Une saisine mal rédigée est une saisine irrecevable ou sans effet. L'administration et les juges ne « devinent » pas : ils lisent. La rigueur formelle est une condition de l'effectivité du droit.

\begin{methode}[Rédaction juridique efficace : la règle F.D.D.]
\begin{enumerate}
    \item \textbf{Faits} : précis (date, heure, lieu), datés, sourcés (procès-verbaux, photos, accusés de réception, témoins). Style sobre, chronologique, impersonnel. Évitez : « Ils sont venus méchamment ». Préférez : « Une équipe d'agents a procédé à la démolition sans présenter de décision de justice. »
    \item \textbf{Droit} : article exact (numéro), convention ratifiée, loi nationale pertinente. Citez le texte, ne le paraphrasez pas.
    \item \textbf{Demandes} : distinctes (urgence / fond), chiffrées si réparation (devis, factures), concrètes si injonction (« reloger les familles sur un site équipé » et non « trouver une solution »).
\end{enumerate}
\end{methode}

\begin{exemplebox}[Structure type d'une saisine CNDHL]
\begin{enumerate}
    \item \textbf{En-tête} : nom, prénom, qualité de l'expéditeur (victime, ayant droit, ONG), adresse, téléphone ; destinataire (Monsieur le Président de la CNDHL, Yaoundé) ; date ; objet précis.
    \item \textbf{Exposé chronologique des faits} : dates, lieux, auteurs, démarches déjà effectuées (plainte, courriers restés sans réponse) ; liste des pièces jointes.
    \item \textbf{Qualification juridique} : droits violés et textes applicables (Constitution, Pactes de 1966, Charte africaine, lois nationales).
    \item \textbf{Demandes} : enquête, médiation, recommandations, mesures urgentes (relogement, réparation).
    \item \textbf{Formule de politesse, date et signature}, liste numérotée des pièces jointes.
\end{enumerate}
\end{exemplebox}

\courssub{Simulation de médiation : la voie non contentieuse à la CNDHL}

La médiation n'est pas une faiblesse, c'est une stratégie. Elle est plus rapide que le procès, gratuite, confidentielle et préserve le lien social (voisinage, relation employeur-employé).

\begin{experience}[Simulation de médiation CNDHL — protocole pédagogique]
\textbf{Objectif :} vivre le processus de médiation pour comprendre le rôle du médiateur (impartialité, écoute, recherche d'options) et la dynamique de l'accord.

\textbf{Rôles (groupes de 3) :}
\begin{itemize}
    \item \textbf{La victime} prépare son récit (faits, préjudice, attentes), ses « lignes rouges » (ce qui est inacceptable) et ses marges de manœuvre.
    \item \textbf{L'auteur présumé} (administration, entreprise, particulier) prépare sa version, ses contraintes (budget, hiérarchie, image) et ses risques en cas de procès.
    \item \textbf{Le médiateur CNDHL} prépare sa grille d'entretien : accueil, règles (confidentialité, bonne foi), techniques de reformulation, recherche d'une solution acceptable par les deux parties.
\end{itemize}

\textbf{Déroulement (30 min) :}
\begin{enumerate}
    \item Séance plénière (5 min) : le médiateur pose le cadre et vérifie la volonté de médier.
    \item Entretiens séparés (2 $\times$ 5 min) : le médiateur reçoit chaque partie seule pour explorer les besoins réels.
    \item Négociation conjointe (10 min) : échanges directs, propositions et contre-propositions.
    \item Clôture (5 min) : rédaction d'un \textbf{procès-verbal d'accord} (engagements précis, délais, signatures) ou d'un PV de non-conciliation (la voie contentieuse reste ouverte).
\end{enumerate}

\textbf{Débriefing (15 min) :} qu'est-ce qui a bloqué ? Quelle solution créative a émergé ? L'accord respecte-t-il le droit (un accord ne peut jamais valider une discrimination) ?
\end{experience}

\begin{aretenir}[Valeur de l'accord de médiation]
L'accord signé devant le médiateur engage les parties. En cas de non-exécution, la victime peut demander au juge de le rendre obligatoire. La médiation n'est donc pas une « parole en l'air » : c'est une voie de règlement reconnue, qui évite la lenteur et le coût du procès.
\end{aretenir}

\courssub{Analyse de recevabilité : filtrer les requêtes avant de saisir}

Saisir une juridiction ou la CNDHL sans vérifier la \cle{recevabilité}, c'est perdre du temps et décrédibiliser la cause. Le juge ou la Commission vérifie la recevabilité \textbf{d'office}, avant même d'examiner le fond.

\begin{propriete}[Conditions classiques de recevabilité]
\begin{enumerate}
    \item \textbf{Qualité pour agir} : être victime directe, ayant droit, ou organisation habilitée (ONG, syndicat pour ses membres).
    \item \textbf{Épuisement des voies de recours internes} : principe général devant les organes internationaux (sauf urgence, déni de justice ou procédure excessivement longue).
    \item \textbf{Délais} : respecter les délais de recours contentieux et saisir dans un délai raisonnable.
    \item \textbf{Objet} : la requête doit porter sur une violation d'un droit humain, non sur un simple litige privé sans dimension de discrimination ou d'abus de pouvoir.
    \item \textbf{Forme} : écrite, signée, accompagnée des pièces, rédigée dans une langue officielle (français ou anglais).
\end{enumerate}
\end{propriete}

\begin{exoresolu}[Analyse de deux requêtes fictives]
\textbf{Requête A :} Mme M. saisit le juge en urgence pour faire cesser la coupure d'eau de son quartier, décidée sans procédure par la société de distribution. Elle joint ses factures payées, ses courriers restés sans réponse et des photos.

\textbf{Requête B :} M. N. saisit directement la Cour africaine des droits de l'homme et des peuples pour contester sa condamnation pénale définitive, en invoquant la violation de son droit à un procès équitable.

\textbf{Travail à faire :} ces deux requêtes sont-elles recevables ? Justifiez.
\tcblower
\textbf{Analyse.}
\begin{itemize}
    \item \textbf{Requête A : RECEVABLE.} Qualité : usagère directement victime. Urgence : l'eau touche à la santé et à la vie. Objet : abus d'un service public. Pièces : complètes.
    \item \textbf{Requête B : IRRECEVABLE à ce stade.} Le Cameroun n'a pas fait la déclaration permettant aux particuliers de saisir directement la Cour africaine. Le particulier doit d'abord passer par la \textbf{Commission africaine des droits de l'homme et des peuples}, après avoir épuisé les voies de recours internes. Saisir directement la Cour est une erreur de voie de recours.
\end{itemize}
\end{exoresolu}

\begin{attention}[Erreur fatale au Probatoire : invoquer des voies de recours inexistantes]
Vérifiez toujours que la voie de recours invoquée est \textbf{ouverte au Cameroun} : un particulier ne saisit pas directement la Cour africaine ; il passe par la Commission africaine après épuisement des recours internes. De même, ne citez jamais un mécanisme international sans vérifier que le Cameroun l'a ratifié. Citer une voie inexistante, c'est perdre tous les points de la question « voies de recours ».
\end{attention}

\begin{savaistu}[La CNDHL aux côtés des victimes]
Créée par la loi n°90/53 du 19 décembre 1990 puis réorganisée par la loi n°2004/016 du 22 juillet 2004, la Commission nationale des droits de l'homme et des libertés (CNDHL) peut recevoir les plaintes, mener des enquêtes, visiter les lieux de détention, proposer des médiations et formuler des recommandations aux pouvoirs publics. C'est un interlocuteur accessible et gratuit pour les victimes qui n'ont pas les moyens de payer un avocat.
\end{savaistu}

\courssub{Construction d'un argumentaire écrit : la méthode P.Q.R.S.}

Que ce soit pour un plaidoyer écrit, une conclusion devant un juge ou une copie d'examen, l'argumentation juridique suit une structure rigoureuse. Nous la résumons par la méthode \cle{P.Q.R.S.} : Problème, Qualification, Raisonnement, Solution.

\begin{methode}[Structure type d'un argumentaire écrit (P.Q.R.S.)]
\begin{enumerate}
    \item \textbf{P — Problème} : formulez la question de droit en une phrase. \emph{Exemple : « L'expulsion de 50 familles sans décision de justice ni relogement constitue-t-elle une violation du droit au logement engageant la responsabilité de l'État ? »}
    \item \textbf{Q — Qualification} : énoncez les règles applicables en respectant la hiérarchie des normes : Constitution (Préambule) $>$ traités ratifiés (art. 45 de la Constitution) $>$ lois $>$ règlements. Citez les numéros d'articles.
    \item \textbf{R — Raisonnement} : appliquez la règle aux faits du dossier (« En l'espèce, le fait que... remplit les conditions de l'article... car... »). Anticipez les arguments de la partie adverse et réfutez-les.
    \item \textbf{S — Solution} : concluez par des demandes claires, numérotées et, si besoin, chiffrées (constat de la violation, cessation, réparation).
\end{enumerate}
\end{methode}

\begin{exemplebox}[Application de la méthode P.Q.R.S. au cas d'expulsion forcée]
\textbf{P :} L'expulsion sans décision de justice ni relogement viole-t-elle les droits des familles ?

\textbf{Q :} Droit au logement (art. 11 du Pacte international relatif aux droits économiques, sociaux et culturels) ; interdiction des traitements dégradants (art. 5 de la Charte africaine) ; Préambule de la Constitution de 1996.

\textbf{R :} En l'espèce, l'opération a été menée de nuit, sans titre ni préavis, avec violence (3 blessés) : les conditions de légalité ne sont pas remplies ; la présence de la police engage la responsabilité de l'État, tenu de protéger les droits.

\textbf{S :} Demander 1°) de constater l'illégalité de l'expulsion ; 2°) d'ordonner le relogement des familles ; 3°) d'allouer des réparations pour les préjudices matériel et moral.
\end{exemplebox}

\courssub{Auto-évaluation : grille de compétences}

Cette grille vous permet de situer votre niveau avant l'évaluation certificative. Soyez honnête : les items « Non acquis » ciblent vos révisions prioritaires.

\begin{center}
\begin{tabularx}{\linewidth}{|l|X|X|X|}\hline
\rowcolor{popblueL}
\textbf{Compétence} & \textbf{Non acquis} & \textbf{En cours} & \textbf{Acquis} \\ \hline
Citer les textes fondamentaux (Constitution, Pactes de 1966, Charte africaine, lois nationales) & Je confonds les textes. & Je cite les grands textes sans précision. & Je cite précisément les instruments et leurs articles clés. \\ \hline
Identifier la hiérarchie des normes au Cameroun & Je ne distingue pas Constitution, traité et loi. & Je connais la hiérarchie mais j'oublie la primauté des traités ratifiés (art. 45). & J'applique correctement la primauté des traités ratifiés sur les lois internes. \\ \hline
Connaître les institutions de protection (CNDHL, tribunaux, Commission africaine, organes ONU) & Je confonds les instances. & Je sais qui fait quoi, avec des hésitations sur les conditions de saisine. & Je maîtrise les compétences et les conditions de saisine de chaque organe. \\ \hline
Qualifier juridiquement une situation concrète & Je décris les faits sans qualification. & Je qualifie le droit mais j'oublie l'instance ou l'urgence. & Qualification complète : droit, texte, instance, urgence. \\ \hline
Rédiger une saisine structurée (F.D.D.) & Rédaction narrative, vague, sans pièces. & Structure présente mais demandes floues. & Rédaction rigoureuse : faits sourcés, droit exact, demandes concrètes. \\ \hline
Construire un argumentaire (P.Q.R.S.) & Pas de structure, mélange faits et droit. & Structure visible mais raisonnement faible. & Raisonnement complet, objections anticipées, conclusion nette. \\ \hline
Vérifier la recevabilité d'une requête & Je ne vérifie pas. & Je vérifie un ou deux critères. & Check-list complète avant tout dépôt. \\ \hline
\end{tabularx}
\end{center}

\begin{methode}[Check-list finale avant tout dépôt de saisine]
\begin{enumerate}
    \item \textbf{Voie de recours vérifiée ?} L'organe saisi est-il compétent et accessible au Cameroun ?
    \item \textbf{Délai respecté ?} Urgence caractérisée ? Recours internes épuisés si nécessaire ?
    \item \textbf{Pièces jointes ?} Copies lisibles, datées, listées (factures, photos, courriers, témoignages).
    \item \textbf{Adresse exacte du destinataire ?} Envoi recommandé avec accusé de réception ou dépôt contre décharge.
    \item \textbf{Double conservé ?} Gardez un exemplaire complet avec la preuve de dépôt.
    \item \textbf{Qualité pour agir prouvée ?} Victime directe, mandat, agrément de l'ONG.
\end{enumerate}
\end{methode}

\begin{aretenir}[Synthèse de la leçon — Les 3 piliers de l'effectivité des droits humains]
\begin{enumerate}
    \item \textbf{Qualifier} : transformer le fait brut en violation juridique fondée sur un texte (Constitution, traités ratifiés, lois).
    \item \textbf{Saisir} : choisir la \textbf{bonne porte} (juge en urgence, CNDHL pour enquête et médiation, Procureur pour les infractions, Commission africaine après épuisement des recours internes) avec un dossier complet, daté et précis.
    \item \textbf{Argumenter} : structurer (P.Q.R.S.), appliquer la règle aux faits, anticiper les objections, conclure par des demandes claires.
\end{enumerate}
Les droits humains ne sont pas une déclamation : c'est une \textbf{technique de protection}. Maîtriser cette technique, c'est rendre les droits réellement \emph{effectifs} pour tous.
\end{aretenir}

\newpage

\coursec{Activité d'intégration}

\courssub{Objectifs de l'activité}
\noindent Cette activité d'intégration vise à mobiliser l'ensemble des savoirs et savoir-faire de l'unité « Les droits humains » à travers une mise en situation professionnelle simulée. Elle permet à l'élève de :
\begin{itemize}
    \item \textbf{Qualifier} juridiquement des faits bruts en violations de droits humains (classification : droits civils et politiques, droits économiques sociaux et culturels, droits des peuples ; caractère absolu ou relatif, dérogeable ou non dérogeable).
    \item \textbf{Identifier} les niveaux de responsabilité (État, collectivités territoriales, acteurs privés) au regard des obligations de \cle{respecter}, de \cle{protéger} et de \cle{réaliser} les droits humains.
    \item \textbf{Hiérarchiser} les normes applicables (Constitution, traités internationaux, lois, règlements) pour fonder une argumentation solide.
    \item \textbf{Choisir} et \textbf{mettre en œuvre} la voie de recours la plus pertinente (CNDHL, juridictions nationales, mécanismes africains, plaidoyer) selon les critères de recevabilité, d'urgence et d'effectivité.
    \item \textbf{Rédiger} un acte clair, structuré et motivé (requête, saisine, recommandation).
    \item \textbf{Argumenter} à l'oral et confronter la légalité formelle à l'effectivité réelle des droits.
\end{itemize}

\courssub{Situation : « Opération \emph{Salongo} à Ngomgham » — Dossier d'instruction CNDHL}
\noindent \textbf{Contexte.} Nous sommes le 20 janvier 2026. La Commission Nationale des Droits de l'Homme et des Libertés (CNDHL) tient une session extraordinaire à Yaoundé pour examiner l'affaire « \emph{Population du quartier Ngomgham c. Préfet du Mfoundi, Maire de Yaoundé III, Société \emph{BTP Moderne} \& Cie} ».

\textbf{Les faits.} Le quartier Ngomgham (Yaoundé III) est un habitat informel dense (environ 2 500 habitants, majoritairement des jeunes, des femmes chefs de famille, des personnes handicapées, des travailleurs du secteur informel). Occupé depuis plus de 20 ans, il fait l'objet d'un litige foncier avec un promoteur privé (Société \emph{BTP Moderne} \& Cie) titulaire d'un titre foncier récent (TF 12458/Mfoundi, établi le 12/12/2025).

\textbf{Chronologie critique.}
\begin{itemize}
    \item \textbf{10 janv. 2026 :} Signature de l'\textbf{Arrêté préfectoral n\textsuperscript{o} 0045/AP/Y3} portant « libération des emprises domaniales et démolition des constructions anarchiques au quartier Ngomgham ». L'arrêté invoque l'« urgence salubre » et l'« intérêt général », ordonne l'exécution sous 48 heures par la force publique, sans mention de relogement ni d'indemnisation.
    \item \textbf{15 janv. 2026 (05 h 30) :} Intervention musclée : 3 bulldozers, 2 camions-bennes, effectifs mixtes (Police d'État, Police municipale, agents de \emph{BTP Moderne} armés de matraques). Pas de notification individuelle préalable aux occupants. Usage de gaz lacrymogène, tirs de sommation, destructions massives de biens (mobilier, documents d'identité, marchandises).
    \item \textbf{Bilan provisoire (18 janv.) :} 427 habitations détruites, 2 100 personnes sans abri (dont 680 enfants, 15 personnes à mobilité réduite). 12 blessés graves (ITT de 15 à 30 jours), 1 décès (M. \textsc{Ebode}, 62 ans, crise cardiaque lors de l'expulsion). 48 arrestations pour « rébellion » et « trouble à l'ordre public ».
\end{itemize}

\textbf{Pièces du dossier (extraits fournis à chaque commission).}

\begin{enumerate}
    \item \textbf{Plainte collective (16 janv.)} signée par 12 chefs de famille : « \emph{Nous n'avons reçu aucun préavis écrit. Nos titres d'occupation coutumiers et nos reçus de taxes municipales (patentes, taxes d'habitation payées à la mairie jusqu'en 2025) sont ignorés. Nos biens ont été pillés par les agents de la société. Les femmes ont été violentées. Nous demandons l'arrêt des opérations, le relogement, la réparation et la sanction des auteurs.} »
    \item \textbf{PV de Police n\textsuperscript{o} 089/YPJ (15 janv.)} : « \emph{Exécution de l'arrêté préfectoral. Résistance active de la population : jets de pierres, barricades. Usage proportionné de la force pour maintenir l'ordre. 48 interpellations. Aucune bavure constatée.} »
    \item \textbf{Arrêté préfectoral n\textsuperscript{o} 0045/AP/Y3 (10 janv.)} : Vu la loi n\textsuperscript{o} 74/1 du 6 juillet 1974 (régime foncier et domanial), le décret n\textsuperscript{o} 76/166 du 27 avril 1976 fixant les modalités de gestion du domaine national. \textbf{Article 1\textsuperscript{er} :} « Il est ordonné la libération immédiate... » \textbf{Article 3 :} « La force publique est réquisitionnée... » \textbf{Article 5 :} « Les frais sont à la charge de l'occupant illégal. » \emph{Aucun article sur le relogement, l'indemnisation ou le délai de préavis raisonnable.}
    \item \textbf{Certificats médicaux (Hôpital Central de Yaoundé, 15-16 janv.)} : 12 certificats mentionnant : « \emph{Traumatismes crâniens, fractures des membres, lésions oculaires par gaz, syndromes post-traumatiques} ». ITT moyenne : 21 jours.
    \item \textbf{Rapport d'enquête de l'ONG \emph{Habitat pour Tous} (18 janv.)} : « \emph{Violation du droit au logement décent (art. 11 PIDESC, art. 14 et 16 CADHP), du droit à la vie privée et familiale (art. 17 PIDCP), de l'interdiction des traitements inhumains (art. 5 CADHP, art. 7 PIDCP), du principe de non-discrimination (femmes, personnes handicapées), du droit de propriété (biens détruits sans indemnisation préalable). Absence de \cle{procédure équitable} : pas de consultation, pas de relogement, préavis déraisonnable (48 h). L'arrêté est entaché d'excès de pouvoir : une mesure d'éviction sans mesure d'accompagnement viole le \cle{principe de proportionnalité} et l'\cle{obligation de protection} de l'État.} »
    \item \textbf{Presse (Mutations du 19 janv., Cameroon Tribune du 20 janv.)} : \emph{Mutations} : « \emph{Ngomgham : la barbarie au nom de la loi. Témoignages accablants.} » \emph{Cameroon Tribune} : « \emph{Salubrité publique : l'État restaure l'ordre foncier. Le Préfet salue le professionnalisme des forces de l'ordre.} »
\end{enumerate}

\courssub{Consignes}
\noindent Chaque commission (3 à 4 élèves) dispose de 2 h 30 pour traiter le dossier et préparer la plénière. Le travail est structuré en 5 étapes obligatoires :

\begin{enumerate}
    \item \textbf{Grille de qualification des violations (30 min).} Complétez le tableau ci-dessous pour chaque fait identifié. Citez la norme violée (Constitution, PIDCP, PIDESC, CADHP, loi nationale) et classez le droit (DCP / DESC / droits des peuples ; absolu / relatif ; dérogeable / non dérogeable).
    \begin{center}
    \begin{tabularx}{\linewidth}{|l|X|X|X|X|}\hline
    \textbf{Fait / Acte} & \textbf{Droit(s) violé(s)} & \textbf{Source(s) juridique(s)} & \textbf{Classification} & \textbf{Caractère (Absolu/Relatif)} \\ \hline
    Ex. : Destruction des maisons sans relogement & Droit au logement décent & Art. 11 PIDESC, art. 14 et 16 CADHP & DESC & Relatif (mais cœur intangible) \\ \hline
    Usage de gaz lacrymogène / tirs / blessés &  &  &  &  \\ \hline
    Arrestations massives sans base légale &  &  &  &  \\ \hline
    Absence de préavis / consultation / indemnisation &  &  &  &  \\ \hline
    Destruction des biens personnels (papiers, marchandises) &  &  &  &  \\ \hline
    Impact discriminatoire (femmes, handicapés, pauvres) &  &  &  &  \\ \hline
    \end{tabularx}
    \end{center}

    \item \textbf{Analyse des responsabilités (20 min).} Identifiez pour chaque violation : l'auteur matériel, l'auteur intellectuel ou hiérarchique, le fondement de la responsabilité (faute de service, faute personnelle détachable du service, responsabilité sans faute pour rupture de l'égalité devant les charges publiques, responsabilité délictuelle, complicité). Distinguez : État (Police, Préfet), Collectivité (Maire), Personne privée (Société \emph{BTP Moderne}).

    \item \textbf{Choix de la stratégie contentieuse et extra-contentieuse (30 min).} Parmi les 4 options ci-dessous, choisissez \textbf{deux} actions prioritaires (une judiciaire/contentieuse, une quasi-judiciaire/institutionnelle) et justifiez le choix par une analyse coûts-avantages (délai, effet suspensif, force exécutoire, portée, risque d'irrecevabilité).
    \begin{itemize}
        \item A. \textbf{Recommandation CNDHL} (loi n\textsuperscript{o} 2019/014 du 24 décembre 2019) : avis motivé, publication, suivi. Non contraignante.
        \item B. \textbf{Référé-suspension / référé-liberté devant le Tribunal Administratif du Centre} (loi n\textsuperscript{o} 2006/022 du 29 décembre 2006 fixant l'organisation et le fonctionnement des tribunaux administratifs) : suspension de l'arrêté, ordonnance de mesures d'urgence. Décision exécutoire. Délai : 48 h à 15 jours.
        \item C. \textbf{Saisine de la Commission Africaine des Droits de l'Homme et des Peuples (CADHP)} (art. 55 et 56 de la Charte africaine) : mesures provisoires possibles. Portée internationale. Délai long (mois, années). Épuisement des voies de recours internes requis (sauf exceptions).
        \item D. \textbf{Plaidoyer médiatique \& mobilisation citoyenne} : pression politique, opinion publique. Risques (poursuites pour diffamation, trouble à l'ordre public).
    \end{itemize}

    \item \textbf{Rédaction de l'acte de saisine principal (45 min).} Rédigez \textbf{une page maximum} de l'acte choisi comme priorité n\textsuperscript{o} 1 (ex. : requête en référé-suspension au TA \textbf{OU} saisine formelle de la CNDHL). Structure imposée : en-tête (qualité, représentation), objet, exposé des faits, discussion (moyens de fait et de droit : violation de la Constitution, des traités, des lois, du principe de proportionnalité, de la procédure équitable), dispositif (demandes chiffrées : suspension, relogement, provision à valoir sur dommages-intérêts, astreinte), signature.

    \item \textbf{Préparation de la plaidoirie orale (15 min).} 5 minutes maximum par commission en plénière. Structure : accroche (1 phrase choc) $\rightarrow$ 3 arguments juridiques forts $\rightarrow$ demande précise $\rightarrow$ réponse anticipée à l'argument adverse (« trouble à l'ordre public »).
\end{enumerate}

\courssub{Ressources à mobiliser}
\noindent \textbf{Rappel synthétique des outils juridiques du cours.}

\begin{definition}[Hiérarchie des normes au Cameroun (art. 45 Const. 1996)]
La Constitution $\succ$ les traités et accords internationaux régulièrement ratifiés (PIDCP, PIDESC, CADHP, CIDE, CEDAW...) $\succ$ les lois organiques $\succ$ les lois ordinaires $\succ$ les ordonnances $\succ$ les décrets $\succ$ les arrêtés. \textbf{Conséquence :} un arrêté préfectoral contraire à la Charte africaine ou à la Constitution est illégal et peut être annulé ou suspendu par le juge.
\end{definition}

\begin{propriete}[Justiciabilité des DESC au Cameroun]
Les Droits Économiques, Sociaux et Culturels (logement, santé, éducation) sont \textbf{invocables devant un juge} :
\begin{itemize}
    \item \textbf{Directement :} via la primauté des traités ratifiés (art. 45 Const.) et l'obligation de \cle{réalisation progressive} des droits (disponibilité, accessibilité, acceptabilité, qualité).
    \item \textbf{Indirectement :} par le biais des droits civils et politiques (droit à la vie, dignité, propriété, non-discrimination, procédure équitable). \emph{Ex. : une expulsion sans relogement porte atteinte au droit à la vie (art. 4 CADHP) et à la dignité (Préambule de la Constitution).}
\end{itemize}
\end{propriete}

\begin{aretenir}[Institutions compétentes \& voies de recours]
\begin{itemize}
    \item \textbf{CNDHL (loi n\textsuperscript{o} 2019/014 du 24 décembre 2019) :} saisine libre (toute victime, ONG, tiers). Avis, recommandations, médiation, rapports publics. \textbf{Atout :} rapidité, gratuité, expertise en droits humains. \textbf{Limite :} recommandations non contraignantes.
    \item \textbf{Tribunal Administratif (loi n\textsuperscript{o} 2006/022 du 29 décembre 2006) :} recours pour excès de pouvoir (annulation de l'arrêté), référé-suspension (urgence + doute sérieux sur la légalité), référé-liberté (atteinte grave et manifestement illégale à une liberté fondamentale). \textbf{Atout :} décision contraignante, effet suspensif possible. \textbf{Condition :} intérêt à agir (direct, personnel, certain).
    \item \textbf{CADHP (Charte africaine de 1981, ratifiée par le Cameroun) :} saisine après épuisement des voies internes (sauf déni de justice ou délai excessif). Mesures provisoires possibles. \textbf{Atout :} portée supranationale, recommandations adressées à l'État partie.
\end{itemize}
\end{aretenir}

\begin{methode}[Rédaction d'un acte de procédure (requête / saisine)]
\begin{enumerate}
    \item \textbf{En-tête :} juridiction ou institution saisie, référence du dossier, qualité des parties (demandeur/défendeur), représentation (avocat ou mandataire).
    \item \textbf{Objet :} précis (ex. : « Requête en référé-suspension de l'arrêté n\textsuperscript{o} 0045/AP/Y3 du 10/01/2026 »).
    \item \textbf{Exposé des faits :} chronologique, précis, sourcé (pièces jointes), sans invective.
    \item \textbf{Moyens (discussion) :} \textbf{moyens de fait} (constatations) + \textbf{moyens de droit} (violation d'une norme supérieure). Classer : 1. incompétence / vice de forme ; 2. violation des droits civils et politiques (vie, dignité, procédure équitable) ; 3. violation des DESC (logement — cœur intangible) ; 4. non-discrimination ; 5. disproportionnalité.
    \item \textbf{Dispositif :} demandes numérotées, chiffrées, avec astreinte (ex. : « 500 000 FCFA par jour de retard »).
    \item \textbf{Formule de clôture et signature.}
\end{enumerate}
\end{methode}

\courssub{Démarche et corrigé commenté}
\begin{exoresolu}[Corrigé type — Commission 1 (Modèle « Référé-TA + Saisine CNDHL »)]
Énoncé : traiter les 5 étapes du dossier « Ngomgham » (qualification, responsabilités, stratégie, acte, plaidoirie).
\tcblower
\textbf{Étape 1 — Grille de qualification (extrait significatif).}

\begin{center}
\begin{tabularx}{\linewidth}{|l|X|X|X|X|}\hline
\textbf{Fait / Acte} & \textbf{Droit(s) violé(s)} & \textbf{Source(s) juridique(s)} & \textbf{Classification} & \textbf{Caractère} \\ \hline
Destruction de 427 maisons sans relogement ni indemnisation préalable & Droit au logement décent ; droit de propriété (biens) ; dignité humaine & Art. 11 PIDESC ; art. 14 et 16 CADHP ; Préambule Const. 1996 & DESC (droit à part entière) ; DCP (propriété) & \textbf{Relatif} (mais \textbf{cœur intangible} : interdiction des expulsions forcées sans protection) \\ \hline
Usage de gaz lacrymogène, tirs, 12 blessés graves, 1 décès & Droit à la vie ; intégrité physique ; interdiction des traitements cruels, inhumains ou dégradants ; sécurité de la personne & Art. 4 et 5 CADHP ; art. 6 et 7 PIDCP ; Préambule Const. & DCP (droits civils et politiques) & \textbf{Absolu} (vie, torture) et \textbf{non dérogeable} (art. 4 PIDCP) \\ \hline
Arrestation de 48 personnes (« rébellion ») sans procès-verbal individuel motivé & Liberté individuelle ; procédure équitable ; présomption d'innocence & Art. 6 et 7 CADHP ; art. 9 PIDCP ; Préambule Const. & DCP & \textbf{Relatif} (mais garanties procédurales \textbf{imprescriptibles} en matière de détention) \\ \hline
Préavis de 48 h, absence de consultation, absence de relogement & Procédure équitable ; droit à l'information ; participation & Art. 7 CADHP ; art. 14 PIDCP ; Principes directeurs ONU sur les expulsions forcées (1997) & DCP (garanties) & \textbf{Absolu} (garanties d'une procédure régulière) \\ \hline
Destruction des papiers d'identité, marchandises, outils de travail & Droit de propriété ; droit au travail ; dignité & Art. 14 CADHP ; art. 15 CADHP ; art. 17 PIDCP & DCP / DESC & \textbf{Relatif} (mais protection immédiate due) \\ \hline
Impact disproportionné sur les femmes chefs de famille, les personnes handicapées, les pauvres & Non-discrimination ; égalité devant la loi ; protection des groupes vulnérables & Art. 2 et 3 CADHP ; art. 2 et 26 PIDCP ; art. 18 CADHP ; Préambule Const. & Transversal (principe général) & \textbf{Absolu} (interdiction de la discrimination) \\ \hline
\end{tabularx}
\end{center}

\textbf{Commentaire :} la qualification croisée (DCP + DESC) est la clé. L'expulsion viole le \textbf{cœur intangible} du droit au logement (interdiction des expulsions forcées sans garanties) et des droits \textbf{absolus et non dérogeables} (vie, interdiction des traitements inhumains, procédure équitable). Cela ouvre la voie du \textbf{référé-liberté}, qui ne demande pas un simple doute sérieux mais une atteinte \emph{grave et manifestement illégale} à une liberté fondamentale.

\vspace{0.5cm}
\textbf{Étape 2 — Analyse des responsabilités.}

\begin{center}
\begin{tabularx}{\linewidth}{|l|X|X|}\hline
\textbf{Acteur} & \textbf{Fondement de responsabilité} & \textbf{Actes imputables} \\ \hline
\textbf{État (Préfet / Police)} & \textbf{Faute de service} (organisation défaillante, absence de mesures d'accompagnement) + \textbf{faute personnelle détachable du service} (violences volontaires). Responsabilité sans faute possible (rupture de l'égalité devant les charges publiques : les habitants supportent seuls le coût de l'intérêt général). & Signature d'un arrêté illégal (excès de pouvoir) ; réquisition de la force publique sans encadrement ; absence de protection des vies et des biens. \\ \hline
\textbf{Commune (Maire de Yaoundé III)} & \textbf{Faute de service} (police municipale) + \textbf{abstention fautive} (défaut de relogement ; perception de taxes sans protection). & Participation de la police municipale ; non-exercice du pouvoir de police administrative pour protéger la population ; perception de taxes (reconnaissance de fait de l'occupation). \\ \hline
\textbf{Société \emph{BTP Moderne} \& Cie} & \textbf{Responsabilité délictuelle} (dommages causés à autrui) + \textbf{complicité} des violences (agents armés). Personne privée participant à l'exécution d'une mission de puissance publique (démolition) $\rightarrow$ responsabilité engagée. & Exécution brutale (bulldozers sur les biens) ; pillages par ses agents ; absence de vérification des droits des occupants (titres coutumiers, taxes payées). \\ \hline
\end{tabularx}
\end{center}

\vspace{0.5cm}
\textbf{Étape 3 — Choix stratégique : Priorité 1 : référé-liberté devant le TA du Centre. Priorité 2 : saisine de la CNDHL.}

\textbf{Justification argumentée :}
\begin{itemize}
    \item \textbf{Référé-liberté (TA) :} \emph{atteinte grave et manifestement illégale} aux libertés fondamentales (vie, art. 4 CADHP ; intégrité, art. 5 ; cœur intangible du droit au logement ; procédure équitable). \textbf{Atouts :} juge unique, délai très court (48 h), ordonnance \emph{exécutoire} (astreinte possible), suspension immédiate de l'arrêté, injonction de relogement d'urgence. \textbf{Risque :} recevabilité (intérêt à agir des 12 plaignants : direct, personnel et certain $\rightarrow$ recevable). \textbf{Écarté :} le référé-suspension exige un « doute sérieux » sur la légalité ; ici l'illégalité est \emph{manifeste} (aucun relogement prévu), le référé-liberté est plus adapté et plus rapide.
    \item \textbf{Saisine de la CNDHL :} parallèle, non exclusive. \textbf{Atouts :} enquête indépendante, médiation, recommandation publique (pression politique), rapport transmis aux plus hautes autorités, saisine possible de la CADHP ultérieurement. \textbf{Complémentarité :} le TA apporte la contrainte juridique immédiate ; la CNDHL apporte la vérité institutionnelle, la réparation morale et le suivi de l'effectivité.
    \item \textbf{Écartées :} la CADHP (voies internes non épuisées, délai long) ; le plaidoyer seul (insécurité juridique, risque de répression).
\end{itemize}

\vspace{0.5cm}
\textbf{Étape 4 — Acte de saisine principal : requête en référé-liberté (modèle réduit).}

\begin{quote}
\small
\textbf{RÉPUBLIQUE DU CAMEROUN}\\
\textbf{Paix — Travail — Patrie}\\
\textbf{TRIBUNAL ADMINISTRATIF DU CENTRE}\\
\textbf{REQUÊTE EN RÉFÉRÉ-LIBERTÉ}\\
\textbf{N\textsuperscript{o} RG : 2026/00458}\\\vspace{0.2cm}

\textbf{VU} la Constitution du 18 janvier 1996 ; \textbf{VU} la loi n\textsuperscript{o} 2006/022 du 29 décembre 2006 fixant l'organisation et le fonctionnement des tribunaux administratifs ; \textbf{VU} le PIDCP, le PIDESC et la CADHP (régulièrement ratifiés) ; \textbf{VU} l'arrêté préfectoral n\textsuperscript{o} 0045/AP/Y3 du 10/01/2026 ; \textbf{VU} la requête enregistrée le 20/01/2026 ; \textbf{VU} les pièces jointes (plainte collective, certificats médicaux, PV de police, rapport d'ONG, photographies).\\\vspace{0.2cm}

\textbf{ENTRE :}\\
\textbf{Demandeurs :} 1. Mme \textsc{Mendouga} Marie, veuve, 3 enfants, résidente à Ngomgham (pièce n\textsuperscript{o} 1 : CNI). 2. M. \textsc{Fouda} Jean, personne à mobilité réduite, résident à Ngomgham. \emph{(et 10 autres signataires, liste en annexe)}.\\
\textbf{Représentés par Me} \textsc{Avocat}, Avocat au Barreau du Cameroun, mandat n\textsuperscript{o} 12/2026.\\\vspace{0.2cm}
\textbf{ET :}\\
\textbf{Défendeurs :} 1. \textbf{État du Cameroun} (Ministère de l'Administration Territoriale), représenté par le Préfet du Mfoundi. 2. \textbf{Commune d'arrondissement de Yaoundé III}, représentée par son Maire. 3. \textbf{Société \emph{BTP Moderne} \& Cie}, représentée par son Directeur Général.\\\vspace{0.3cm}

\textbf{OBJET :} \textbf{Requête en référé-liberté} — Suspension de l'exécution de l'arrêté n\textsuperscript{o} 0045/AP/Y3 — Injonction de relogement d'urgence — Provision à valoir sur dommages-intérêts.\\\vspace{0.3cm}

\textbf{EXPOSÉ DES FAITS}\\
Les 15 et 16 janvier 2026, sur le fondement de l'arrêté attaqué, les forces de l'ordre et des agents de la société \emph{BTP Moderne} ont procédé à la démolition violente du quartier Ngomgham (427 habitations, 2 100 personnes). Aucune notification individuelle, aucun préavis raisonnable (48 h), aucune mesure de relogement, aucune indemnisation préalable. Bilan : 1 décès, 12 blessés graves (ITT de 15 à 30 jours), 48 arrestations arbitraires, destructions massives de biens personnels (papiers d'identité, marchandises). Les demandeurs, occupants de bonne foi depuis plus de 20 ans (taxes payées à la mairie), se retrouvent sans abri, en pleine saison des pluies, exposés aux épidémies.\\\vspace{0.3cm}

\textbf{DISCUSSION}\\\vspace{0.1cm}
\textbf{SUR LA RECEVABILITÉ :} les demandeurs justifient d'un intérêt direct, personnel, certain et né (destruction du domicile et des biens, atteinte à l'intégrité physique). La qualité pour agir est établie.\\\vspace{0.1cm}
\textbf{SUR LE FONDEMENT (RÉFÉRÉ-LIBERTÉ) :} l'atteinte est \textbf{grave} (vie, intégrité, domicile, biens, dignité) et \textbf{manifestement illégale} :\\
\textbf{1. Violation du droit à la vie et à l'intégrité physique (droits absolus, non dérogeables).} L'usage disproportionné de la force (gaz lacrymogène, tirs, bulldozers sur des habitations occupées à l'aube) a causé la mort de M. \textsc{Ebode} et 12 blessés graves. Violation des art. 4 et 5 CADHP, art. 6 et 7 PIDCP, Préambule de la Constitution. L'État a manqué à son obligation de \cle{protéger}.\\
\textbf{2. Violation du cœur intangible du droit au logement (DESC invocable).} L'expulsion forcée sans relogement, sans indemnisation et sans consultation viole l'obligation de \cle{respecter} (ne pas détruire) et de \cle{protéger} (contre les tiers). Art. 11 PIDESC, art. 14 et 16 CADHP, Principes directeurs des Nations Unies sur les expulsions forcées. En ne prévoyant aucune mesure d'accompagnement, l'arrêté est entaché d'\textbf{excès de pouvoir} et de \textbf{détournement de pouvoir} (servir l'intérêt privé du promoteur).\\
\textbf{3. Violation de la procédure équitable et des garanties des personnes (droits absolus).} Préavis de 48 h déraisonnable, absence de notification individuelle, absence de voie de recours effective préalable. Violation des art. 7 CADHP et 14 PIDCP.\\
\textbf{4. Discrimination indirecte.} Impact majeur sur les femmes chefs de famille, les personnes handicapées et les plus pauvres. Violation des art. 2, 3 et 18 CADHP, art. 26 PIDCP.\\
\textbf{5. Responsabilité de la société privée.} Exécution d'une mission de puissance publique sans encadrement, violences et pillages $\rightarrow$ responsabilité solidaire avec l'État.\\\vspace{0.3cm}

\textbf{DISPOSITIF}\\
\textbf{Les demandeurs demandent au Juge des référés, statuant en audience publique et contradictoirement :}\\
\textbf{1.} De \textbf{suspendre} immédiatement l'exécution de l'arrêté n\textsuperscript{o} 0045/AP/Y3 du 10/01/2026 et toute opération de démolition ou d'expulsion sur le site de Ngomgham.\\
\textbf{2.} D'\textbf{enjoindre} à l'État (Préfet) et à la Commune (Maire) de mettre en œuvre, sous astreinte de \textbf{500 000 FCFA par jour de retard} à compter de la signification, un plan de relogement d'urgence décent (abris, eau potable, latrines, sécurité) pour les 2 100 victimes, dans un délai de \textbf{72 heures}.\\
\textbf{3.} D'\textbf{enjoindre} à l'État de faire cesser toute violence, de libérer les 48 personnes interpellées (sauf procédure judiciaire régulière distincte) et de garantir l'intégrité physique des occupants.\\
\textbf{4.} De \textbf{condamner} solidairement l'État, la Commune et la Société \emph{BTP Moderne} à verser aux demandeurs une provision de \textbf{5 000 000 FCFA} chacun à valoir sur les dommages-intérêts (préjudice moral, matériel, frais médicaux, perte de revenus), sans préjudice de l'évaluation définitive.\\
\textbf{5.} D'\textbf{ordonner} l'exécution provisoire de la présente ordonnance, nonobstant tout recours.\\
\textbf{6.} De \textbf{mettre} les dépens à la charge des défendeurs.\\\vspace{0.3cm}
Fait à Yaoundé, le 20 janvier 2026.\\
\textbf{Me} \textsc{Avocat}, Conseil des demandeurs.\\
\textit{(Signature)}
\end{quote}

\textbf{Commentaire sur la rédaction :} l'acte respecte la structure imposée. Les moyens sont hiérarchisés : les droits absolus (vie, traitements inhumains, procédure équitable) en premier, pour verrouiller les conditions du référé-liberté. Le droit au logement (DESC) est invoqué via son \emph{cœur intangible} (interdiction des expulsions forcées) pour écarter l'argument de la « relativité » des droits sociaux. L'astreinte (500 000 FCFA/jour) et la provision (5 000 000 FCFA) sont chiffrées pour garantir l'effectivité. La solidarité des défendeurs (État, Commune, privé) est revendiquée.

\vspace{0.5cm}
\textbf{Étape 5 — Plaidoirie orale (simulation, 5 min).}

\begin{quote}
\small
\textbf{Accroche :} « Monsieur le Juge, à 5 h 30 du matin, pendant que la loi dormait, des bulldozers ont écrasé le droit à la vie de M. Ebode et le toit de 2 100 Camerounais. »\\
\textbf{Argument 1 (légalité) :} « L'arrêté est un leurre juridique. Il invoque l'urgence salubre mais organise une expulsion sans relogement. C'est une violation manifeste de l'art. 16 de la Charte africaine et du principe de proportionnalité. L'illégalité saute aux yeux : pas de relogement, pas de légalité. »\\
\textbf{Argument 2 (droits absolus) :} « On nous oppose le trouble à l'ordre public. Mais la vie de M. Ebode, les fractures de 12 personnes, la terreur des enfants : ce ne sont pas des dommages collatéraux, ce sont des atteintes à des droits absolus (art. 4 et 5 CADHP). Les droits absolus ne se négocient pas, ils s'imposent. »\\
\textbf{Argument 3 (effectivité) :} « Une annulation dans six mois ne ressuscitera pas M. Ebode et ne reconstruira pas les maisons. Seul le référé-liberté, par son effet immédiat et son astreinte, peut rendre la justice effective \emph{aujourd'hui}. »\\
\textbf{Demande :} « Suspension immédiate, relogement sous 72 heures sous astreinte, provision de 5 000 000 FCFA. »\\
\textbf{Réponse adverse anticipée :} « Le Préfet dira "trouble à l'ordre public". Nous répondons : l'ordre public, c'est d'abord le respect de la Constitution et des traités. L'État ne peut pas créer le trouble (expulsion illégale) pour s'en prévaloir ensuite (résistance) et justifier la répression. C'est un cercle vicieux que le Juge doit briser. »
\end{quote}

\end{exoresolu}

\courssub{Bilan de l'activité}
\noindent Cette simulation a permis de vérifier l'acquisition des compétences cibles du Probatoire technique :
\begin{itemize}
    \item \textbf{Savoir mobilisé :} la classification fine (DESC/DCP, absolu/relatif, cœur intangible) n'est pas académique : elle dicte le choix du juge (référé-liberté plutôt que référé-suspension) et la force de l'argumentaire. La hiérarchie des normes (art. 45 Const.) est l'outil qui permet d'invalider l'arrêté préfectoral.
    \item \textbf{Savoir-faire « Qualifier » :} la grille oblige à sortir de l'émotion pour nommer juridiquement chaque fait (ex. : « destruction des papiers » = atteinte à la propriété + dignité + procédure équitable). La responsabilité est partagée (État, Commune, privé) selon des régimes distincts (faute de service, faute personnelle détachable, responsabilité délictuelle, complicité).
    \item \textbf{Savoir-faire « Saisir » :} le choix stratégique (TA + CNDHL) montre la compréhension de la \cle{complémentarité des mécanismes} : contrainte juridique immédiate (TA) + vérité institutionnelle et suivi (CNDHL). La rédaction de l'acte (en-tête, objet, faits, moyens classés, dispositif chiffré) respecte le formalisme exigé.
    \item \textbf{Savoir-faire « Argumenter » :} l'oral impose la synthèse, la hiérarchie des arguments (les droits absolus d'abord) et l'anticipation de la partie adverse (l'argument du « trouble à l'ordre public » retourné).
    \item \textbf{Posture citoyenne :} le débat final « Effectivité contre Légalité » révèle que le droit n'est pas auto-exécutant. L'astreinte, la provision, la médiation de la CNDHL et le plaidoyer sont les leviers qui transforment la norme en réalité pour les 2 100 victimes de Ngomgham.
\end{itemize}

\begin{aretenir}[Grille d'évaluation des compétences (indicative — sur 20 points)]
\begin{center}
\begin{tabularx}{\linewidth}{|l|X|c|}\hline
\textbf{Compétence évaluée} & \textbf{Indicateurs de réussite (niveau Probatoire technique)} & \textbf{Pts} \\ \hline
\textbf{Qualification juridique} & Grille complète, sources exactes (articles des traités, Constitution, lois), classification rigoureuse (absolu/relatif, cœur intangible), distinction DCP/DESC. & 5 \\ \hline
\textbf{Analyse des responsabilités} & Identification des acteurs, fondements juridiques distincts (faute de service, faute personnelle détachable, responsabilité sans faute, complicité), articulation État/privé. & 4 \\ \hline
\textbf{Choix stratégique} & Justification pertinente (délai, effet, recevabilité, portée), articulation contentieux / quasi-judiciaire, exclusion motivée des autres voies. & 3 \\ \hline
\textbf{Rédaction de l'acte} & Structure formelle respectée, moyens hiérarchisés (absolus $\rightarrow$ cœur des DESC $\rightarrow$ procédure), dispositif précis, chiffré, avec astreinte, style juridique. & 5 \\ \hline
\textbf{Argumentation orale} & Temps respecté (5 min), structure claire (accroche, 3 arguments, demande, réponse adverse), vocabulaire précis, posture professionnelle, réactivité. & 3 \\ \hline
\textbf{Travail collaboratif} & Répartition des rôles (rapporteur, rédacteur, plaideur, juriste chargé des sources), cohérence du rendu final, respect des délais. & \textbf{--} \\ \hline
\end{tabularx}
\end{center}
\noindent \emph{Note :} le travail collaboratif est apprécié par l'enseignant (observation) et peut bonifier la note individuelle. L'activité valide les attendus du TD1 (étude d'une structure locale de promotion des droits de l'homme) par la simulation concrète de la saisine de la CNDHL.
\end{aretenir}

\newpage

\coursec{Méthodes et exercices résolus}

\coursec{Leçon 01 : Les droits humains — Méthodes, Cas pratiques et Exercices résolus}

\begin{methode}[Qualifier une situation au regard des caractéristiques et catégories des droits de l'homme]
    \textbf{Objectif :} Appliquer les notions de l'unité (universalité, indivisibilité, interdépendance, classification) à une situation concrète.
    \begin{enumerate}
        \item \textbf{Identifier les faits :} Relever les atteintes ou les revendications mentionnées dans le cas pratique (ex. : arrestation arbitraire, interdiction de manifester, non-paiement de salaire, pollution d'un cours d'eau).
        \item \textbf{Qualifier le droit violé ou revendiqué :} Nommer le droit précis (ex. : droit à la liberté et à la sécurité, liberté d'association, droit au travail, droit à un environnement sain).
        \item \textbf{Classer le droit :} Rattacher ce droit à sa génération/catégorie :
            \begin{itemize}
                \item \cle{1\textsuperscript{re} génération (Droits civils et politiques) :} Libertés individuelles, droits procéduraux (vie, liberté, opinion, vote). \emph{Devoir d'abstention de l'État}.
                \item \cle{2\textsuperscript{e} génération (Droits économiques, sociaux et culturels - DESC) :} Droit au travail, santé, éducation, logement, culture. \emph{Devoir d'action/prestation de l'État}.
                \item \cle{3\textsuperscript{e} génération (Droits de solidarité/peuples) :} Droit au développement, paix, environnement sain, patrimoine commun. \emph{Devoir de coopération internationale}.
            \end{itemize}
        \item \textbf{Vérifier les caractéristiques :} Préciser si le droit est \cle{universel} (s'applique à tous), \cle{inaliénable} (ne peut être cédé), \cle{indivisible/interdépendant} (violation de l'un affecte les autres).
        \item \textbf{Conclure :} Formuler une phrase de qualification juridique complète.
    \end{enumerate}
\end{methode}

\begin{exoresolu}[Qualification juridique d'une situation d'atelier]
    \textbf{Énoncé :} Dans une entreprise de BTP à Douala, M. Ngono, délégué du personnel, est licencié sans préavis ni indemnité après avoir organisé une grève pour réclamer le port d'équipements de protection individuelle (EPI) conformes aux normes. La direction invoque une « faute grave » pour trouble à l'ordre public. Au même moment, les riverains se plaignent de la pollution sonore nocturne du chantier.
    \tcblower
    \textbf{Solution détaillée :}
    \begin{enumerate}
        \item \textbf{Analyse des faits :} Trois situations distinctes : (1) Licenciement d'un délégué suite à une action syndicale. (2) Revendication d'EPI (sécurité au travail). (3) Pollution sonore affectant les riverains.
        \item \textbf{Qualification et classification :}
            \begin{itemize}
                \item \textbf{Situation 1 :} Atteinte à la \cle{liberté syndicale} (droit de grève, protection des délégués) et au \cle{droit au travail} (procédure de licenciement abusive). Ce sont des \cle{droits civils et politiques (1\textsuperscript{re} génération)} pour la liberté syndicale et \cle{DESC (2\textsuperscript{e} génération)} pour la protection contre le licenciement abusif.
                \item \textbf{Situation 2 :} Revendication d'EPI = \cle{droit à la santé} et \cle{droit à des conditions de travail justes et satisfaisantes} (Art. 7 PIDESC, Code du travail camerounais). C'est un \cle{Droit Économique, Social et Culturel (2\textsuperscript{e} génération)}. L'État a une obligation de \emph{due diligence} pour faire respecter ce droit par l'employeur privé.
                \item \textbf{Situation 3 :} Pollution sonore = Atteinte au \cle{droit à un environnement sain} (Art. 24 Charte africaine, Loi camerounaise sur l'environnement). C'est un \cle{droit de solidarité / 3\textsuperscript{e} génération} (droit des peuples/riverains).
            \end{itemize}
        \item \textbf{Caractéristiques mobilisées :}
            \begin{itemize}
                \item \cle{Indivisibilité :} La violation de la liberté syndicale (1\textsuperscript{re} gén.) empêche la défense des DESC (2\textsuperscript{e} gén.).
                \item \cle{Interdépendance :} L'absence d'EPI (DESC) menace le droit à la vie/santé (1\textsuperscript{re} gén.).
                \item \cle{Universalité :} Les riverains (non-employés) sont titulaires du droit à l'environnement.
            \end{itemize}
        \item \textbf{Conclusion :} L'entreprise viole simultanément des droits de trois générations. Le licenciement de M. Ngono est une violation de la liberté syndicale (1\textsuperscript{re} gén.) et du droit au travail (2\textsuperscript{e} gén.). L'absence d'EPI viole le droit à la santé (2\textsuperscript{e} gén.). La pollution viole le droit à l'environnement (3\textsuperscript{e} gén.).
    \end{enumerate}
\end{exoresolu}

\begin{methode}[Analyser la hiérarchie des normes et la justiciabilité des droits au Cameroun]
    \textbf{Objectif :} Situer les textes juridiques dans la pyramide de Kelsen camerounaise et déterminer la voie de recours (justiciabilité).
    \begin{enumerate}
        \item \textbf{Rappeler la hiérarchie constitutionnelle (Art. 45 Const. 1996 modifiée) :}
            \begin{enumerate}
                \item \cle{La Constitution} (sommet).
                \item \cle{Les Traités et Accords internationaux} régulièrement ratifiés/publiés (supériorité sur les lois, sous réserve de réciprocité).
                \item \cle{Les Lois} (votées par le Parlement).
                \item \cle{Les Décrets d'application} (Premier Ministre/Président).
                \item \cle{Les Arrêtés} (Ministres, Gouverneurs, Préfets, Maires).
            \end{enumerate}
        \item \textbf{Identifier la norme invoquée :} Est-ce un article de la Constitution ? Une disposition d'un traité (PIDCP, PIDESC, Charte Africaine, CEDH) ? Une loi nationale (Code du travail, Code pénal, Loi sur l'environnement) ?
        \item \textbf{Vérifier la justiciabilité (possibilité de saisir un juge) :}
            \begin{itemize}
                \item \cle{Droits civils/politiques (1\textsuperscript{re} gén.) :} Généralement \textbf{justiciables immédiatement} (obligation de respect/protection). Ex: Habeas corpus, liberté d'expression.
                \item \cle{DESC (2\textsuperscript{e} gén.) :} \textbf{Justiciabilité progressive} (Art. 2 PIDESC : « s'engage à agir... au maximum de ses ressources »). Au Cameroun, le Préambule de la Constitution (valeur constitutionnelle) les rend justiciables, mais le juge contrôle la \emph{raisonnabilité} de l'action étatique.
                \item \cle{Droits de solidarité (3\textsuperscript{e} gén.) :} Souvent \textbf{programmatiques}, justiciabilité plus difficile devant les juridictions classiques, mais possibles via les recours collectifs (action populaire) ou devant la Cour Africaine.
            \end{itemize}
        \item \textbf{Choisir la juridiction compétente :} Tribunal de première instance (TPI), Cour d'appel, Cour Suprême (Chambre administrative/Judiciaire/Constitutionnelle), Cour Africaine des Droits de l'Homme (si déclaration Art. 34.6 faite — \textbf{Attention : Cameroun n'a pas fait cette déclaration pour les individus/ONG}), Commission Africaine, Comité ONU (si Protocole facultatif ratifié).
    \end{enumerate}
\end{methode}

\begin{exoresolu}[Hiérarchie des normes et choix du recours - Cas pratique]
    \textbf{Énoncé :} Une association de défense des droits des femmes (ONG) souhaite contester une disposition du Code pénal camerounais (Art. 297 sur l'adultère) qui punit différemment l'homme et la femme. Elle estime que cette loi viole la Constitution (Préambule : « tous les êtres humains sont égaux en droits ») et la CEDAW (Convention sur l'élimination de toutes les formes de discrimination à l'égard des femmes) ratifiée par le Cameroun en 1994. Quelle est la hiérarchie des normes applicable ? Devant quelle juridiction nationale l'ONG peut-elle agir ? La saisine de la Cour Africaine des Droits de l'Homme et des Peuples (CADHP) est-elle possible directement ?
    \tcblower
    \textbf{Solution détaillée :}
    \begin{enumerate}
        \item \textbf{Hiérarchie des normes (Pyramide de Kelsen au Cameroun) :}
            \begin{enumerate}
                \item \textbf{Constitution (1996)} : Sommet. Le Préambule a valeur constitutionnelle (Déc. Cons. Const. n° 92/007). L'égalité homme/femme y est consacrée.
                \item \textbf{Traités internationaux (Art. 45 Const.)} : La \cle{CEDAW}, ratifiée et publiée, a une autorité supérieure à la loi nationale (Code pénal), sous réserve de réciprocité (remplie pour la CEDAW).
                \item \textbf{La Loi (Code pénal)} : L'Art. 297 est une loi ordinaire. Elle doit être conforme à la Constitution et aux traités.
            \end{enumerate}
            \textbf{Conséquence :} En cas de conflit, le traité (CEDAW) et la Constitution priment sur l'Art. 297 du Code pénal. Le juge national doit écarter la loi contraire au traité (contrôle de conventionalité).
        \item \textbf{Justiciabilité et Juridiction nationale :}
            \begin{itemize}
                \item L'ONG peut saisir le \textbf{Tribunal de Première Instance (TPI)} ou la \textbf{Cour d'Appel} via une \cle{exception d'inconstitutionnalité} (Art. 46 Const. : « Toute partie à un procès peut soulever l'inconstitutionnalité d'une loi... »). La juridiction sursoit à statuer et saisit le Conseil Constitutionnel.
                \item Ou saisine directe du \textbf{Conseil Constitutionnel} par 1/3 des députés/sénateurs (pas par l'ONG directement).
                \item \textbf{Voie la plus adaptée pour l'ONG :} Action en justice devant le TPI/Cour d'Appel $\rightarrow$ Exception d'inconstitutionnalité $\rightarrow$ Renvoi au Conseil Constitutionnel.
            \end{itemize}
        \item \textbf{Saisine de la Cour Africaine (CADHP) :}
            \begin{itemize}
                \item \textbf{Non possible directement par l'ONG.} Le Cameroun a ratifié le Protocole de 1998 créant la Cour, mais \textbf{n'a pas fait la déclaration au titre de l'Article 34.6} acceptant la compétence de la Cour pour les requêtes des individus et des ONG.
                \item Seuls la Commission Africaine, les États parties et les institutions africaines peuvent saisir la Cour contre le Cameroun.
                \item \textbf{Alternative :} Saisine de la \textbf{Commission Africaine des Droits de l'Homme et des Peuples (CADHP - Commission)} qui elle, est accessible aux ONG (Art. 55 Charte Africaine). La Commission peut ensuite renvoyer l'affaire à la Cour.
            \end{itemize}
        \item \textbf{Conclusion :} Hiérarchie : Const. = Traité (CEDAW) $>$ Loi (Code pénal). Recours national : Exception d'inconstitutionnalité devant juridiction de droit commun $\rightarrow$ Conseil Constitutionnel. Recours régional : Commission Africaine (oui), Cour Africaine (non, sauf si Commission y renvoie).
    \end{enumerate}
\end{exoresolu}

\begin{methode}[Identifier et situer les institutions de protection (Architecture multiniveau)]
    \textbf{Objectif :} Distinguer les niveaux de protection (ONU, Union Africaine, National) et leurs mécanismes spécifiques (rapports, communications, visites, contentieux).
    \begin{enumerate}
        \item \textbf{Niveau Universel (ONU - Siège Genève/New York) :}
            \begin{itemize}
                \item \cle{Conseil des Droits de l'Homme (CDH) :} Examen Périodique Universel (EPU) tous les 4-5 ans (État + société civile), Procédures spéciales (Rapporteurs thématiques/pays), Procédure de plainte (confidentielle).
                \item \cle{Comités des Traités (Organes conventionnels) :} Ex. Comité des Droits de l'Homme (PIDCP), Comité DESC (PIDESC), CEDAW, CAT, CDE. \textbf{Fonctions :} Examen des rapports périodiques États, \textbf{Communications individuelles} (si Protocole facultatif ratifié), Enquêtes, Observations générales.
                \item \cle{Haut-Commissariat aux Droits de l'Homme (HCDH) :} Secrétariat, assistance technique, terrain.
            \end{itemize}
        \item \textbf{Niveau Régional (Union Africaine - Siège Addis-Abeba / Cour à Arusha) :}
            \begin{itemize}
                \item \cle{Commission Africaine des Droits de l'Homme et des Peuples (CADHP - Banjul) :} Promotion (rapports États, missions), Protection (Communications États/Individus/ONG \textbf{Art. 55}), Interprétation Charte. Accessible aux ONG (statut observateur).
                \item \cle{Cour Africaine des Droits de l'Homme et des Peuples (CADHP - Arusha) :} Contentieux. \textbf{Condition :} Déclaration Art. 34.6 (Cameroun : NON). Accessible : Commission Africaine, États, Institutions UA.
                \item \cle{Comité d'experts (Droits de l'enfant) / Mécanismes pairs (PAP) :} Spécialisés.
            \end{itemize}
        \item \textbf{Niveau National (Cameroun - Yaoundé) :}
            \begin{itemize}
                \item \cle{Commission Nationale des Droits de l'Homme et des Libertés (CNDHL) :} Institution nationale (Statut A GANHRI). \textbf{Missions :} Promotion (éducation, formation), Protection (enquêtes, médiation, avis, saisine procureur), Conseil (gouvernement, Parlement). \textbf{Pouvoirs :} Autonomie financière/opérationnelle, accès lieux de détention, pouvoir d'auto-saisine.
                \item \cle{Juridictions :} Conseil Constitutionnel (constitutionnalité), Cour Suprême (admin/judiciaire), TPI/Cour d'Appel (droit commun).
                \item \cle{Autres :} Médiateur de la République (non créé formellement au Cameroun, fonctions proches assurées par CNDHL/SGPR), ONG agréées.
            \end{itemize}
        \item \textbf{Articulation :} Principe de \cle{subsidiarité} : Épuiser les voies de recours internes (national) avant saisine internationale (sauf urgence/ineffectivité). Complémentarité : CNDHL soumet rapports parallèles aux Comités ONU/Commission UA.
    \end{enumerate}
\end{methode}

\begin{exoresolu}[Architecture institutionnelle et stratégie de saisine]
    \textbf{Énoncé :} Un journaliste est détenu au secret depuis 10 jours au Secrétariat d'État à la Défense (SED) pour « atteinte à la sûreté de l'État » après un article sur la gestion des fonds Covid-19. Sa famille n'a pas de nouvelles. Son avocat n'a pas accès au dossier. La loi camerounaise (Code de procédure pénale) limite la garde à vue à 48h renouvelable une fois (96h max) sous contrôle du Procureur. Quelles institutions saisir en urgence aux trois niveaux (National, Africain, ONU) et selon quelles procédures spécifiques ?
    \tcblower
    \textbf{Solution détaillée :}
    \begin{enumerate}
        \item \textbf{Niveau National (Urgence absolue - Heures/Jours) :}
            \begin{itemize}
                \item \textbf{CNDHL :} \cle{Saisine d'urgence / Auto-saisine.} Pouvoir d'accès aux lieux de détention (Art. 17 Loi 2004/016 modifiée). Demande de visite au SED, vérification de la légalité de la détention (garde à vue > 96h = illégale), saisine du Procureur Général / Ministre de la Justice. \emph{Action la plus rapide et efficace sur le terrain.}
                \item \textbf{Juge des libertés / Procureur de la République (TPI) :} Requête en \cle{Habeas Corpus} (Art. 14 Const. / Art. 9 PIDCP / Art. 6 Charte Africaine) ou demande de mainlevée de garde à vue irrégulière. \emph{Via judiciaire contraignante.}
                \item \textbf{Avocat :} Saisine du Bâtonnier pour violation du droit à la défense (accès au dossier/client).
            \end{itemize}
        \item \textbf{Niveau Africain (Urgence - Jours/Semaines) :}
            \begin{itemize}
                \item \textbf{Commission Africaine (CADHP - Banjul) :} \cle{Communication d'urgence (Règle 111 du Règlement intérieur).} L'ONG/avocat saisit la Commission pour « mesures provisoires » (demande de libération/accès avocat/vie) \textbf{sans épuiser les voies de recours internes} si risque de préjudice irréparable (vie, torture, détention arbitraire). Accessible directement (Art. 55 Charte).
                \item \textbf{Rapporteur Spécial sur les Défenseurs des Droits de l'Homme / Liberté d'Expression (UA) :} Appel urgent (Urgent Appeal).
            \end{itemize}
        \item \textbf{Niveau ONU (Urgence - Jours/Semaines) :}
            \begin{itemize}
                \item \textbf{Groupe de travail de l'ONU sur la détention arbitraire (GTDA) :} \cle{Appel urgent (Urgent Appeal).} Procédure humanitaire, non contentieuse, rapide (semaines). Ne nécessite pas épuisement voies internes. Basé sur PIDCP (ratifié par Cameroun, \textbf{Protocole facultatif OP1 ratifié en 1984} $\rightarrow$ communications individuelles classiques possibles au Comité DH).
                \item \textbf{Comité des Droits de l'Homme (ONU) :} \cle{Communication individuelle (Art. 5 Prot. Fac. PIDCP).} Recevable car Cameroun a ratifié l'OP1 (1984). Épuisement voies internes requis (sauf si ineffectives/prolongées).
                \item \textbf{Rapporteurs Spéciaux ONU (Liberté d'expression, Indépendance juges/avocats, Torture) :} \cle{Allégations urgentes (Urgent Appeals) / Lettres d'allégation.} Mécanisme non juridictionnel, pression politique/diplomatique.
                \item \textbf{EPU (Examen Périodique Universel) :} Pas d'urgence (cycle 4-5 ans), mais soumission rapport parallèle (shadow report) pour documentation long terme.
            \end{itemize}
        \item \textbf{Stratégie intégrée :} 1. CNDHL + Habeas Corpus (National - effet juridique contraignant immédiat). 2. Commission Africaine (Mesures provisoires - effet juridique régional). 3. GTDA + Rapporteurs ONU + Comité DH (Pression diplomatique/juridique - effet politique/juridique). \textbf{Complémentarité :} Les décisions nationales (libération) rendent les saisines internationales sans objet ; les saisines internationales documentent et protègent si la justice nationale faillit.
    \end{enumerate}
\end{exoresolu}

\begin{methode}[Conduire l'étude de terrain d'une structure locale (TD1 - CNDHL ou ONG)]
    \textbf{Objectif :} Réaliser une enquête de terrain structurée (observation, entretien, analyse documentaire) et produire un rapport d'analyse critique.
    \begin{enumerate}
        \item \textbf{Préparation (Avant terrain) :}
            \begin{itemize}
                \item \textbf{Choix de la structure :} CNDHL (antenne régionale/départementale) ou ONG agréée (ex. REDHAC, ACDH, Réseau des Défenseurs DH).
                \item \textbf{Revue documentaire :} Textes fondateurs (Loi 2004/016 pour CNDHL, Statuts pour ONG), Rapports d'activité annuels, Rapports parallèles (ONU/UA), Site web.
                \item \textbf{Grille d'entretien/observation :} Thèmes : Mandat, Indépendance (financière, nomination), Moyens (humains, logistiques, budget), Actions concrètes (plaidoyer, formation, contentieux, visites prisons), Résultats/Indicateurs, Difficultés, Relation avec autorités/partenaires.
                \item \textbf{Demande d'autorisation/RDV :} Courrier officiel (chef d'établissement $\rightarrow$ Président CNDHL / Coordonnateur ONG).
            \end{itemize}
        \item \textbf{Terrain (Pendant) :}
            \begin{itemize}
                \item \textbf{Observation participante/non participante :} Locaux, accessibilité, affichage, accueil public, registre de plaintes.
                \item \textbf{Entretiens semi-directifs :} Responsable, enquêteurs/juristes, plaignants/bénéficiaires (si éthique/confidentialité respectée), partenaires locaux (Préfet, Procureur, Mairie).
                \item \textbf{Collecte de pièces :} Brochures, formulaires de plainte, affiches, organigramme, budget (si accessible).
            \end{itemize}
        \item \textbf{Analyse et Rédaction du Rapport (Après) :}
            \begin{itemize}
                \item \textbf{Structure type :} 1. Introduction (Contexte, Objectifs, Méthodologie, Limites). 2. Présentation de la structure (Cadre légal, Organigramme, Ressources). 3. Analyse des activités (Promotion, Protection, Conseil - \emph{avec exemples concrets}). 4. Évaluation critique (Forces / Faiblesses / Indépendance / Efficacité / Impact). 5. Recommandations (Réalistes, priorisées). 6. Conclusion. 7. Annexes (Grille, Photos, Documents).
                \item \textbf{Critères d'évaluation (Grille d'analyse) :} Conformité Principes de Paris (CNDHL), Indépendance, Accessibilité, Effectivité (traitement plaintes), Visibilité, Durabilité.
            \end{itemize}
        \item \textbf{Restitution orale :} Synthèse 10 min, support visuel (diaporama), débat.
    \end{enumerate}
\end{methode}

\begin{exoresolu}[Analyse critique d'un rapport de terrain (Simulation)]
    \textbf{Énoncé :} Votre groupe a visité l'antenne départementale de la CNDHL du Wouri. Vous avez noté : (1) Locaux dans l'enceinte de la Préfecture, pas de panneau indicatif extérieur. (2) 2 enquêteurs pour 3 millions d'habitants, pas de véhicule de service. (3) Registre de plaintes : 15 plaintes en 6 mois, 12 classées sans suite « faute de moyens », 3 transmises au Procureur (sans suite connue). (4) Président de l'antenne nommé par le Préfet. (5) Activité visible : 1 conférence/débat par an pour la Journée des DH (10 déc.). Rédigez la partie « Évaluation critique » de votre rapport en mobilisant les Principes de Paris et le mandat légal.
    \tcblower
    \textbf{Solution détaillée (Extrait - Évaluation critique) :}
    \textbf{1. Indépendance institutionnelle compromise (Principe de Paris § 2 \& 3 / Art. 2 Loi 2004/016) :}
    \begin{itemize}
        \item \textbf{Localisation :} Siège dans l'enceinte de la Préfecture $\rightarrow$ confusion avec l'administration, inaccessibilité psychologique pour les victimes de violations policières/administratives. Absence de signalétique $\rightarrow$ invisibilité.
        \item \textbf{Nomination :} Président nommé par le Préfet (autorité administrative) $\rightarrow$ violation du principe d'indépendance \emph{vis-à-vis du gouvernement}. La Loi 2004/016 (modifiée 2019) prévoit une désignation par les pairs ou élection au sein du bureau national, pas une nomination administrative locale.
    \end{itemize}
    \textbf{2. Insuffisance criante des moyens (Principe de Paris § 9 / Art. 14 Loi) :}
    \begin{itemize}
        \item Ratio 2 enquêteurs / 3 millions d'habitants = 1 pour 1,5 million (norme internationale : 1 pour 100 000 \emph{minimum}). Absence de véhicule $\rightarrow$ impossibilité matérielle d'enquêter sur le terrain (visites prisons, enquêtes terrain).
        \item Budget non communiqué mais implicitement nul (frais de mission non payés).
    \end{itemize}
    \textbf{3. Ineffectivité de la mission de protection (Mandat Art. 3, 5, 17 Loi) :}
    \begin{itemize}
        \item Taux de traitement : 80\% classés sans suite « faute de moyens » $\rightarrow$ déni de justice institutionnalisé.
        \item Transmission au Procureur sans suivi (pas de rôle de « watchdog » / chien de garde) $\rightarrow$ absence de pouvoir de suite / monitoring.
        \item Aucune auto-saisine documentée, aucune visite de lieu de détention (prison centrale, commissariats) sur la période.
    \end{itemize}
    \textbf{4. Promotion réduite à l'événementiel (Mandat Art. 3, 4 Loi) :}
    \begin{itemize}
        \item Une seule conférence/an $\rightarrow$ pas de programme continu d'éducation aux DH (écoles, entreprises, médias), pas de formation des acteurs (police, justice, administration).
    \end{itemize}
    \textbf{5. Forces (Rares) :} Existence légale, cadre physique (même précaire), registre tenu, lien formel avec Procureur.
    \textbf{Conclusion de l'évaluation :} L'antenne est \textbf{structurellement dysfonctionnelle}. Elle existe \emph{de jure} mais pas \emph{de facto}. Elle ne remplit pas les critères du Statut A (GANHRI) au niveau déconcentré. Risque de « décoration » institutionnelle.
    \textbf{Recommandations prioritaires :} (1) Délocalisation hors Préfecture + signalétique. (2) Recrutement 5 enquêteurs + dotation 1 véhicule/moto. (3) Formation enquêteurs (techniques enquête, droit pénal/procédure). (4) Mise en place suivi plaintes (tableau de bord, relances Procureur). (5) Plan d'action annuel promotion (écoles, marchés, radio locale).
\end{exoresolu}

\begin{methode}[Intégrer les savoir-faire : Qualifier, Saisiner, Argumenter (Démarche globale type Bac)]
    \textbf{Objectif :} Résoudre un cas complexe mobilisant l'ensemble de l'unité : qualification juridique, choix de l'institution/recours, argumentation juridique structurée (IRAC : Issue, Rule, Analysis, Conclusion).
    \begin{enumerate}
        \item \textbf{ÉTAPE 1 : QUALIFIER (Le « Quoi ») :}
            \begin{itemize}
                \item Lire le cas $\rightarrow$ Identifier les \textbf{faits juridiques} (pas seulement narratifs).
                \item Lister les \textbf{droits violés} (terminologie exacte : « droit à... », « liberté de... »).
                \item Classer par \textbf{génération} (CP, DESC, Solidarité) et noter \textbf{titulaires} (individus, groupes, peuples) et \textbf{obligés} (État, agents publics, privés).
            \end{itemize}
        \item \textbf{ÉTAPE 2 : HIÉRARCHISER LES NORMES (Le « Fondement ») :}
            \begin{itemize}
                \item Citer la \textbf{hiérarchie} : Constitution (Préambule/Art.) $>$ Traités (PIDCP, PIDESC, Charte Africaine, CEDAW, CIDE, CAT...) $>$ Lois nationales (Code Pénal, Travail, Procédure Pénale, Environnement...).
                \item Vérifier \textbf{ratification + publication} (J.O.) pour les traités.
                \item Identifier la \textbf{justiciabilité} : Immédiate (CP) / Progressive (DESC) / Programmatiques (Solidarité).
            \end{itemize}
        \item \textbf{ÉTAPE 3 : CHOISIR LA VOIE DE RECOURS / INSTITUTION (Le « Où » et « Comment ») :}
            \begin{itemize}
                \item \textbf{Interne (Priorité - Subsidiarité) :} CNDHL (urgence, médiation, enquête), Juridictions (TPI - Exception inconstitutionnalité, Référé liberté, Habeas Corpus, Action en responsabilité), Médiateur (si existant).
                \item \textbf{Régional :} Commission Africaine (Communications Art. 55, Mesures provisoires Règle 111), Cour Africaine (si déclaration Art. 34.6 - \textbf{Non pour CM}).
                \item \textbf{International :} Comités ONU (Communications si Prot. Fac. ratifié - \textbf{Vérifier pour CM : PIDCP OUI (OP1 1984), PIDESC NON, CEDAW OUI (OP-CEDAW 2004), CAT OUI (Art. 22 1987), CIDE OUI (OP3 2013) }), Procédures Spéciales (Appels urgents), EPU (Rapports parallèles).
            \end{itemize}
        \item \textbf{ÉTAPE 4 : ARGUMENTER (La « Démonstration ») - Structure IRAC :}
            \begin{itemize}
                \item \textbf{I - Issue (Problème juridique) :} « La détention au-delà de 96h sans présentation au juge viole-t-elle le droit à la liberté ? »
                \item \textbf{R - Rule (Règle de droit) :} Citer les articles précis : Art. 14 Const., Art. 9 PIDCP, Art. 6 Charte Africaine, Art. 119 CPP Cameroun. \emph{Ne pas citer de mémoire, citer le texte.}
                \item \textbf{A - Analysis (Application) :} Confronter faits $\rightarrow$ droit. « En l'espèce, M. X est détenu depuis 10 jours (fait) > 96h (règle CPP) $\rightarrow$ violation Art. 119 CPP + Art. 9 PIDCP (arbitraire). L'absence d'accès à l'avocat viole Art. 14 PIDCP + Art. 7 Charte. »
                \item \textbf{C - Conclusion (Demande/Dispositif) :} « Par conséquent, il y a lieu de demander la libération immédiate / la régularisation / la condamnation de l'État à réparer le préjudice. »
            \end{itemize}
        \item \textbf{ÉTAPE 5 : RÉDIGER LA RÉPONSE FINALE :} Phrases complètes, connecteurs logiques (donc, par ailleurs, en outre, toutefois), vocabulaire juridique précis (violation, grief, recevabilité, bien-fondé, réparation, mesures de non-répétition).
    \end{enumerate}
\end{methode}

\begin{exoresolu}[Cas d'intégration type Épreuve Probatoire/Bac]
    \textbf{Énoncé :} Suite à un éboulement dans une carrière artisanale non déclarée à Bafoussam, 5 jeunes mineurs (15-17 ans) trouvent la mort. L'exploitation était tolérée par les autorités locales (mairie, préfecture) contre redevances informelles. Les familles, non informées des risques, portent plainte au Tribunal de Grande Instance (TGI) pour homicide involontaire et demandent réparation. Le Procureur classe sans suite « auteur inconnu ». Les familles saisissent la CNDHL et une ONG internationale.
    \textbf{Questions :}
    1. Qualifiez les violations de droits humains subies par les victimes (enfants) et leurs familles. Classez-les.
    2. Identifiez les normes juridiques violées (Hiérarchie : Constitution, Traités, Lois nationales).
    3. Évaluez la recevabilité et la pertinence des saisines : TGI (classement sans suite), CNDHL, ONG internationale $\rightarrow$ Commission Africaine, Comité ONU (CIDE).
    4. Rédigez l'argumentaire juridique (structure IRAC) pour la saisine de la Commission Africaine aux fins de mesures provisoires et communication sur le fond.
    \tcblower
    \textbf{Solution détaillée :}
    \textbf{1. Qualification et Classification :}
    \begin{itemize}
        \item \textbf{Enfants victimes :}
            \begin{itemize}
                \item \cle{Droit à la vie (Art. 6 CIDE, Art. 4 Charte Africaine, Art. 6 PIDCP)} $\rightarrow$ \textbf{1\textsuperscript{re} génération (CP)}. Violation par omission (absence protection État).
                \item \cle{Droit à la survie et au développement (Art. 6 CIDE)} $\rightarrow$ \textbf{2\textsuperscript{e} génération (DESC) / Spécifique enfant}.
                \item \cle{Protection contre l'exploitation économique / travail dangereux (Art. 32 CIDE, Art. 15 Charte Africaine Enfant, Convention 182 OIT)} $\rightarrow$ \textbf{2\textsuperscript{e} génération (DESC)}. Travail des enfants dans carrière = pire forme de travail.
                \item \cle{Droit à la santé / environnement sain (Art. 24 CIDE, Art. 16 Charte Africaine, Art. 12 PIDESC)} $\rightarrow$ \textbf{2\textsuperscript{e} / 3\textsuperscript{e} génération}.
            \end{itemize}
        \item \textbf{Familles :}
            \begin{itemize}
                \item \cle{Droit à la vérité, justice, réparation (Art. 2 PIDCP, Principes Bassiouni/Van Boven)} $\rightarrow$ \textbf{1\textsuperscript{re} génération (Procédural)}.
                \item \cle{Droit à la protection de la famille (Art. 18 Charte Africaine, Art. 10 PIDESC)} $\rightarrow$ \textbf{2\textsuperscript{e} génération}.
            \end{itemize}
        \item \textbf{Caractères :} Indivisibilité (vie + travail + justice), Interdépendance (pauvreté $\rightarrow$ travail enfant $\rightarrow$ mort), Responsabilité de l'État (obligation de \emph{due diligence} : prévenir, enquêter, punir, réparer).
    \end{itemize}
    \textbf{2. Hiérarchie des Normes Violées :}
    \begin{enumerate}
        \item \textbf{Constitution (Préambule) :} Droit à la vie, protection de l'enfance, devoir de l'État de protéger les mineurs.
        \item \textbf{Traités (Supériorité Art. 45 Const.) :}
            \begin{itemize}
                \item \cle{CIDE (Ratifiée 1993) + Protocole Facultatif 3 (Comm. Indiv. - Ratifié 2013) :} Art. 3 (Intérêt supérieur), 6 (Vie/Développement), 19 (Protection violence), 32 (Travail).
                \item \cle{Charte Africaine Droits de l'Enfant (Ratifiée 1997) :} Art. 4 (Intérêt supérieur), 5 (Survie), 15 (Travail), 16 (Santé).
                \item \cle{PIDCP (Ratifiée 1984 + Prot. Fac. OP1 1984) :} Art. 2 (Recours effectif), 6 (Vie), 24 (Protection enfant).
                \item \cle{PIDESC (Ratifiée 1984) :} Art. 10 (Protection famille/enfants), 12 (Santé).
                \item \cle{Convention OIT 182 (Ratifiée 2002) :} Pires formes de travail.
            \end{itemize}
        \item \textbf{Lois Nationales :} Code Pénal (Art. 276 Homicide involontaire, Art. 350 Travail enfants), Code du Travail (Art. 86 Age minimum, Art. 87 Travaux interdits), Loi 2010/002 Protection Enfance, Loi sur l'Environnement (Étude d'impact absente).
    \end{enumerate}
    \textbf{3. Évaluation des Saisines :}
    \begin{itemize}
        \item \textbf{TGI (Classement sans suite) :} \textbf{Violation du droit à un recours effectif (Art. 2 PIDCP, Art. 7 Charte Africaine).} « Auteur inconnu » est fallacieux (autorités locales complices/tolérantes). \textbf{Voie de recours :} Appel du classement devant Chambre d'Accusation Cour d'Appel / Saisine Procureur Général / Constitution de partie civile (Art. 89 CPP) pour forcer instruction.
        \item \textbf{CNDHL :} \textbf{Pertinente et prioritaire.} Auto-saisine possible. Pouvoir d'enquête (accès carrière, audition autorités), médiation, rapport public, saisine Procureur Général / Ministre Justice. \emph{Action nationale la plus complète.}
        \item \textbf{Commission Africaine (ONG) :} \textbf{Recevable (Art. 55 Charte, Règle 111).} Épuisement voies internes \textbf{NON EXIGÉ} pour mesures provisoires (Règle 111) si urgence/préjudice irréparable (vie, impunité persistante). Sur le fond, épuisement requis (sauf si voies internes ineffectives/longues - argument recevable ici vu classement sans suite). \textbf{Compétence ratione materiae :} Charte Africaine + Charte Africaine Enfant.
        \item \textbf{Comité ONU Droits de l'Enfant (CIDE - Prot. Fac. 3) :} \textbf{Recevable (Protocole Facultatif 3 ratifié 2013).} Communication individuelle possible. Épuisement voies internes requis (sauf si déraisonnablement prolongées/ineffectives). Compétence sur CIDE uniquement.
    \end{itemize}
    \textbf{4. Argumentaire IRAC pour Commission Africaine (Mesures Provisoires + Fond) :}
    \begin{itemize}
        \item \textbf{I - ISSUE :} Le Cameroun a-t-il violé ses obligations de protection de la vie (Art. 4), interdiction travail dangereux enfants (Art. 15 Charte Enfant), droit à la justice/réparation (Art. 7 Charte, Art. 25 Charte Enfant) en tolérant une carrière illicite tuant 5 mineurs et en classant l'affaire sans suite ?
        \item \textbf{R - RULES :}
            \begin{itemize}
                \item Art. 1 Charte Africaine : Obligation de respecter et garantir les droits.
                \item Art. 4 Charte Africaine + Art. 5 Charte Enfant : Droit à la vie, survie, développement. \emph{Obligation positive de protection (Due diligence).}
                \item Art. 15 Charte Enfant : Protection contre exploitation économique/travail dangereux.
                \item Art. 7 Charte Africaine + Art. 25 Charte Enfant : Droit à un procès équitable / recours effectif.
                \item Jurisprudence Commission : \textit{Commission Nationale des Droits de l'Homme et des Libertés c. Tchad} (obligation protéger), \textit{Institute for Human Rights and Development in Africa c. Angola} (due diligence), \textit{Minority Rights Group c. Kenya} (peuples autochtones/environnement).
                \item Règle 111 Règlement Intérieur : Mesures provisoires pour éviter préjudice irréparable.
            \end{itemize}
        \item \textbf{A - ANALYSIS :}
            \begin{itemize}
                \item \textbf{Faits :} 5 morts, mineurs, travail illégal/dangereux, tolérance administrative (redevances), classement sans suite fallacieux.
                \item \textbf{Application :}
                    \begin{itemize}
                        \item \textbf{Vie/Survie :} L'État savait (redevances) ou devait savoir le danger. Absence de régulation/inspection/fermeture = \textbf{manquement à l'obligation positive de protection (due diligence)}. Violation Art. 4 Charte + Art. 5 Charte Enfant.
                        \item \textbf{Travail Enfant :} Tolérance administrative d'activité interdite (Code Travail, C182, Charte Enfant Art. 15) = \textbf{violation directe} obligation de prévention/répression.
                        \item \textbf{Justice/Réparation :} Classement sans suite « auteur inconnu » alors que responsables identifiables (exploitant, autorités tolérantes) = \textbf{déni de justice}. Violation Art. 7 Charte + Art. 25 Charte Enfant. Voies internes ineffectives (classement sans suite, impunité).
                    \end{itemize}
                \item \textbf{Urgence Mesures Provisoires (Règle 111) :} Risque de destruction preuves (carrière), intimidation familles, continuation exploitation illicite ailleurs. Préjudice irréparable pour familles (vérité, réparation) et autres enfants (risque réitération). $\rightarrow$ Demande : Enquête indépendante, suspension activités carrières illicites zone, protection familles/témoins, accès dossier judiciaire.
            \end{itemize}
        \item \textbf{C - CONCLUSION / DEMANDES :}
            \begin{enumerate}
                \item \textbf{Mesures Provisoires (Urgent) :} Ordonner au Cameroun : (a) Enquête judiciaire indépendante/impartiale sur l'éboulement et la chaîne de responsabilité (autorités locales). (b) Suspension immédiate de toute exploitation artisanale non conforme dans le département. (c) Protection physique/juridique des familles et témoins. (d) Accès au dossier judiciaire pour les avocats des familles.
                \item \textbf{Sur le Fond :} Déclarer le Cameroun responsable de violations des Art. 4, 7, 15 (Charte Enfant), 1, 25 (Charte Enfant). Ordonner : (a) Poursuite et sanction des auteurs (exploitants, complices administratifs). (b) Réparation intégrale aux familles (dommages-intérêts, mesures de réhabilitation psychosociale). (c) Mesures de non-répétition : Renforcement inspection travail/mines, campagne sensibilisation travail enfants, formation magistrats/police droits enfants, mise en place fonds indemnisation victimes travail enfants.
            \end{enumerate}
    \end{itemize}
\end{exoresolu}

\begin{attention}[Les 5 erreurs qui coûtent des points au Bac / Probatoire]
    \begin{enumerate}
        \item \textbf{Confondre classification des droits et simple énumération :} Dire « c'est le droit à la santé » sans préciser \cle{2\textsuperscript{e} génération (DESC)} et l'\cle{obligation de l'État (action/prestation vs abstention)}. \emph{Reflexe :} Toujours citer la génération + le type d'obligation étatique.
        \item \textbf{Ignorer la hiérarchie des normes (Art. 45 Const.) :} Citer seulement la loi nationale (Code pénal, Code du travail) sans vérifier la supériorité du Traité (PIDCP, CIDE, Charte) ou de la Constitution. \emph{Reflexe :} Pyramide : Const. $>$ Traité $>$ Loi $>$ Décret/Arrêté. Le juge écarte la loi contraire au traité.
        \item \textbf{Saisir la mauvaise institution / Procédure irrecevable :}
            \begin{itemize}
                \item Saisir la \textbf{Cour Africaine (Arusha)} directement pour un individu/ONG contre le Cameroun $\rightarrow$ \textbf{IRRECEVABLE} (Pas de déclaration Art. 34.6). Cible : \textbf{Commission Africaine (Banjul)}.
                \item Saisir le \textbf{Comité DESC (PIDESC)} pour communication individuelle $\rightarrow$ \textbf{IRRECEVABLE} (Cameroun n'a pas ratifié le Protocole facultatif au PIDESC).
                \item Oublier l'\textbf{épuisement des voies de recours internes} (sauf urgence/ineffectivité prouvée) pour les communications au fond (Commission Africaine, Comités ONU).
            \end{itemize}
        \item \textbf{Rédiger une argumentation « littéraire » sans structure juridique (IRAC) :} Raconter l'histoire au lieu de : \textbf{Problème de droit} $\rightarrow$ \textbf{Règle de droit (Articles précis)} $\rightarrow$ \textbf{Application faits/droit} $\rightarrow$ \textbf{Conclusion/Dispositif}. \emph{Sanction :} Hors sujet / Note faible en « Rédiger et justifier ».
        \item \textbf{Négliger la responsabilité de l'État pour actes d'acteurs non étatiques (Due Diligence) :} Dire « c'est l'entreprise qui a tué, pas l'État ». \emph{Erreur majeure.} L'État a l'obligation de \textbf{protéger} contre les violations par des tiers (entreprises, individus) : légiférer, réguler, enquêter, punir, réparer. L'inaction ou la tolérance (redevances informelles) = violation imputable à l'État.
    \end{enumerate}
\end{attention}

\newpage

\coursec{Exercices}

\courssub{Je vérifie mes connaissances}

\begin{exercice}[QCM : caractéristiques et générations des droits humains]
Pour chaque question, une seule réponse est exacte. Entoure la lettre correspondante. (1 point par question)

\textbf{1.} Le caractère \emph{universel} des droits humains signifie qu'ils :
\begin{itemize}
\item[a)] s'appliquent uniquement aux citoyens d'un État ;
\item[b)] appartiennent à tout être humain, sans distinction de race, de sexe ou de nationalité ;
\item[c)] sont identiques dans toutes les constitutions du monde ;
\item[d)] ne concernent que les adultes.
\end{itemize}

\textbf{2.} La liberté de la presse est un droit de :
\begin{itemize}
\item[a)] première génération (droits civils et politiques) ;
\item[b)] deuxième génération (droits économiques, sociaux et culturels) ;
\item[c)] troisième génération (droits de solidarité) ;
\item[d)] aucune génération, c'est une simple liberté publique.
\end{itemize}

\textbf{3.} Le droit à un environnement sain appartient à :
\begin{itemize}
\item[a)] la première génération ; \item[b)] la deuxième génération ; \item[c)] la troisième génération ; \item[d)] aucune génération.
\end{itemize}

\textbf{4.} Dire que les droits humains sont \emph{indivisibles} signifie que :
\begin{itemize}
\item[a)] on ne peut pas les diviser entre plusieurs personnes ;
\item[b)] toutes les catégories de droits sont liées et se renforcent mutuellement ;
\item[c)] ils ne peuvent pas être modifiés par la loi ;
\item[d)] ils sont réservés aux personnes indivisibles.
\end{itemize}

\textbf{5.} La Déclaration universelle des droits de l'homme a été adoptée en :
\begin{itemize}
\item[a)] 1945 ; \item[b)] 1948 ; \item[c)] 1960 ; \item[d)] 1996.
\end{itemize}
\end{exercice}

\begin{exercice}[Vrai ou faux : justifie ta réponse]
Pour chaque affirmation, indique si elle est \textbf{vraie} ou \textbf{fausse}, puis justifie en une ou deux phrases. (0,5 point pour la réponse + 0,5 point pour la justification)

\begin{enumerate}
\item Les droits humains peuvent être retirés à une personne qui a commis un délit.
\item La Constitution camerounaise de 1996 fait référence aux droits humains dans son préambule.
\item La CNDHL est un tribunal qui juge les violations des droits humains.
\item La Charte africaine des droits de l'homme et des peuples a été adoptée à Banjul.
\item Tout individu peut saisir directement la Cour pénale internationale.
\item Les droits économiques, sociaux et culturels comprennent le droit au travail et le droit à l'éducation.
\item Un traité international ratifié par le Cameroun peut être invoqué devant les juridictions nationales.
\item Le Médiateur de la République a exactement les mêmes compétences que la CNDHL.
\end{enumerate}
\end{exercice}

\begin{exercice}[Définitions express]
Définis en deux ou trois phrases chacun des termes suivants, en donnant un exemple concret tiré de la vie d'un atelier, d'une entreprise ou d'un quartier camerounais :
\begin{enumerate}
\item Droit humain ;
\item Justiciabilité ;
\item Saisine ;
\item Droit de l'homme de deuxième génération ;
\item Institution nationale des droits de l'homme.
\end{enumerate}
\end{exercice}

\begin{exercice}[Associe l'institution à sa mission]
Relie chaque institution (colonne A) à sa mission principale (colonne B) et précise son niveau d'action (international, régional ou national).

\begin{center}
\begin{tabularx}{\linewidth}{|l|X|}
\hline
\rowcolor{popblueL} \textbf{Colonne A : institution} & \textbf{Colonne B : mission} \\
\hline
1. CNDHL & a) Juge les crimes les plus graves (génocide, crimes contre l'humanité) \\
\hline
2. Commission africaine des droits de l'homme et des peuples & b) Reçoit les plaintes des citoyens contre les administrations \\
\hline
3. Cour pénale internationale & c) Promeut et protège les droits humains au niveau national, sans pouvoir de jugement \\
\hline
4. Médiateur de la République & d) Veille à l'application de la Charte africaine sur le continent \\
\hline
5. Conseil des droits de l'homme de l'ONU & e) Examine périodiquement la situation des droits humains dans les États membres \\
\hline
\end{tabularx}
\end{center}
\end{exercice}

\courssub{Je m'entraîne}

\begin{exercice}[Classe ces droits par génération]
Voici une liste de droits : liberté d'expression ; droit à la santé ; droit à un environnement sain ; droit de vote ; droit au travail ; droit à la paix ; droit à la vie ; droit à l'éducation ; droit au développement ; droit à un procès équitable.

\begin{enumerate}
\item Classe chacun de ces droits dans un tableau à trois colonnes : première, deuxième, troisième génération.
\item Choisis deux droits de la liste et explique, pour chacun, comment sa violation peut apparaître dans la vie quotidienne d'un apprenti ou d'un ouvrier au Cameroun.
\item Explique en quelques lignes pourquoi on dit que ces générations de droits sont \emph{interdépendantes}.
\end{enumerate}
\end{exercice}

\begin{exercice}[Qui saisir ?]
Pour chacune des situations suivantes, indique l'institution la plus adaptée à saisir (CNDHL, Médiateur de la République, tribunal judiciaire, Commission africaine des droits de l'homme et des peuples) et justifie ton choix en deux phrases.

\begin{enumerate}
\item Madame A., commerçante à Douala, attend depuis trois ans une décision de la mairie concernant sa licence, sans aucune réponse.
\item Des élèves d'un lycée technique dénoncent des châtiments corporels infligés par un surveillant.
\item Un État africain expulse massivement une communauté et refuse tout recours aux victimes, après épuisement des voies de recours internes.
\item Monsieur B. estime que son licenciement viole son contrat de travail et veut obtenir réparation.
\end{enumerate}
\end{exercice}

\begin{exercice}[Hiérarchie des normes]
On te donne les textes suivants : la Constitution camerounaise de 1996 ; un arrêté municipal de la ville de Bafoussam ; le Pacte international relatif aux droits civils et politiques (PIDCP) ; une loi ordinaire votée par l'Assemblée nationale.

\begin{enumerate}
\item Range ces quatre textes du plus au moins important dans la hiérarchie des normes, en justifiant chaque rang.
\item Un arrêté municipal interdit à toute personne de tenir une réunion publique dans la ville. Ce texte est-il conforme à la hiérarchie des normes ? Argumente en t'appuyant sur les libertés reconnues par les textes supérieurs.
\item Explique ce que signifie l'expression \og le traité ratifié a une valeur supérieure à la loi \fg{} (article 45 de la Constitution).
\end{enumerate}
\end{exercice}

\begin{exercice}[Analyse d'une situation d'atelier]
Dans une entreprise de menuiserie de Yaoundé, le patron impose aux apprentis des journées de 12 heures sans jour de repos, sans contrat écrit et sans couverture santé. Un apprenti qui proteste est renvoyé sur-le-champ.

\begin{enumerate}
\item Identifie dans cette situation trois droits humains bafoués et précise à quelle génération chacun appartient.
\item Pour chaque droit violé, cite le texte national ou international qui le protège (Constitution, Code du travail, PIDESC).
\item Propose deux démarches concrètes que l'apprenti renvoyé peut entreprendre pour faire valoir ses droits, en précisant les structures à saisir.
\end{enumerate}
\end{exercice}

\courssub{Je me prépare au Bac}

\begin{exercice}[Sujet type Probatoire : étude de document]
\textbf{Document 1 — Extrait de la Constitution du 18 janvier 1996 (article 45)} : \og Les traités ou accords internationaux dûment approuvés ou ratifiés ont, dès leur publication, une autorité supérieure à celle des lois \fg.

\textbf{Document 2 — Extrait (fictif mais réaliste) du rapport annuel de la CNDHL} :

\begin{center}
\begin{tabularx}{\linewidth}{|X|c|c|}
\hline
\rowcolor{popblueL} \textbf{Nature des plaintes reçues} & \textbf{Nombre} & \textbf{Part en \%} \\
\hline
Atteintes à la liberté (arrestations, détentions irrégulières) & 420 & 35 \\
\hline
Violations des droits économiques et sociaux (travail, logement) & 300 & 25 \\
\hline
Atteintes à la propriété foncière & 240 & 20 \\
\hline
Violences faites aux femmes et aux enfants & 120 & 10 \\
\hline
Autres plaintes & 120 & 10 \\
\hline
\textbf{Total} & \textbf{1200} & \textbf{100} \\
\hline
\end{tabularx}
\end{center}

La CNDHL recommande notamment : l'accélération des procédures judiciaires, la formation des forces de l'ordre aux droits humains et l'ouverture de nouvelles antennes régionales.

\textbf{Questions :}
\begin{enumerate}
\item Présente la CNDHL : nature, missions, pouvoirs. (4 points)
\item À l'aide de l'article 45, explique la place des traités internationaux relatifs aux droits humains dans l'ordre juridique camerounais. (3 points)
\item Analyse le tableau : quelle catégorie de plaintes domine ? Calcule la part cumulée des deux premières catégories et commente ce que cela révèle de l'état des droits humains au Cameroun. (4 points)
\item Relève deux recommandations de la CNDHL et explique comment chacune peut améliorer la protection des droits humains. (4 points)
\item La CNDHL peut-elle à elle seule garantir le respect des droits humains ? Justifie ta réponse en évoquant ses limites et les autres institutions complémentaires. (5 points)
\end{enumerate}
\end{exercice}

\begin{exercice}[Cas pratique noté : la détention de monsieur X]
Monsieur X, journaliste à Douala, publie un article sur la mauvaise gestion d'un marché public. Le lendemain, il est arrêté à son domicile \textbf{sans mandat}, conduit au commissariat et maintenu en garde à vue pendant \textbf{72 heures} sans pouvoir contacter ni sa famille ni un avocat. Il est ensuite inculpé de \og propagation de fausses nouvelles \fg{} sur la base de la loi de 2010 sur la cybercriminalité et la cybersécurité.

\textbf{Rappels de textes :} article 9 du PIDCP (droit à la liberté, interdiction de la détention arbitraire, droit d'être informé des charges) ; préambule et article 13 de la Constitution camerounaise ; loi de 2005 portant Code de procédure pénale (durée légale de la garde à vue, droit à l'assistance d'un avocat).

\textbf{Questions :}
\begin{enumerate}
\item Qualifie juridiquement la détention de monsieur X : relève trois irrégularités commises par les forces de l'ordre et cite pour chacune le texte violé. (6 points)
\item La liberté de la presse est-elle absolue ? Explique l'équilibre entre liberté d'expression et lutte contre les fausses informations. (4 points)
\item Rédige, en une quinzaine de lignes, une requête en référé-liberté que l'avocat de monsieur X adresse au président du tribunal administratif : tu y exposeras les faits, les moyens juridiques (textes violés) et la demande (libération immédiate). (6 points)
\item Au-delà du juge, quelles autres institutions monsieur X ou sa famille peuvent-ils saisir ? Présente deux d'entre elles et leur rôle possible dans cette affaire. (4 points)
\end{enumerate}
\end{exercice}

\begin{exercice}[Sujet type Baccalauréat technique : dissertation argumentée]
\textbf{Sujet :} \og L'effectivité des droits économiques, sociaux et culturels au Cameroun dépend-elle uniquement de la volonté politique ? \fg

\textbf{Consignes :} Après avoir défini les droits économiques, sociaux et culturels (droit au travail, à l'éducation, à la santé, au logement) et rappelé les textes qui les consacrent (PIDESC, préambule de la Constitution de 1996), tu discuteras la thèse proposée en mobilisant notamment : le rôle du budget de l'État et des ressources disponibles, l'action des institutions de protection (CNDHL, tribunaux), la contribution des ONG et de la société civile, et la responsabilité des citoyens eux-mêmes.

\textbf{Questions directrices :}
\begin{enumerate}
\item Définis les droits économiques, sociaux et culturels et montre en quoi ils se distinguent des droits civils et politiques. (4 points)
\item Montre que la volonté politique de l'État joue un rôle central dans la concrétisation de ces droits (lois, budgets sociaux, infrastructures scolaires et sanitaires). (6 points)
\item Montre ensuite que d'autres acteurs et d'autres facteurs interviennent : institutions de contrôle, ONG, partenaires internationaux, comportements des citoyens et des employeurs. (6 points)
\item Conclus en proposant deux mesures concrètes susceptibles de rendre ces droits plus effectifs pour les jeunes en formation technique au Cameroun. (4 points)
\end{enumerate}
\end{exercice}



\newpage

\coursec{Corrigés des exercices}

\begin{corrige}[Exercice 1 --- QCM : caractéristiques et générations des droits humains]
\textbf{1. b)} L'\cle{universalité} signifie que les droits humains appartiennent à tout être humain du seul fait de sa naissance, sans distinction de race, de sexe, de langue, de religion ou de nationalité. La réponse a) est fausse car les droits humains ne dépendent pas de la citoyenneté ; c) est fausse car toutes les constitutions ne les consacrent pas de la même manière ; d) est fausse car les enfants aussi ont des droits.

\textbf{2. a)} La liberté de la presse découle de la liberté d'expression, droit civil et politique de \cle{première génération} (article 19 de la DUDH et du PIDCP).

\textbf{3. c)} Le droit à un environnement sain est un droit de solidarité de \cle{troisième génération}, comme le droit au développement et le droit à la paix.

\textbf{4. b)} L'\cle{indivisibilité} signifie que les droits civils, politiques, économiques, sociaux et culturels sont liés et se renforcent mutuellement : on ne peut pas sacrifier une catégorie au profit d'une autre.

\textbf{5. b)} La DUDH a été adoptée par l'Assemblée générale des Nations unies le \textbf{10 décembre 1948} à Paris (1945 : création de l'ONU ; 1960 : indépendance du Cameroun ; 1996 : Constitution camerounaise).

\textbf{Points de vigilance :} ne pas confondre universalité (tous les humains) et uniformité (mêmes lois partout) ; bien mémoriser les trois dates charnières : 1948 (DUDH), 1966 (Pactes), 1981 (Charte africaine).
\end{corrige}

\begin{corrige}[Exercice 2 --- Vrai ou faux : justifie ta réponse]
\begin{enumerate}
\item \textbf{Faux.} Les droits humains sont \cle{inaliénables} : même un délinquant conserve ses droits fondamentaux (dignité, interdiction de la torture, procès équitable). Seules certaines libertés peuvent être limitées par une décision de justice rendue dans le respect de la loi (ex. : la liberté de circulation pour un condamné à une peine de prison).
\item \textbf{Vrai.} Le préambule de la Constitution du 18 janvier 1996 affirme l'attachement du peuple camerounais aux libertés fondamentales inscrites dans la DUDH, la Charte des Nations unies et la Charte africaine des droits de l'homme et des peuples.
\item \textbf{Faux.} La CNDHL est une institution nationale de \emph{promotion} et de \emph{protection} des droits humains : elle enquête, recommande et médiatise, mais elle \textbf{ne juge pas}. Le pouvoir de juger appartient aux tribunaux.
\item \textbf{Faux.} La Charte africaine a été adoptée à \cle{Nairobi} (Kenya) en 1981 et est entrée en vigueur en 1986 ; c'est son organe de contrôle, la Commission africaine, qui a son siège à Banjul (Gambie). La confusion Nairobi/Banjul est un piège classique.
\item \textbf{Faux.} Seuls le Procureur de la CPI, le Conseil de sécurité de l'ONU ou un État partie peuvent déclencher une procédure ; l'individu ne saisit pas directement la Cour pénale internationale.
\item \textbf{Vrai.} Le droit au travail, le droit à l'éducation, le droit à la santé et le droit au logement sont des droits de deuxième génération, consacrés par le PIDESC (1966).
\item \textbf{Vrai.} Selon l'article 45 de la Constitution, les traités dûment ratifiés ont, dès leur publication, une autorité supérieure à celle des lois : le citoyen peut donc les invoquer devant un juge camerounais.
\item \textbf{Faux.} Le Médiateur de la République règle les litiges entre les citoyens et les administrations (médiation administrative), tandis que la CNDHL traite spécifiquement des violations des droits humains : leurs champs de compétence sont différents, même s'ils se rejoignent parfois.
\end{enumerate}

\textbf{Points de vigilance :} la justification est notée autant que la réponse (0,5 + 0,5 point) : une réponse juste sans justification ne rapporte que la moitié des points. Toujours citer le principe (inaliénabilité, hiérarchie des normes) ou le texte (article 45) qui fonde la réponse.
\end{corrige}

\begin{corrige}[Exercice 3 --- Définitions express]
\begin{enumerate}
\item \textbf{Droit humain} : prérogative reconnue à tout être humain du seul fait de sa naissance ; il est universel, inaliénable et imprescriptible. \emph{Exemple :} un apprenti menuisier a droit au respect de sa dignité, même sans contrat écrit.
\item \textbf{Justiciabilité} : possibilité d'invoquer un droit devant un juge et d'obtenir une décision de justice. \emph{Exemple :} un ouvrier licencié abusivement peut saisir le tribunal du travail : son droit est justiciable.
\item \textbf{Saisine} : acte par lequel une personne porte une affaire devant une autorité (juge, CNDHL, Médiateur) pour qu'elle l'examine. \emph{Exemple :} une mère de famille saisit la CNDHL après l'arrestation irrégulière de son fils.
\item \textbf{Droit de deuxième génération} : droit économique, social ou culturel dont la réalisation exige une action positive de l'État (écoles, hôpitaux, emplois). \emph{Exemple :} le droit à la santé d'un ouvrier couvert par un dispensaire d'entreprise.
\item \textbf{Institution nationale des droits de l'homme} : organe public indépendant chargé de promouvoir et de protéger les droits humains dans un pays, sans pouvoir de jugement. \emph{Exemple :} la CNDHL au Cameroun, créée en 1990 et réorganisée par la loi de 2004.
\end{enumerate}

\textbf{Points de vigilance :} une définition complète comporte trois éléments : la nature (qu'est-ce que c'est ?), la caractéristique essentielle (qu'est-ce qui le distingue ?) et un exemple concret, exigé ici par l'énoncé. Ne pas confondre \emph{justiciabilité} (pouvoir aller devant le juge) et \emph{juridiction} (le tribunal lui-même).
\end{corrige}

\begin{corrige}[Exercice 4 --- Associe l'institution à sa mission]
\begin{itemize}
\item \textbf{1 -- c)} CNDHL : promeut et protège les droits humains sans pouvoir de jugement ; niveau \cle{national} (Cameroun).
\item \textbf{2 -- d)} Commission africaine des droits de l'homme et des peuples : veille à l'application de la Charte africaine ; niveau \cle{régional} (Union africaine, siège à Banjul).
\item \textbf{3 -- a)} Cour pénale internationale : juge les crimes les plus graves ; niveau \cle{international} (siège à La Haye).
\item \textbf{4 -- b)} Médiateur de la République : reçoit les plaintes des citoyens contre les administrations ; niveau \cle{national}.
\item \textbf{5 -- e)} Conseil des droits de l'homme de l'ONU : examine périodiquement la situation des États membres (Examen périodique universel) ; niveau \cle{international} (Genève).
\end{itemize}

\textbf{Points de vigilance :} l'énoncé demande deux choses : l'association \emph{et} le niveau d'action. Oublier le niveau fait perdre des points. Attention à ne pas attribuer un pouvoir de jugement à la CNDHL ou à la Commission africaine : seules les juridictions (tribunaux nationaux, CPI, Cour africaine) jugent.
\end{corrige}

\begin{corrige}[Exercice 5 --- Classe ces droits par génération]
\textbf{1.} Classement :

\begin{center}
\begin{tabularx}{\linewidth}{|X|X|X|}
\hline
\rowcolor{popblueL} \textbf{1\textsuperscript{re} génération (civils et politiques)} & \textbf{2\textsuperscript{e} génération (économiques, sociaux et culturels)} & \textbf{3\textsuperscript{e} génération (solidarité)} \\
\hline
Liberté d'expression ; droit de vote ; droit à la vie ; droit à un procès équitable & Droit à la santé ; droit au travail ; droit à l'éducation & Droit à un environnement sain ; droit à la paix ; droit au développement \\
\hline
\end{tabularx}
\end{center}

Vérification : $4 + 3 + 3 = 10$ droits, toute la liste est classée.

\textbf{2.} Exemples de violations : le \emph{droit au travail} est bafoué lorsqu'un apprenti travaille 12 heures par jour sans contrat ni rémunération ; la \emph{liberté d'expression} est violée lorsqu'un ouvrier est licencié pour avoir dénoncé des conditions de travail dangereuses.

\textbf{3.} Les générations sont interdépendantes car aucune ne peut pleinement exister sans les autres : sans éducation (2\textsuperscript{e} génération), on ne peut pas voter en connaissance de cause (1\textsuperscript{re} génération) ; sans paix (3\textsuperscript{e} génération), ni le travail ni les élections ne sont possibles.

\textbf{Points de vigilance :} le droit à la vie et le procès équitable relèvent de la 1\textsuperscript{re} génération (ce sont des droits civils), pas de la 2\textsuperscript{e} ; le droit au développement est de 3\textsuperscript{e} génération même s'il a une dimension économique. Pour la question 3, le mot-clé attendu est \emph{interdépendance} ou \emph{indivisibilité}, illustré par au moins un exemple.
\end{corrige}

\begin{corrige}[Exercice 6 --- Qui saisir ?]
\begin{enumerate}
\item \textbf{Médiateur de la République} : il s'agit d'un litige entre une citoyenne et une administration (la mairie) qui reste silencieuse depuis trois ans. Le Médiateur est compétent pour la médiation administrative et peut intervenir pour débloquer le dossier.
\item \textbf{CNDHL} : les châtiments corporels sont une atteinte à la dignité et à l'intégrité physique des enfants, donc une violation des droits humains. La CNDHL peut enquêter et recommander des mesures ; le procureur peut aussi être informé des faits.
\item \textbf{Commission africaine des droits de l'homme et des peuples} : les voies de recours internes étant épuisées, les victimes peuvent adresser une communication à la Commission, organe de contrôle de la Charte africaine, qui interdit expressément les expulsions massives (article 12).
\item \textbf{Tribunal judiciaire (tribunal du travail)} : le litige porte sur l'exécution d'un contrat de travail, matière régie par le Code du travail. Seul le juge peut ordonner une réparation (dommages et intérêts, voire réintégration).
\end{enumerate}

\textbf{Points de vigilance :} le choix de l'institution dépend de deux critères à citer dans la justification : la \emph{nature du litige} (administratif, droits humains, contrat de travail) et le \emph{niveau} (national d'abord, régional seulement après épuisement des recours internes). La condition d'épuisement des voies de recours internes est un point méthodologique essentiel pour la question 3.
\end{corrige}

\begin{corrige}[Exercice 7 --- Hiérarchie des normes]
\textbf{1.} Ordre décroissant :
\begin{enumerate}
\item la \textbf{Constitution de 1996} : norme suprême de l'État, toutes les autres doivent lui être conformes ;
\item le \textbf{PIDCP} : traité dûment ratifié, dont l'autorité est supérieure à celle des lois dès sa publication (article 45 de la Constitution) ;
\item la \textbf{loi ordinaire} : votée par le Parlement dans le respect de la Constitution et des traités ;
\item l'\textbf{arrêté municipal} : acte administratif local, qui doit respecter toutes les normes supérieures.
\end{enumerate}

\textbf{2.} Non, cet arrêté est \textbf{contraire} à la hiérarchie des normes : la liberté de réunion est garantie par le préambule de la Constitution et par l'article 21 du PIDCP. Une interdiction \emph{générale et absolue} décidée par un simple arrêté municipal viole donc des normes supérieures ; seules des restrictions prévues par la loi et nécessaires à l'ordre public ou à la sécurité sont admises.

\textbf{3.} L'article 45 signifie qu'en cas de conflit entre une loi camerounaise et un traité dûment ratifié et publié, c'est le \textbf{traité qui s'impose} : le juge doit écarter la loi contraire et appliquer le traité. C'est un mécanisme essentiel de protection des droits humains au Cameroun, car il permet d'invoquer directement le PIDCP ou la Charte africaine devant un tribunal national.

\textbf{Points de vigilance :} deux conditions de l'article 45 sont souvent oubliées et doivent être citées : le traité doit être \emph{dûment ratifié} \textbf{et} \emph{publié}. Attention aussi : le traité est supérieur à la loi mais reste \emph{inférieur à la Constitution}.
\end{corrige}

\begin{corrige}[Exercice 8 --- Analyse d'une situation d'atelier]
\textbf{1. et 2.} Trois droits bafoués :
\begin{itemize}
\item \textbf{Droit à des conditions de travail justes et favorables} (journées de 12 heures, absence de repos hebdomadaire) : droit de 2\textsuperscript{e} génération, protégé par l'article 7 du PIDESC et par le Code du travail camerounais (loi n° 92/007 du 14 août 1992 : durée légale du travail, repos hebdomadaire obligatoire).
\item \textbf{Droit à la santé et à la sécurité sociale} (aucune couverture santé) : droit de 2\textsuperscript{e} génération, protégé par les articles 9 et 12 du PIDESC et par le préambule de la Constitution de 1996.
\item \textbf{Droit à un procès équitable et liberté d'expression} (licenciement immédiat pour avoir protesté, sans procédure) : droits de 1\textsuperscript{re} génération, protégés par le préambule de la Constitution et par le Code du travail (procédure de licenciement encadrée).
\end{itemize}

\textbf{3.} Démarches concrètes : (a) saisir l'\textbf{inspection du travail} territorialement compétente pour faire constater les infractions et tenter une conciliation avec l'employeur ; (b) saisir le \textbf{tribunal du travail} pour contester le licenciement abusif et réclamer des indemnités. Il peut en outre déposer une plainte à la \textbf{CNDHL} pour dénoncer les conditions de travail dégradantes.

\textbf{Points de vigilance :} l'énoncé exige, pour chaque droit, trois éléments : l'identification, la génération \emph{et} le texte protecteur. Ne pas citer la CNDHL comme voie de réparation du licenciement : elle ne juge pas ; la réparation relève du tribunal du travail. Les démarches doivent être concrètes (structure nommée + résultat attendu).
\end{corrige}

\begin{corrige}[Exercice 9 --- Sujet type Probatoire : étude de document]
\textbf{1. Présentation de la CNDHL (4 points).} La Commission nationale des droits de l'homme et des libertés est une institution nationale indépendante, créée en 1990 et réorganisée par la loi n° 2004/016 du 22 juillet 2004. Ses missions : \emph{promouvoir} les droits humains (éducation, sensibilisation, formation) et les \emph{protéger} (réception des plaintes, enquêtes, visites des lieux de détention, médiation). Elle émet des recommandations aux pouvoirs publics mais ne juge pas. \emph{Barème indicatif : nature et cadre légal 1 pt ; missions de promotion 1 pt ; missions de protection 1 pt ; absence de pouvoir de jugement 1 pt.}

\textbf{2. Place des traités (3 points).} Selon l'article 45, les traités ratifiés et publiés (DUDH, PIDCP, PIDESC, Charte africaine) ont une autorité \emph{supérieure aux lois} camerounaises, tout en restant inférieurs à la Constitution. Tout citoyen peut donc les invoquer directement devant un juge national, qui doit écarter la loi contraire. \emph{Barème indicatif : supériorité sur les lois 1 pt ; conditions (ratification, publication) 1 pt ; conséquence pratique (invoquabilité devant le juge) 1 pt.}

\textbf{3. Analyse du tableau (4 points).} La catégorie dominante est celle des atteintes à la liberté : 420 plaintes, soit 35 \%. Part cumulée des deux premières catégories :
\[ 35\% + 25\% = 60\% \qquad \text{soit} \quad 420 + 300 = 720 \text{ plaintes sur } 1200. \]
Vérification : $720 / 1200 = 0{,}6 = 60\%$. Cela révèle que les violations les plus fréquentes touchent à la fois les droits civils (détentions irrégulières) et les droits économiques et sociaux : les deux premières générations de droits restent fragiles au Cameroun. \emph{Barème indicatif : catégorie dominante 1 pt ; calcul exact 1 pt ; commentaire 2 pts.}

\textbf{4. Recommandations (4 points).} Exemples : l'\emph{accélération des procédures judiciaires} réduit les détentions prolongées sans jugement ; la \emph{formation des forces de l'ordre} prévient les arrestations arbitraires et les mauvais traitements ; l'\emph{ouverture d'antennes régionales} rapproche la Commission des populations éloignées de Yaoundé. \emph{Barème indicatif : 1 pt par recommandation relevée + 1 pt par explication.}

\textbf{5. Limites de la CNDHL (5 points).} Non. La CNDHL n'a \textbf{pas de pouvoir de jugement} ni de contrainte : ses recommandations ne sont pas obligatoires, et ses moyens matériels sont limités. La protection effective exige l'action complémentaire des \textbf{tribunaux} (qui jugent et sanctionnent), du \textbf{Médiateur de la République} (médiation administrative), des institutions internationales (Conseil des droits de l'homme, Commission africaine) et de la \textbf{société civile} (ONG, médias, citoyens). \emph{Barème indicatif : réponse argumentée 1 pt ; limites (pas de jugement, recommandations non contraignantes) 2 pts ; deux institutions complémentaires avec leur rôle 2 pts.}

\textbf{Points de vigilance :} en étude de document, toujours appuyer la réponse sur les documents (citer l'article 45, utiliser les chiffres du tableau) ; un calcul non vérifié ou un pourcentage faux coûte cher : vérifier que la somme des parts fait bien 100 \%.
\end{corrige}

\begin{corrige}[Exercice 10 --- Cas pratique noté : la détention de monsieur X]
\textbf{1. Qualification juridique (6 points).} Trois irrégularités :
\begin{itemize}
\item \emph{Arrestation sans mandat} hors flagrant délit : détention arbitraire, contraire à l'article 9 du PIDCP et au Code de procédure pénale de 2005 ;
\item \emph{Garde à vue de 72 heures} sans renouvellement légal ni présentation à un magistrat : violation du Code de procédure pénale, qui encadre strictement la durée de la garde à vue ;
\item \emph{Refus de contacter la famille et un avocat} : violation du droit à l'assistance d'un conseil et du droit d'être informé des charges (article 9 du PIDCP, article 13 de la Constitution).
\end{itemize}
\emph{Barème indicatif : 1 pt par irrégularité relevée + 1 pt par texte correctement cité.}

\textbf{2. Liberté de la presse (4 points).} Non, elle n'est pas absolue : l'article 19 du PIDCP admet des restrictions prévues par la loi et nécessaires (respect de la réputation d'autrui, ordre public, sécurité nationale). La lutte contre les fausses informations est donc légitime, mais elle ne peut servir à réduire au silence un journaliste qui dénonce un fait d'intérêt public : la restriction doit être légale, nécessaire et proportionnée. \emph{Barème indicatif : réponse nuancée 1 pt ; fondement (article 19) 1 pt ; conditions des restrictions 1 pt ; application au cas 1 pt.}

\textbf{3. Requête en référé-liberté (6 points).} Modèle attendu : l'avocat expose les faits (arrestation sans mandat, garde à vue irrégulière, isolement), invoque les moyens (violation de l'article 9 du PIDCP, de l'article 13 de la Constitution et du Code de procédure pénale, textes supérieurs à la loi en vertu de l'article 45), souligne l'urgence et l'atteinte grave à la liberté, puis conclut en demandant au président du tribunal administratif d'ordonner la \emph{libération immédiate} de monsieur X. \emph{Barème indicatif : exposé des faits 1 pt ; moyens juridiques cités 2 pts ; forme de la requête (destinataire, urgence, demande) 2 pts ; clarté et correction 1 pt.}

\textbf{4. Autres institutions (4 points).} La \textbf{CNDHL}, qui peut ouvrir une enquête, visiter le lieu de détention et recommander la libération ; les \textbf{organisations de défense de la presse et ONG de droits humains}, qui peuvent médiatiser l'affaire et appuyer la défense. La famille peut aussi alerter le Procureur de la République sur l'irrégularité de la détention. \emph{Barème indicatif : 1 pt par institution nommée + 1 pt par rôle expliqué.}

\textbf{Points de vigilance :} dans un cas pratique, la méthode est toujours la même : \emph{faits -- qualification -- texte -- conséquence}. Ne jamais écrire qu'un texte est violé sans le nommer ; ne pas confondre référé-liberté (procédure d'urgence devant le juge administratif) et procès au fond.
\end{corrige}

\begin{corrige}[Exercice 11 --- Sujet type Baccalauréat technique : dissertation argumentée]
\textbf{1. Définition et distinction (4 points).} Les droits économiques, sociaux et culturels (DESC) sont les droits de deuxième génération : droit au travail, à l'éducation, à la santé, au logement, à la sécurité sociale. Ils sont consacrés par le PIDESC (1966) et par le préambule de la Constitution de 1996. Contrairement aux droits civils et politiques (1\textsuperscript{re} génération), qui exigent surtout l'\emph{abstention} de l'État (ne pas torturer, ne pas censurer), les DESC exigent une \emph{action positive} : construire des écoles, financer des hôpitaux, créer des emplois. \emph{Barème indicatif : définition 1 pt ; textes cités 1 pt ; distinction abstention/action positive 2 pts.}

\textbf{2. Rôle central de la volonté politique (6 points).} Seul l'État vote les lois sociales (Code du travail de 1992), consacre des budgets à l'éducation et à la santé, construit les lycées techniques et les centres de santé, et ratifie les traités internationaux. Sans décision politique, aucun de ces droits ne se concrétise à grande échelle : la gratuité de l'école primaire publique, par exemple, est d'abord une décision de l'État. \emph{Barème indicatif : trois arguments étayés (lois, budgets, infrastructures) 2 pts chacun.}

\textbf{3. Autres acteurs et facteurs (6 points).} La volonté politique ne suffit pas : les \emph{ressources budgétaires} limitées contraignent les choix de l'État ; les \emph{tribunaux} et la \emph{CNDHL} contrôlent et dénoncent les violations ; les \emph{ONG} et la société civile sensibilisent, assistent les victimes et construisent parfois des infrastructures ; les \emph{partenaires internationaux} (ONU, coopérations) financent des programmes ; enfin, les \emph{employeurs} doivent respecter le Code du travail et les \emph{citoyens} doivent réclamer leurs droits et entretenir les biens publics. \emph{Barème indicatif : au moins trois acteurs correctement présentés avec leur rôle, 2 pts chacun.}

\textbf{4. Conclusion et mesures concrètes (4 points).} Mesures pour les jeunes en formation technique : (a) augmenter le nombre de places, d'ateliers équipés et de stages conventionnés dans les lycées techniques, avec un suivi renforcé de l'inspection du travail ; (b) généraliser l'information sur les droits (clubs de citoyenneté, antennes régionales de la CNDHL) et faciliter l'accès gratuit à la justice et à la médiation pour les apprentis victimes d'abus. En définitive, l'effectivité des DESC résulte d'une \emph{responsabilité partagée} entre l'État, les institutions, la société civile et les citoyens. \emph{Barème indicatif : 1,5 pt par mesure concrète et réaliste + 1 pt pour une conclusion qui répond clairement à la question posée.}

\textbf{Points de vigilance :} la dissertation exige une introduction (définition des termes du sujet), un développement en deux temps (thèse puis nuances) et une conclusion personnelle. Éviter deux écueils : répondre seulement \og oui \fg{} ou \og non \fg{} sans discussion, et oublier d'ancrer les arguments dans le contexte camerounais (exemples nationaux valorisés).
\end{corrige}

\newpage

\coursec{Fiche bilan}

\begin{popfigure}
\begin{tikzpicture}[mindmap, grow cyclic, every node/.style={concept}, root concept/.append style={font=\bfseries\large, fill=popink, text=white, minimum size=3.2cm}, level 1 concept/.append style={sibling angle=60, level distance=4.6cm, minimum size=2.6cm, font=\small\bfseries}, level 2 concept/.append style={level distance=2.8cm, minimum size=2.1cm, font=\scriptsize}]
\node[root concept]{Les droits humains}
  child[concept color=popblue, text=white]{ node {Fondements}
    child[concept color=popblueL, text=popdark]{ node {Dignité, liberté, égalité} }
    child[concept color=popblueL, text=popdark]{ node {DUDH 1948} }
  }
  child[concept color=poporange, text=white]{ node {Caractéristiques}
    child[concept color=poporangeL, text=popdark]{ node {Universels, inaliénables} }
    child[concept color=poporangeL, text=popdark]{ node {Indivisibles, imprescriptibles} }
  }
  child[concept color=popgreen, text=white]{ node {Classification}
    child[concept color=popgreenL, text=popdark]{ node {1\iere{} génération : civils et politiques} }
    child[concept color=popgreenL, text=popdark]{ node {2\ieme{} génération : éco., soc., cult.} }
    child[concept color=popgreenL, text=popdark]{ node {3\ieme{} génération : solidarité} }
  }
  child[concept color=popgold, text=popdark]{ node {Cadre juridique}
    child[concept color=popgoldL, text=popdark]{ node {Préambule Constitution 1996} }
    child[concept color=popgoldL, text=popdark]{ node {Hiérarchie des normes} }
  }
  child[concept color=poppurple, text=white]{ node {Institutions}
    child[concept color=poppurpleL, text=popdark]{ node {ONU : CDH, Haut-Commissariat} }
    child[concept color=poppurpleL, text=popdark]{ node {UA : Commission et Cour africaines} }
    child[concept color=poppurpleL, text=popdark]{ node {National : CNDHL, tribunaux} }
  }
  child[concept color=poppink, text=white]{ node {Savoir-faire}
    child[concept color=poppinkL, text=popdark]{ node {Qualifier une violation} }
    child[concept color=poppinkL, text=popdark]{ node {Saisir la bonne structure} }
  };
\end{tikzpicture}
\legende{Carte mentale de la leçon : les droits humains (fondements, caractéristiques, classification, cadre juridique, institutions de protection et savoir-faire exigibles).}
\end{popfigure}

\begin{bilan}
\textbf{Définitions clés}
\begin{itemize}
  \item \cle{Droits humains} : droits fondamentaux reconnus à tout être humain, du seul fait de sa \cle{dignité}, sans distinction de race, de sexe, de langue, de religion ou d'opinion.
  \item \cle{DUDH} : Déclaration universelle des droits de l'homme, adoptée par l'ONU le 10 décembre 1948 ; texte fondateur des droits humains.
  \item \cle{Justiciabilité} : possibilité d'invoquer un droit devant un juge pour le faire respecter.
  \item \cle{CNDHL} : Commission nationale des droits de l'homme et des libertés, institution camerounaise de promotion et de protection des droits humains.
  \item \cle{Hiérarchie des normes} : ordre de valeur des règles de droit, de la Constitution (sommet) jusqu'aux décisions administratives (base).
\end{itemize}

\textbf{Repères à connaître par c\oe ur}
\begin{center}
\begin{tabularx}{\linewidth}{|l|X|}
\hline
\rowcolor{popblueL} \textbf{Notion} & \textbf{À retenir} \\
\hline
Caractéristiques & Universels, inaliénables, indivisibles, imprescriptibles, interdépendants \\
\hline
1\iere{} génération & Droits civils et politiques (vie, liberté, expression, vote) \\
\hline
2\ieme{} génération & Droits économiques, sociaux et culturels (travail, santé, éducation) \\
\hline
3\ieme{} génération & Droits de solidarité (paix, développement, environnement sain) \\
\hline
Niveau international & ONU : Conseil des droits de l'homme, Haut-Commissariat \\
\hline
Niveau régional & Union africaine : Commission africaine des droits de l'homme et des peuples, Cour africaine \\
\hline
Niveau national & Constitution de 1996 (préambule), tribunaux, CNDHL, ONG de défense \\
\hline
\end{tabularx}
\end{center}

\textbf{Méthodes express}
\begin{enumerate}
  \item \textbf{Qualifier} : identifier le droit atteint, nommer sa génération et le texte qui le garantit.
  \item \textbf{Saisir} : choisir la structure compétente (tribunal, CNDHL, ONG, institution internationale) selon le cas.
  \item \textbf{Argumenter} : fait $\rightarrow$ droit violé $\rightarrow$ texte de référence $\rightarrow$ institution compétente $\rightarrow$ conclusion justifiée.
\end{enumerate}
\end{bilan}

\begin{aretenir}[Checklist avant le Bac]
\begin{itemize}
  \item[$\square$] Je sais définir les droits humains en une phrase complète.
  \item[$\square$] Je cite la DUDH et sa date (10 décembre 1948).
  \item[$\square$] Je donne au moins quatre caractéristiques des droits humains.
  \item[$\square$] Je distingue les trois générations de droits avec un exemple chacune.
  \item[$\square$] Je situe la protection des droits humains dans le préambule de la Constitution camerounaise de 1996.
  \item[$\square$] Je décris la hiérarchie des normes du sommet à la base.
  \item[$\square$] Je nomme les institutions de protection aux trois niveaux (ONU, UA, national).
  \item[$\square$] Je connais le rôle de la CNDHL au Cameroun.
  \item[$\square$] Je sais qualifier une violation à partir d'une situation concrète (atelier, quartier, réseaux sociaux).
  \item[$\square$] Je sais rédiger une réponse en quatre étapes : fait, droit, texte, institution.
  \item[$\square$] Je justifie toujours ma conclusion par un texte ou une institution citée.
  \item[$\square$] Je relis ma copie : orthographe des sigles (DUDH, ONU, UA, CNDHL) vérifiée.
\end{itemize}
\end{aretenir}

\findecours

\end{document}
